% Les principes de la déplomatie camerounaise — ECM Tle Tech — Cours signé OnBuch+
% Programme officiel MINESEC. Compiler avec : tectonic N02-principes-deplomatie-camerounaise.tex
\newcommand{\DOCMATIERE}{ECM}
\newcommand{\DOCNIVEAU}{Tle Tech}
\newcommand{\DOCMODULE}{ECM – classe de Terminale technique}
\newcommand{\DOCLECON}{N02}
\newcommand{\DOCTITRE}{Les principes de la déplomatie camerounaise}
% OnBuch+ — préambule « cours signés » ENRICHI (compilé avec tectonic / XeLaTeX)
% Système visuel « Pop » d'OnBuch+ : fond crème (jamais blanc), bordures encre
% épaisses, ombres « dures » décalées, titres Archivo Black en capitales.
% Version MATHÉMATIQUES : graphes (pgfplots/tikz), boîtes pédagogiques
% (définition, propriété, méthode, attention, le savais-tu, exercice résolu,
% exercice, TP, bilan), page de garde et signature OnBuch+ ancrée en bas.
%
% Macros attendues dans main.tex AVANT \input{preamble} :
%   \DOCMATIERE  \DOCNIVEAU  \DOCMODULE  \DOCLECON (numéro)  \DOCTITRE
\documentclass[11pt]{article}

\usepackage{fontspec}
\usepackage[french]{babel}
\usepackage[a4paper,margin=2cm,top=3.4cm,bottom=2.8cm,headheight=1.4cm,footskip=1.3cm]{geometry}
\usepackage{amsmath,amssymb,amsfonts}
\usepackage{mathtools} % \xmapsto, \coloneqq... souvent utilisés par les modèles
\usepackage{cancel} % simplifications barrées (bilans de forces)
% \abs{z} (module, valeur absolue) et \norm{u} (norme), utilisés par les modèles
\providecommand{\abs}{}\let\abs\relax\DeclarePairedDelimiter{\abs}{\lvert}{\rvert}
\providecommand{\norm}{}\let\norm\relax\DeclarePairedDelimiter{\norm}{\lVert}{\rVert}
\usepackage[table]{xcolor} % [table] : fournit aussi \rowcolors (alternance de lignes)
\usepackage{pagecolor}
\usepackage{tikz}
\usetikzlibrary{arrows.meta,positioning,calc,shapes.geometric,shapes.misc,shapes.arrows,decorations.pathmorphing,decorations.pathreplacing,decorations.markings,patterns,shadows,mindmap,trees,fit,backgrounds,matrix,intersections,angles,quotes}
\usepackage{pgfplots}
\usepackage[version=4]{mhchem} % équations nucléaires \ce{^{235}_{92}U}
\usepackage[european,siunitx]{circuitikz} % schémas électriques
\pgfplotsset{compat=1.18}
\usepgfplotslibrary{fillbetween,groupplots} % aires entre courbes (calcul intégral)
\usepackage{enumitem}
\usepackage{fancyhdr}
% [most] charge toutes les bibliothèques tcolorbox (listings, minted, raster,
% documentation...) et gonfle inutilement la mémoire TeX sur un document long
% (~300 boîtes) : on ne charge que ce qui est réellement utilisé.
\usepackage[skins,breakable]{tcolorbox}
\usepackage{graphicx}
\usepackage{colortbl}
\usepackage{booktabs}
\usepackage{tabularx}
\usepackage{multirow}
\usepackage{diagbox} % case d'en-tête diagonale des tableaux à double entrée
% Colonnes extensibles centrée / gauche / droite (C, L, R), souvent utilisées par les modèles
\newcolumntype{C}{>{\centering\arraybackslash}X}
\newcolumntype{L}{>{\raggedright\arraybackslash}X}
\newcolumntype{R}{>{\raggedleft\arraybackslash}X}
\usepackage{array}
\usepackage{multicol}
\usepackage{siunitx}
% parse-numbers=false : accepte aussi \SI{2,5 \times 10^{-3}}{...}
\sisetup{locale=FR,output-decimal-marker={,},group-minimum-digits=4,parse-numbers=false}
\DeclareSIUnit\ohm{\ensuremath{\symup{\Omega}}} % U+2126 absent de la police
\DeclareSIUnit\division{div} % réglages d'oscilloscope (ms/div, V/div)
\DeclareSIUnit\tr{tr} % tours (vitesse de rotation, ex. \SI{1500}{\tr\per\minute})
\DeclareSIUnit\molar{M} % concentration molaire (ex. \SI{3}{\molar}, \SI{1}{\micro\molar})
\DeclareSIUnit\an{an} % année (le modèle confond parfois avec \year, un compteur TeX) — ex. \SI{5}{\kilo\meter\per\an}
\DeclareSIUnit\FCFA{FCFA} % franc CFA (ex. \SI{500}{\FCFA\per\kilo\gram})
\DeclareSIUnit\degreeBrix{\textdegree Brix} % teneur en sucre d'un sirop/jus (réfractométrie)
\DeclareSIUnit\month{mois} % le modèle utilise parfois \month, un compteur TeX, comme unité de durée
\DeclareSIUnit\microliter{\micro\liter} % le modèle écrit parfois \microliter en un seul token
\DeclareSIUnit\million{millions} % dénombrement (ex. \SI{15}{\million\per\mL} spermatozoïdes)
\usepackage{qrcode}
% \degree : commande fréquemment utilisée par les modèles pour un angle ou une
% température en dehors de \SI{}{} (non fournie par les paquets chargés).
\providecommand{\degree}{^{\circ}}
% \qed (fin de démonstration), utilisé par les modèles sans amsthm
\providecommand{\qed}{\hfill$\square$}
% \barre : double barre des valeurs interdites dans un tableau de variations (tkz-tab)
\providecommand{\barre}{\ensuremath{\|}}
% environnement proof (amsthm non chargé) utilisé par les modèles
\ifdefined\proof\else\newenvironment{proof}[1][Démonstration]{\par\noindent\textit{#1.}\ }{\qed\par}\fi
\usepackage{lastpage}
\usepackage{hyperref}
\hypersetup{hidelinks}

% ============================================================
% Palette « Pop » OnBuch+ — mêmes hex que la classe P (plus_ui.dart)
% ============================================================
\definecolor{poporange}{HTML}{F59321}
\definecolor{popdark}{HTML}{E07A0C}
\definecolor{popcream}{HTML}{FFF6E8}
\definecolor{popink}{HTML}{1C1714}
\definecolor{poppurple}{HTML}{7A5AE0}
\definecolor{popgreen}{HTML}{2E9E6A}
\definecolor{poppink}{HTML}{E0457B}
\definecolor{popblue}{HTML}{2F7DE1}
\definecolor{popgold}{HTML}{C98A12}
\definecolor{popmuted}{HTML}{8A7F76}
\definecolor{popline}{HTML}{EADFD2}
\definecolor{popgray}{HTML}{8A7F76}
\colorlet{popgrey}{popgray}
% Alias tolérants (le modèle nomme parfois les couleurs différemment)
\colorlet{popwhite}{white}
\colorlet{popblack}{popink}
\colorlet{popred}{poppink}
\colorlet{popyellow}{popgold}
% Teintes claires (fonds de tableaux/encadrés) — le modèle nomme parfois
% les couleurs avec un suffixe L (ex. popblueL) sans les avoir définies.
\colorlet{poporangeL}{poporange!20}
\colorlet{popdarkL}{popdark!20}
\colorlet{poppurpleL}{poppurple!20}
\colorlet{popgreenL}{popgreen!20}
\colorlet{poppinkL}{poppink!20}
\colorlet{popblueL}{popblue!20}
\colorlet{popgoldL}{popgold!20}
\colorlet{popredL}{popred!20}
\colorlet{popgrayL}{popgray!20}
% Teintes claires pour fonds de tableaux / zones
\colorlet{popblueL}{popblue!10!white}
\colorlet{poporangeL}{poporange!14!white}
\colorlet{popgreenL}{popgreen!12!white}
\colorlet{poppurpleL}{poppurple!12!white}
\colorlet{poppinkL}{poppink!12!white}
\colorlet{popgoldL}{popgold!14!white}

\pagecolor{popcream}

% ============================================================
% Typographie
% ============================================================
\defaultfontfeatures{Ligatures=TeX}
\setmainfont{PlusJakartaSans-Medium.ttf}[Path=../../../fonts/,BoldFont=PlusJakartaSans-Bold.ttf]
% Maths en sans-serif (Fira Math) avec chiffres et lettres droites de Plus
% Jakarta, pour que les formules se fondent dans le texte.
\usepackage[math-style=ISO,bold-style=ISO]{unicode-math}
\setmathfont{FiraMath-Regular.otf}[Path=../../../fonts/]
\setmathfont{PlusJakartaSans-Medium.ttf}[Path=../../../fonts/,range={up/num,up/{Latin,latin},\mathup}]
\usepackage{icomma} % virgule décimale sans espace en mode math
% Caractères mathématiques Unicode absents des polices de texte (Plus Jakarta,
% Archivo Black) : rendus par leur équivalent mathématique, en texte comme en
% formule — sinon ils s'affichent en carré vide (ex. « Arithmétique dans ℤ »).
\usepackage{newunicodechar}
\newunicodechar{ℝ}{\ensuremath{\mathbb{R}}}
\newunicodechar{ℤ}{\ensuremath{\mathbb{Z}}}
\newunicodechar{ℂ}{\ensuremath{\mathbb{C}}}
\newunicodechar{ℕ}{\ensuremath{\mathbb{N}}}
\newunicodechar{ℚ}{\ensuremath{\mathbb{Q}}}
\newunicodechar{𝔻}{\ensuremath{\mathbb{D}}}
\newunicodechar{α}{\ensuremath{\alpha}}
\newunicodechar{β}{\ensuremath{\beta}}
\newunicodechar{γ}{\ensuremath{\gamma}}
\newunicodechar{δ}{\ensuremath{\delta}}
\newunicodechar{ε}{\ensuremath{\varepsilon}}
\newunicodechar{θ}{\ensuremath{\theta}}
\newunicodechar{λ}{\ensuremath{\lambda}}
\newunicodechar{μ}{\ensuremath{\mu}}
\newunicodechar{σ}{\ensuremath{\sigma}}
\newunicodechar{τ}{\ensuremath{\tau}}
\newunicodechar{φ}{\ensuremath{\varphi}}
\newunicodechar{ψ}{\ensuremath{\psi}}
\newunicodechar{ω}{\ensuremath{\omega}}
\newunicodechar{Σ}{\ensuremath{\Sigma}}
% majuscules produites par \MakeUppercase dans les titres (α-aminés -> Α)
\newunicodechar{Α}{\ensuremath{\alpha}}
\newunicodechar{Β}{\ensuremath{\beta}}
\newunicodechar{Γ}{\ensuremath{\Gamma}}
\newunicodechar{Δ}{\ensuremath{\Delta}}
\newunicodechar{Ε}{\ensuremath{\varepsilon}}
\newunicodechar{Θ}{\ensuremath{\Theta}}
\newunicodechar{Λ}{\ensuremath{\Lambda}}
\newunicodechar{Μ}{\ensuremath{\mu}}
\newunicodechar{Τ}{\ensuremath{\tau}}
\newunicodechar{Φ}{\ensuremath{\Phi}}
\newunicodechar{Ψ}{\ensuremath{\Psi}}
\newunicodechar{Ω}{\ensuremath{\Omega}}
\newunicodechar{↦}{\ensuremath{\mapsto}}
\newunicodechar{∈}{\ensuremath{\in}}
\newunicodechar{∉}{\ensuremath{\notin}}
\newunicodechar{∧}{\ensuremath{\wedge}}
\newunicodechar{⇔}{\ensuremath{\Leftrightarrow}}
\newunicodechar{⟺}{\ensuremath{\Longleftrightarrow}}
\newunicodechar{⇒}{\ensuremath{\Rightarrow}}
\newunicodechar{⟹}{\ensuremath{\Longrightarrow}}
\newunicodechar{∘}{\ensuremath{\circ}}
\newunicodechar{∪}{\ensuremath{\cup}}
\newunicodechar{∩}{\ensuremath{\cap}}
\newunicodechar{≡}{\ensuremath{\equiv}}
\newunicodechar{⊂}{\ensuremath{\subset}}
\newunicodechar{⊥}{\ensuremath{\perp}}
\newunicodechar{∃}{\ensuremath{\exists}}
\newunicodechar{∀}{\ensuremath{\forall}}
\newunicodechar{∛}{\ensuremath{\sqrt[3]{\,}}}
\newunicodechar{∜}{\ensuremath{\sqrt[4]{\,}}}
\newunicodechar{⃗}{\ensuremath{{}^{\to}}}
\newunicodechar{ⁿ}{\ifmmode^{n}\else\textsuperscript{\ensuremath{n}}\fi}
\newunicodechar{ˣ}{\ifmmode^{x}\else\textsuperscript{\ensuremath{x}}\fi}
\newunicodechar{ᵇ}{\ifmmode^{b}\else\textsuperscript{\ensuremath{b}}\fi}
\newunicodechar{ʳ}{\ifmmode^{r}\else\textsuperscript{\ensuremath{r}}\fi}
\newunicodechar{ᵘ}{\ifmmode^{u}\else\textsuperscript{\ensuremath{u}}\fi}
\newunicodechar{ʸ}{\ifmmode^{y}\else\textsuperscript{\ensuremath{y}}\fi}
\newunicodechar{ᵃ}{\ifmmode^{a}\else\textsuperscript{\ensuremath{a}}\fi}
\newunicodechar{ᵉ}{\ifmmode^{e}\else\textsuperscript{\ensuremath{e}}\fi}
\newunicodechar{ᵏ}{\ifmmode^{k}\else\textsuperscript{\ensuremath{k}}\fi}
\newunicodechar{ᵖ}{\ifmmode^{p}\else\textsuperscript{\ensuremath{p}}\fi}
\newunicodechar{ᴺ}{\ifmmode^{N}\else\textsuperscript{\ensuremath{N}}\fi}
\newunicodechar{ᶜ}{\ifmmode^{c}\else\textsuperscript{\ensuremath{c}}\fi}
\newunicodechar{ʰ}{\ifmmode^{h}\else\textsuperscript{\ensuremath{h}}\fi}
\newunicodechar{ᵠ}{\ifmmode^{q}\else\textsuperscript{\ensuremath{q}}\fi}
\newunicodechar{ₙ}{\ifmmode_{n}\else\textsubscript{\ensuremath{n}}\fi}
\newunicodechar{ₐ}{\ifmmode_{a}\else\textsubscript{\ensuremath{a}}\fi}
\newunicodechar{ᵢ}{\ifmmode_{i}\else\textsubscript{\ensuremath{i}}\fi}
\newunicodechar{ᵣ}{\ifmmode_{r}\else\textsubscript{\ensuremath{r}}\fi}
\newunicodechar{ₖ}{\ifmmode_{k}\else\textsubscript{\ensuremath{k}}\fi}
\newunicodechar{ᵦ}{\ifmmode_{\beta}\else\textsubscript{\ensuremath{\beta}}\fi}
\newunicodechar{ₓ}{\ifmmode_{x}\else\textsubscript{\ensuremath{x}}\fi}
\newfontfamily\popdisplay{ArchivoBlack-Regular.ttf}[Path=../../../fonts/]
\newfontfamily\poplogo{SpaceGrotesk-Bold.ttf}[Path=../../../fonts/]
\newfontfamily\popsemi{PlusJakartaSans-SemiBold.ttf}[Path=../../../fonts/]
\newfontfamily\popxbold{PlusJakartaSans-ExtraBold.ttf}[Path=../../../fonts/]
\newfontfamily\popxboldls{PlusJakartaSans-ExtraBold.ttf}[Path=../../../fonts/,LetterSpace=3.0]

\setlist[enumerate]{leftmargin=1.7em,itemsep=5pt,topsep=4pt}
\setlist[itemize]{leftmargin=1.7em,itemsep=4pt,topsep=4pt,label={\color{poporange}\textbullet}}
\setlength{\parindent}{0pt}
\setlength{\parskip}{5pt}
\renewcommand{\arraystretch}{1.25}

% pgfplots : style Pop par défaut
\pgfplotsset{
  every axis/.append style={axis line style={popink,line width=1pt},tick style={popink},
    grid style={popline,line width=0.5pt},label style={font=\small},tick label style={font=\footnotesize}},
  popcurve/.style={poporange,line width=1.8pt,no markers,smooth},
}
% Même style disponible dans un \draw TikZ simple (hors pgfplots)
\tikzset{popcurve/.style={poporange,line width=1.8pt}}
\tikzset{popgrid/.style={popline,very thin}}
\colorlet{popcurve}{poporange} % le nom du style est parfois utilisé comme couleur

% ============================================================
% Pill / squiggle
% ============================================================
\newcommand{\poppill}[3]{%
  \tikz[baseline=-0.55ex]\node[draw=popink,line width=1.2pt,rounded corners=2.6pt,%
    fill=#2,inner xsep=6.5pt,inner ysep=3pt]{%
    \popxboldls\fontsize{7}{7}\selectfont\color{#3}\MakeUppercase{#1}};%
}
\newcommand{\popsquiggle}[1]{%
  \tikz[baseline=-0.4ex]{
    \draw[#1,line width=2.6pt,line cap=round]
      plot [smooth] coordinates {(0,0.09) (0.34,0.01) (0.68,0.21) (1.02,0.01) (1.36,0.10)};
  }%
}
% Mot-clé surligné (marqueur orange sous le texte)
\newcommand{\cle}[1]{{\popxbold\color{popdark}#1}}

% ============================================================
% En-tête / pied
% ============================================================
\pagestyle{fancy}
\fancyhf{}
\renewcommand{\headrulewidth}{0pt}
\renewcommand{\footrulewidth}{0pt}
\lhead{\raisebox{-0.15ex}{\poplogo\fontsize{13}{13}\selectfont\color{popink}OnBuch\color{poporange}+}}
\rhead{\poppill{\DOCMATIERE\ \textbullet\ \DOCNIVEAU\ \textbullet\ Leçon \DOCLECON}{white}{popink}}
\lfoot{\popxbold\fontsize{7.6}{7.6}\selectfont\color{popmuted}\MakeUppercase{Cours signé OnBuch+}\ \textbullet\ {\popsemi\DOCTITRE}}
\rfoot{\tikz[baseline=-0.6ex]\node[rounded corners=8.5pt,fill=popink,minimum height=17pt,minimum width=17pt,inner xsep=5pt,inner sep=0]{\popxbold\fontsize{8}{8}\selectfont\color{popcream}\thepage\,/\,\pageref*{LastPage}};}

% ============================================================
% Titres
% ============================================================
\newcommand{\chaptitre}[2]{%
  \begin{tcolorbox}[enhanced,colback=poporange,colframe=popink,boxrule=2.6pt,arc=11pt,
      drop shadow=popink,left=15pt,right=15pt,top=12pt,bottom=11pt,before skip=3pt,after skip=18pt]
    \poppill{#2}{popink}{popcream}\par\vspace{9pt}
    {\raggedright\hyphenpenalty=10000\exhyphenpenalty=10000\popdisplay\fontsize{22}{24}\selectfont\color{popcream}\MakeUppercase{#1}\par}
  \end{tcolorbox}
}
% Section numérotée : pastille orange avec le numéro + titre Archivo
\newcounter{popsec}
\newcommand{\coursec}[1]{%
  \stepcounter{popsec}\setcounter{popsub}{0}%
  \par\vspace{16pt}\noindent
  \begin{minipage}{\linewidth}
  \tikz[baseline=-0.7ex]\node[draw=popink,line width=1.6pt,fill=poporange,rounded corners=4pt,minimum size=22pt,inner sep=0]{\popdisplay\fontsize{12}{12}\selectfont\color{popcream}\thepopsec};%
  \hspace{8pt}{\popdisplay\fontsize{15}{17}\selectfont\color{popink}\MakeUppercase{#1}}\par
  \vspace{4pt}\noindent{\color{poporange}\rule{36pt}{3.6pt}}
  \end{minipage}\par\nopagebreak\vspace{8pt}
}
\newcounter{popsub}
\newcommand{\courssub}[1]{%
  \stepcounter{popsub}%
  \par\vspace{9pt}\noindent{\popxbold\fontsize{11.5}{13}\selectfont\color{popink}{\color{poporange}\thepopsec.\thepopsub}\hspace{6pt}#1}\par\nopagebreak\vspace{4pt}
}

% ============================================================
% Boîtes pédagogiques — même recette : fond blanc, bordure épaisse colorée,
% ombre dure encre, badge plein. Titre optionnel [#1].
% ============================================================
\tcbset{popbase/.style={breakable,enhanced,colback=white,boxrule=2.3pt,arc=9pt,
  shadow={3pt}{-3pt}{0pt}{popink},left=14pt,right=14pt,top=11pt,bottom=11pt,before skip=11pt,after skip=15pt,
  fontupper=\popsemi\fontsize{10.6}{14.5}\selectfont\color{popink},
  fontlower=\popsemi\fontsize{10.6}{14.5}\selectfont\color{popink}}}
% Coche dessinée (le glyphe ✓ n'existe pas dans les polices du document)
\newcommand{\popcheck}[1]{\tikz[baseline=-0.5ex]\draw[#1,line width=1.6pt,line cap=round,line join=round](-0.13,0.0)--(-0.04,-0.09)--(0.13,0.1);}
\providecommand{\coche}{\popcheck{popgreen}}
\providecommand{\resultat}[1]{\par\smallskip\noindent\fcolorbox{popgreen}{popgreenL}{\parbox{\dimexpr\linewidth-2\fboxsep-2\fboxrule\relax}{\textbf{Résultat.} #1}}\par\smallskip}
\newcommand{\popboxhead}[3]{\poppill{#1}{#2}{white}\ifx\relax#3\relax\else\hspace{6pt}{\popxbold\fontsize{10.6}{12}\selectfont\color{popink}#3}\fi\par\vspace{7pt}}

\newtcolorbox{aretenir}[1][]{popbase,colframe=popgreen,before upper={\popboxhead{À retenir}{popgreen}{#1}}}
\newtcolorbox{exemplebox}[1][]{popbase,colframe=poporange,before upper={\popboxhead{Exemple}{poporange}{#1}}}
\newtcolorbox{definition}[1][]{popbase,colframe=popblue,before upper={\popboxhead{Définition}{popblue}{#1}}}
\newtcolorbox{propriete}[1][]{popbase,colframe=poppurple,before upper={\popboxhead{Propriété}{poppurple}{#1}}}
\newtcolorbox{methode}[1][]{popbase,colframe=popgold,before upper={\popboxhead{Méthode}{popgold}{#1}}}
\newtcolorbox{attention}[1][]{popbase,colframe=poppink,before upper={\popboxhead{Attention}{poppink}{#1}}}
\newtcolorbox{experience}[1][]{popbase,colframe=popblue,colback=popblueL,before upper={\popboxhead{Expérience}{popblue}{#1}}}
% « Le savais-tu ? » : carte violette pleine (contexte, culture, Cameroun)
\newtcolorbox{savaistu}[1][]{popbase,colback=poppurple,colframe=popink,
  fontupper=\popsemi\fontsize{10.4}{14.2}\selectfont\color{white},
  before upper={\poppill{Le savais-tu\,?}{white}{poppurple}\ifx\relax#1\relax\else\hspace{6pt}{\popxbold\color{white}#1}\fi\par\vspace{7pt}}}
% Exercice résolu : énoncé en haut, \tcblower puis la solution sur fond crème
\newcounter{popexo}\newcounter{popexr}
\newtcolorbox{exoresolu}[1][]{popbase,colframe=poporange,colbacklower=popcream,
  segmentation style={popink,line width=1.2pt,dashed},
  before upper={\stepcounter{popexr}\popboxhead{Exercice résolu \thepopexr}{poporange}{#1}},
  before lower={\poppill{Solution}{popink}{popcream}\par\vspace{6pt}}}
% Exercice (non résolu en place) — la correction est dans la partie Corrigés
\newtcolorbox{exercice}[1][]{popbase,colframe=popink,
  before upper={\stepcounter{popexo}\popboxhead{Exercice \thepopexo}{popink}{#1}}}
% Corrigé
\newtcolorbox{corrige}[1][]{popbase,colframe=popgreen,colback=popgreenL,
  before upper={\popboxhead{Corrigé}{popgreen}{#1}}}
% Objectifs / prérequis / bilan
\newtcolorbox{objectifs}{popbase,colframe=popink,colback=poporangeL,
  before upper={\popboxhead{Objectifs de la leçon}{popink}{}}}
\newtcolorbox{prerequis}{popbase,colframe=popink,colback=popblueL,
  before upper={\popboxhead{Prérequis}{popblue}{}}}
\newtcolorbox{bilan}{popbase,colframe=popink,colback=white,boxrule=2.8pt,
  before upper={\popboxhead{Fiche bilan}{poporange}{L'essentiel à connaître pour l'examen}}}
% Figure encadrée avec légende
\newtcolorbox{popfigure}[1][]{enhanced,breakable=false,colback=white,colframe=popline,boxrule=1.4pt,arc=8pt,
  left=10pt,right=10pt,top=10pt,bottom=8pt,before skip=10pt,after skip=12pt,halign=center,
  lower separated=false,halign lower=center,
  fontlower=\popsemi\fontsize{9}{11.5}\selectfont\color{popmuted},
  #1}
\newcommand{\legende}[1]{\par\vspace{6pt}{\popsemi\fontsize{9}{11.5}\selectfont\color{popmuted}#1}}

% ============================================================
% Signature OnBuch+
% ============================================================
\newcommand{\poplogoinline}[1]{{\poplogo\fontsize{#1}{#1}\selectfont\color{popink}OnBuch\color{poporange}+}}
\newcommand{\signatureonbuch}{%
  \begin{tcolorbox}[enhanced,colback=popink,colframe=popink,boxrule=2.2pt,arc=10pt,
      drop shadow=poporange,left=14pt,right=14pt,top=10pt,bottom=10pt,before skip=0pt,after skip=0pt]
    \tikz[baseline=-0.7ex]\node[circle,fill=popgreen,draw=popcream,line width=1.3pt,minimum size=22pt,inner sep=0]{\popcheck{white}};%
    \hspace{9pt}\begin{minipage}[c]{0.62\linewidth}
      {\popxbold\fontsize{12}{13}\selectfont\color{popcream}Signé\ \textbullet\ {\poplogo OnBuch\color{poporange}+}}\par\vspace{2pt}
      {\raggedright\popsemi\fontsize{8}{10.5}\selectfont\color{popline}\MakeUppercase{Équipe pédagogique OnBuch+}\\\MakeUppercase{Conforme au programme officiel MINESEC}\par}
    \end{minipage}\hfill
    {\poplogo\fontsize{22}{22}\selectfont\color{popcream}OnBuch\color{poporange}+}
  \end{tcolorbox}%
}

% ============================================================
% Page de garde
% #1 titre · #2 matière · #3 niveau · #4 module · #5 n° leçon · #6 durée ·
% #7 description / objectifs (texte ou itemize)
% ============================================================
\newcommand{\pagegarde}[7]{%
  \thispagestyle{empty}
  \newgeometry{a4paper,margin=1.7cm,top=1.6cm,bottom=1.4cm}%
  \begin{tikzpicture}[remember picture,overlay]
    \fill[poporange!18] ([xshift=-3.2cm,yshift=-3.6cm]current page.north east) circle (4.6cm);
    \fill[poppurple!14] ([xshift=2.4cm,yshift=7.6cm]current page.south west) circle (3.8cm);
  \end{tikzpicture}%
  {\poplogo\fontsize{30}{30}\selectfont\color{popink}OnBuch\color{poporange}+}\hfill
  \raisebox{0.6ex}{\poppill{Cours signé \textbullet\ Premium}{popink}{popcream}}\par
  \vspace{1pt}\popsquiggle{poppurple}\par
  \vspace{8pt}
  \begin{tcolorbox}[enhanced,colback=poporange,colframe=popink,boxrule=2.8pt,arc=13pt,
      drop shadow=popink,left=18pt,right=18pt,top=12pt,bottom=13pt,after skip=10pt]
    \poppill{#2 \textbullet\ #3}{popink}{popcream}\hspace{5pt}\poppill{#4}{white}{popink}\par\vspace{8pt}
    {\popdisplay\fontsize{14}{14}\selectfont\color{popink}LEÇON #5}\par\vspace{4pt}
    {\raggedright\hyphenpenalty=10000\exhyphenpenalty=10000\popdisplay\fontsize{25}{27}\selectfont\color{popcream}\MakeUppercase{#1}\par}
    \vspace{8pt}
    {\popxbold\fontsize{9.5}{9.5}\selectfont\color{popink}\MakeUppercase{Durée conseillée : #6}}
  \end{tcolorbox}
  \begin{tcolorbox}[enhanced,colback=white,colframe=popink,boxrule=2.2pt,arc=10pt,
      drop shadow=popink,left=16pt,right=16pt,top=10pt,bottom=10pt,after skip=8pt]
    \poppill{Dans cette leçon}{poppurple}{white}\par\vspace{5pt}
    \popsemi\fontsize{9.3}{12.6}\selectfont\color{popink}#7
  \end{tcolorbox}
  \noindent\begin{minipage}[t]{0.487\linewidth}
    \begin{tcolorbox}[enhanced,colback=popgreen,colframe=popink,boxrule=2.2pt,arc=10pt,
        drop shadow=popink,left=11pt,right=11pt,top=10pt,bottom=10pt,height=3.1cm,valign=center]
      \tcbox[colback=white,colframe=popink,boxrule=1.3pt,arc=5pt,boxsep=0pt,nobeforeafter,
        left=3pt,right=3pt,top=3pt,bottom=3pt,valign=center]{\qrcode[height=1.9cm]{https://play.google.com/store/apps/details?id=com.luvvix.onbuch&hl=fr}}%
      \hspace{8pt}\begin{minipage}[c]{0.5\linewidth}
        {\popdisplay\fontsize{11}{12.5}\selectfont\color{white}\MakeUppercase{Télécharge OnBuch}}\par\vspace{4pt}
        {\raggedright\popsemi\fontsize{8.4}{11}\selectfont\color{white}Gratuit \textbullet\ Google Play\\Tuteur IA, annales, concours, résultats\par}
      \end{minipage}
    \end{tcolorbox}
  \end{minipage}\hfill
  \begin{minipage}[t]{0.487\linewidth}
    \begin{tcolorbox}[enhanced,colback=poppurple,colframe=popink,boxrule=2.2pt,arc=10pt,
        drop shadow=popink,left=11pt,right=11pt,top=10pt,bottom=10pt,height=3.1cm,valign=center]
      \tikz[baseline=-0.7ex]\node[circle,fill=white,draw=popink,line width=1.3pt,minimum size=1.6cm,inner sep=0]{\popdisplay\fontsize{24}{24}\selectfont\color{poppurple}+};%
      \hspace{8pt}\begin{minipage}[c]{0.6\linewidth}
        {\popdisplay\fontsize{11}{12.5}\selectfont\color{white}\MakeUppercase{Passe à OnBuch+}}\par\vspace{4pt}
        {\raggedright\popsemi\fontsize{8.4}{11}\selectfont\color{white}Tous les cours signés, Tuteur IA illimité, fiches PDF — onglet OnBuch+ de l'app\par}
      \end{minipage}
    \end{tcolorbox}
  \end{minipage}
  \vspace{8pt}
  \popcontenu
  \vfill
  \signatureonbuch
  \restoregeometry
  \newpage
}

% Rangée « Ce cours contient » (page de garde)
\newcommand{\poptile}[3]{% #1 couleur #2 gros texte #3 légende
  \begin{tcolorbox}[enhanced,colback=#1,colframe=popink,boxrule=2pt,arc=9pt,shadow={2.5pt}{-2.5pt}{0pt}{popink},
      left=6pt,right=6pt,top=9pt,bottom=9pt,halign=center,valign=center,height=1.75cm,width=\linewidth,nobeforeafter]
    {\popdisplay\fontsize{12.5}{13}\selectfont\color{white}\MakeUppercase{#2}}\par\vspace{3pt}
    {\centering\popsemi\fontsize{7.8}{9.5}\selectfont\color{white}#3\par}
  \end{tcolorbox}}
\newcommand{\popcontenu}{%
  \par\noindent{\popxboldls\fontsize{7.5}{7.5}\selectfont\color{popmuted}CE COURS SIGNÉ CONTIENT}\par\vspace{7pt}
  \noindent\begin{minipage}[t]{0.235\linewidth}\poptile{popblue}{Cours}{complet, illustré, pas à pas}\end{minipage}\hfill
  \begin{minipage}[t]{0.235\linewidth}\poptile{poporange}{Méthodes}{et exercices résolus}\end{minipage}\hfill
  \begin{minipage}[t]{0.235\linewidth}\poptile{poppink}{Exercices}{type examen + corrigés}\end{minipage}\hfill
  \begin{minipage}[t]{0.235\linewidth}\poptile{popgreen}{Fiche bilan}{l'essentiel à retenir}\end{minipage}\par}

% Sommaire
\newtcolorbox{sommaire}{enhanced,colback=white,colframe=popink,boxrule=2.2pt,arc=10pt,
  drop shadow=popink,left=18pt,right=18pt,top=13pt,bottom=13pt,before skip=6pt,after skip=20pt,
  before upper={\poppill{Sommaire}{poppurple}{white}\par\vspace{9pt}\popsemi\fontsize{10.6}{14.5}\selectfont\color{popink}}}

% Signature de fin de cours — poussée tout en bas de la dernière page
\newcommand{\findecours}{%
  \par\vfill\nopagebreak
  \begin{center}\popsquiggle{poporange}\end{center}\vspace{4pt}
  \signatureonbuch
}
\AtBeginDocument{\def\llbracket{\mathopen{[\mkern-3.5mu[}}\def\rrbracket{\mathclose{]\mkern-3.5mu]}}} % intervalles d'entiers (glyphe ⟦ absent de la police)
% Glyphes absents de Fira Math : redéfinis à partir de glyphes disponibles
\AtBeginDocument{%
  \def\setminus{\mathbin{\backslash}}\let\smallsetminus\setminus
  \def\vdots{\mathord{\vbox{\baselineskip3.2pt\lineskiplimit0pt\kern4pt\hbox{.}\hbox{.}\hbox{.}}}}%
  \def\ddots{\mathinner{\mkern1mu\raise6.5pt\hbox{.}\mkern2mu\raise3.6pt\hbox{.}\mkern2mu\raise0.7pt\hbox{.}\mkern1mu}}%
  \def\triangle{\mathord{\scalebox{0.9}{$\symup{\Delta}$}}}%
  \def\nabla{\mathord{\raisebox{\depth}{\scalebox{1}[-1]{$\symup{\Delta}$}}}}%
  \def\checkmark{\popcheck{popdark}}%
  \def\star{\mathbin{\ast}}\def\bigstar{\mathord{\ast}}%
  \def\diamond{\mathbin{\scalebox{0.6}{$\Diamond$}}}%
  \def\bigcup{\mathop{\mathchoice{\raisebox{-0.3em}{\scalebox{1.7}{$\cup$}}}{\scalebox{1.25}{$\cup$}}{\cup}{\cup}}\displaylimits}%
  \def\bigcap{\mathop{\mathchoice{\raisebox{-0.3em}{\scalebox{1.7}{$\cap$}}}{\scalebox{1.25}{$\cap$}}{\cap}{\cap}}\displaylimits}%
  \def\lesseqqgtr{\mathrel{\lessgtr}}\def\bigoplus{\mathop{\oplus}\displaylimits}%
}
\providecommand{\ssi}{si et seulement si} % abréviation fréquente
\AtBeginDocument{\providecommand{\wideparen}{\overparen}} % arc de cercle

%% FIGFIX-BEGIN v5 — correctifs globaux des figures (docs/onbuch-generation/tools/figfix)
\usepackage{adjustbox}\usepackage{etoolbox}\tcbuselibrary{raster}
% toute figure TikZ/pgfplots plus large que la ligne est réduite à la largeur disponible
\BeforeBeginEnvironment{tikzpicture}{\begin{adjustbox}{max width=\linewidth}}
\AfterEndEnvironment{tikzpicture}{\end{adjustbox}}
% légendes pgfplots lisibles par-dessus les courbes
\pgfplotsset{every axis/.append style={legend style={fill=white,fill opacity=0.92,text opacity=1,draw=black!15,inner sep=3pt}}}
% étiquettes d'arêtes (syntaxe edge["..."]) avec fond blanc
\tikzset{every edge quotes/.append style={fill=white,inner sep=1.2pt}}
%% FIGFIX-END
\begin{document}

\pagegarde{Les principes de la déplomatie camerounaise}{ECM}{Tle Tech}{ECM – classe de Terminale technique}{N02}{4 séances (fiche de progression harmonisée)}{Cette leçon permet aux élèves de maîtriser les concepts fondamentaux des relations internationales, les formes de coopération, l'assistance humanitaire et les principes directeurs de la diplomatie du Cameroun, tout en développant des compétences d'analyse documentaire et de rédaction argumentée.
\vspace{4pt}
\begin{multicols}{2}\small
\begin{itemize}
\item Définir le concept de relations internationales et en identifier les acteurs principaux.
\item Classer les différentes formes de coopération internationale (bilatérale, multilatérale, décentralisée, publique‑privée).
\item Analyser un dossier d'assistance humanitaire en extrayant les données chiffrées et les acteurs impliqués.
\item Énoncer et expliquer les cinq principes fondamentaux de la diplomatie camerounaise.
\item Appliquer ces principes à une situation concrète (ex. : négociation d'un accord de pêche).
\item Rédiger une note diplomatique structurée en justifiant chaque argument par des références aux principes étudiés.
\end{itemize}\end{multicols}}

\begin{sommaire}
\begin{multicols}{2}
\begin{enumerate}
\item 1. Le concept de Relations Internationales
\item 2. Aspects et formes de la coopération internationale
\item 3. Dossier 1 : L'assistance humanitaire
\item 4. Les principes de la diplomatie camerounaise
\item 5. Synthèse et mise en situation : rédaction d'une note diplomatique
\item Activité d'intégration
\item Méthodes et exercices résolus
\item Exercices
\item Corrigés des exercices
\item Fiche bilan
\end{enumerate}
\end{multicols}
\end{sommaire}

\begin{objectifs}
\begin{itemize}
\item Définir le concept de relations internationales et en identifier les acteurs principaux.
\item Classer les différentes formes de coopération internationale (bilatérale, multilatérale, décentralisée, publique‑privée).
\item Analyser un dossier d'assistance humanitaire en extrayant les données chiffrées et les acteurs impliqués.
\item Énoncer et expliquer les cinq principes fondamentaux de la diplomatie camerounaise.
\item Appliquer ces principes à une situation concrète (ex. : négociation d'un accord de pêche).
\item Rédiger une note diplomatique structurée en justifiant chaque argument par des références aux principes étudiés.
\item Construire un croquis cartographique localisant les missions diplomatiques du Cameroun et les principaux flux d'aide humanitaire.
\item Évaluer la pertinence d'une politique de coopération au regard des objectifs de développement durable (ODD).
\end{itemize}
\end{objectifs}

\begin{prerequis}
\begin{itemize}
\item Notion de souveraineté et d'État (programme de 4ème).
\item Repérage des principaux organismes internationaux (ONU, UA, CEMAC) vu en début d'année.
\item Lecture de cartes thématiques (flux, réseaux) acquise en géographie physique.
\item Techniques de base de la rédaction administrative (lettre formelle, rapport).
\end{itemize}
\end{prerequis}

\newpage

\coursec{Le concept de Relations Internationales}

Imaginez un jeune ingénieur de Douala, fraîchement diplômé de l'ENSET, qui bénéficie d'un accord de coopération technique entre le Cameroun et l'Allemagne. Il part se former à Munich dans une entreprise de pointe en mécatronique, découvre l'industrie 4.0, puis revient créer une start-up locale qui embauche une dizaine de jeunes techniciens. Ce parcours, loin d'être un cas isolé, illustre concrètement comment les \cle{relations internationales} et la \cle{diplomatie camerounaise} façonnent l'avenir professionnel des élèves de Terminale technique. Elles ne sont pas l'apanage des chancelleries : elles déterminent l'accès aux bourses, aux marchés, aux technologies, aux normes de certification et aux opportunités d'emploi.

\courssub{Définition académique et évolution historique}

\begin{definition}[Relations internationales]
Les \textbf{relations internationales (RI)} désignent l'ensemble des interactions régulières, structurées et normées qui s'établissent entre les acteurs de la scène mondiale — États, organisations internationales, organisations non gouvernementales, entreprises multinationales, collectivités territoriales — en vue de défendre des intérêts, de résoudre des conflits ou de produire des biens publics mondiaux (sécurité, climat, santé, commerce, éducation).
\end{definition}

\emph{Intuition première :} pensez aux RI comme à un immense « réseau social » planétaire où chaque profil (État, ONU, ONG, entreprise) publie, commente, négocie ou rompt le contact. La différence avec un réseau social classique ? Ici, les « abonnements » sont des traités, les « commentaires » des résolutions, les « blocages » des sanctions, et les règles de fonctionnement sont le \cle{droit international}.

\begin{savaistu}[Du concert européen à la société mondiale]
Historiquement, les RI étaient l'affaire quasi exclusive des souverains (Traités de Westphalie, 1648 : principe de \cle{souveraineté} et d'\cle{non-ingérence}). Au XIX\textsuperscript{e} siècle, le « Concert européen » gère l'équilibre des puissances. La Société des Nations (1920) puis l'ONU (1945) institutionnalisent le \cle{multilatéralisme}. Depuis les années 1970, la \cle{mondialisation} fait entrer massivement des acteurs non étatiques (ONG, firmes, villes, réseaux scientifiques). Aujourd'hui, on parle de \emph{gouvernance mondiale} : l'État reste central, mais il n'est plus seul.
\end{savaistu}

\begin{propriete}[Fondement consensuel — \emph{pacta sunt servanda}]
Toute relation internationale repose sur le \textbf{principe du consentement mutuel} : \emph{pacta sunt servanda} (« les conventions doivent être respectées »). Un traité n'oblige que les parties qui l'ont ratifié ; une coutume lie ceux qui ne s'y sont pas opposés de manière persistante. C'est la base de la \cle{sécurité juridique} dans l'ordre international. Ce principe est consacré par l'article 26 de la Convention de Vienne sur le droit des traités (1969).
\end{propriete}

\begin{exemplebox}[Du traité à la start-up]
Les accords de coopération technique et culturelle entre le Cameroun et l'Allemagne prévoient des échanges d'étudiants, d'enseignants et d'experts. Ce sont des \cle{traités bilatéraux} au sens de la Convention de Vienne de 1969. L'ingénieur de notre accroche en est un bénéficiaire direct : son visa « chercheur », sa bourse du DAAD (organisme allemand d'échanges universitaires), la reconnaissance de son diplôme en Allemagne puis au Cameroun découlent de clauses précises de ces accords.
\end{exemplebox}

\begin{attention}[Diplomatie $\neq$ Relations internationales]
Ne confondez pas :
\begin{itemize}
\item \textbf{Diplomatie} : action extérieure de l'État (négociation, représentation, protection des ressortissants, information). C'est un \emph{moyen} étatique.
\item \textbf{Relations internationales} : champ d'étude et réalité empirique qui englobe la diplomatie \emph{et} l'action des OI, ONG, firmes, diasporas, villes, etc. C'est le \emph{champ} global.
\end{itemize}
\end{attention}

\courssub{Les acteurs du système international}

Le système international contemporain est \cle{multi-acteurs}. Identifions les cinq catégories majeures, leurs logiques et leurs interactions avec le Cameroun.

\begin{center}
\begin{tabularx}{\linewidth}{|l|X|X|X|}
\hline
\rowcolor{popblueL}
\textbf{Acteur} & \textbf{Nature / Statut} & \textbf{Logique d'action} & \textbf{Exemple camerounais} \\
\hline
\textbf{État} & Souverain, personnalité juridique de droit international pleine & Intérêt national, sécurité, puissance, bien-être de la population & Cameroun (Présidence, MINREX, ambassades) \\
\hline
\textbf{Organisations internationales (OI)} & Créées par traité, personnalité juridique propre & Mise en œuvre de mandats collectifs (paix, développement, normes) & ONU, UA, CEMAC, OIF, OMC, OMS, UNESCO \\
\hline
\textbf{ONG internationales} & Droit privé, but non lucratif, indépendance vis-à-vis des États & Plaidoyer, assistance, expertise, témoignage & MSF, Croix-Rouge, Amnesty International, Care \\
\hline
\textbf{Entreprises multinationales (EMN)} & Droit privé, but lucratif, capital multi-pays & Profit, accès aux ressources et marchés, stabilité réglementaire & Bolloré, Orange, TotalEnergies, Dangote, Huawei \\
\hline
\textbf{Collectivités territoriales} & Droit interne, personnalité morale, pas de souveraineté externe & Coopération décentralisée, attractivité, services locaux & Ville de Douala (jumelages avec des villes étrangères) \\
\hline
\end{tabularx}
\end{center}

\begin{popfigure}
\begin{tikzpicture}[
  node distance=2.2cm and 1.8cm,
  every node/.style={font=\small},
  state/.style={circle, draw=popblue, thick, fill=popblueL, minimum size=1.2cm, align=center},
  org/.style={rectangle, rounded corners, draw=popgreen, thick, fill=popgreenL, minimum width=2.2cm, minimum height=1cm, align=center},
  ong/.style={diamond, draw=poporange, thick, fill=poporangeL, aspect=1.5, minimum width=1.8cm, align=center},
  ent/.style={rectangle, draw=poppurple, thick, fill=poppurpleL, minimum width=2cm, minimum height=0.9cm, align=center},
  local/.style={ellipse, draw=poppink, thick, fill=poppinkL, minimum width=1.8cm, minimum height=0.9cm, align=center},
  arrow/.style={-Stealth, thick, popline},
  bidir/.style={Stealth-Stealth, thick, popdark}
]
% Nœuds
\node[state] (CM) {Cameroun};
\node[org, above left=of CM] (UA) {Union Africaine};
\node[org, above=of CM] (ONU) {ONU};
\node[org, above right=of CM] (CEMAC) {CEMAC};
\node[ent, left=of CM] (FR) {France};
\node[ent, right=of CM] (CN) {Chine};
\node[ong, below left=of CM] (MSF) {MSF};
\node[ong, below=of CM] (CICR) {CICR};
\node[local, below right=of CM] (Douala) {Ville de Douala};

% Liens bilatéraux (plein)
\draw[bidir] (CM) -- node[midway, above, sloped, font=\tiny] {Accords cadre, défense, culture} (FR);
\draw[bidir] (CM) -- node[midway, above, sloped, font=\tiny] {Infrastructures, mines, numérique} (CN);
\draw[bidir] (CM) -- node[midway, below, sloped, font=\tiny] {Coopération décentralisée} (Douala);

% Liens multilatéraux (pointillés)
\draw[arrow, densely dashed] (CM) -- node[midway, left, font=\tiny] {Contributions, votes} (ONU);
\draw[arrow, densely dashed] (CM) -- node[midway, left, font=\tiny] {Paix, Agenda 2063} (UA);
\draw[arrow, densely dashed] (CM) -- node[midway, right, font=\tiny] {Marché commun, monnaie} (CEMAC);

% Liens ONG
\draw[arrow] (MSF) -- node[midway, left, font=\tiny] {Santé d'urgence} (CM);
\draw[arrow] (CICR) -- node[midway, right, font=\tiny] {DIH, détention} (CM);
\end{tikzpicture}
\legende{Réseau d'interactions du Cameroun : traits pleins à double flèche = relations bilatérales directes ; traits pointillés = relations multilatérales via organisations ; losanges = ONG ; rectangles violets = partenaires économiques majeurs ; ellipse rose = collectivité territoriale.}
\end{popfigure}

\begin{methode}[Grille d'analyse d'une relation internationale : « Acteurs – Intérêts – Instruments – Résultats »]
Pour décrypter toute situation (accord, crise, sommet, projet), appliquez systématiquement ces quatre questions :
\begin{enumerate}
\item \textbf{Acteurs} : qui sont les parties prenantes ? (État, OI, ONG, firme, ville, individu). Sont-ils souverains ou non ?
\item \textbf{Intérêts} : que cherche chacun ? (Sécurité, profit, influence, normes, légitimité, développement).
\item \textbf{Instruments} : quels leviers mobilisent-ils ? (Traité, résolution, financement, sanction, expertise, médiatisation, réseau diaspora, standard technique).
\item \textbf{Résultats} : quels effets observables ? (Accord signé, flux financier, changement de loi, infrastructure construite, conflit gelé, norme adoptée).
\end{enumerate}
\end{methode}

\begin{exoresolu}[Analyse de l'adhésion du Cameroun à la ZLECAf]
\textbf{Énoncé :} Le Cameroun a signé en 2019 puis ratifié en 2020 l'accord instituant la Zone de libre-échange continentale africaine (ZLECAf), entrée en vigueur commerciale le 1\textsuperscript{er} janvier 2021. Utilisez la grille « Acteurs – Intérêts – Instruments – Résultats » pour analyser cette décision.
\tcblower
\textbf{Solution détaillée :}
\begin{itemize}
\item \textbf{Acteurs :} État du Cameroun (ratifiant), Union Africaine (dépositaire), États parties (marché unique d'environ 1,3 milliard de consommateurs), secteur privé camerounais (exportateurs, PME, agro-industrie), douanes, ministères du Commerce et de l'Économie.
\item \textbf{Intérêts :} pour le Cameroun : accès à un vaste marché continental, diversification des exportations hors pétrole et bois brut, attraction d'investissements, intégration régionale (complémentarité avec la CEMAC). Pour l'UA : crédibilité du projet panafricain, croissance du commerce intra-africain. Pour le secteur privé : démantèlement progressif des droits de douane, règles d'origine claires.
\item \textbf{Instruments :} Accord ZLECAf (2018), protocoles sur le commerce des biens et des services, mécanisme de règlement des différends, stratégie nationale d'implémentation, réduction progressive des droits de douane.
\item \textbf{Résultats :} démarrage effectif des échanges sous régime préférentiel (2021), premières exportations camerounaises sous le régime ZLECAf, formation d'opérateurs économiques aux règles d'origine ; défis persistants : infrastructures logistiques (port de Kribi, corridor Douala–N'Djamena), barrières non tarifaires, mise à niveau qualité des entreprises.
\end{itemize}
\end{exoresolu}

\courssub{Relations bilatérales et relations multilatérales}

\begin{definition}[Bilatéralisme]
Relation structurée \textbf{entre deux États} (ou un État et une OI, une ONG ou une firme) par des accords, des échanges et des représentations diplomatiques permanentes. C'est la forme la plus ancienne et la plus directe des relations internationales.
\end{definition}

\begin{definition}[Multilatéralisme]
Relation structurée \textbf{entre trois États ou plus}, généralement au sein d'une organisation internationale à vocation universelle ou régionale, selon des règles communes (charte, statuts, vote, secrétariat).
\end{definition}

\begin{center}
\begin{tabularx}{\linewidth}{|l|X|X|}
\hline
\rowcolor{popblueL}
\textbf{Critère} & \textbf{Bilatéral} & \textbf{Multilatéral} \\
\hline
\textbf{Nombre de parties} & 2 (principalement) & $\ge$ 3 (souvent dizaines ou centaines) \\
\hline
\textbf{Négociation} & Directe, souple, parfois confidentielle & Codifiée, transparente, par coalitions \\
\hline
\textbf{Engagement} & Sur mesure, plus aisément réversible & Normes communes s'imposant à tous les membres \\
\hline
\textbf{Exemples camerounais} & Accords Cameroun–France (défense, culture), Cameroun–Chine (infrastructures), Cameroun–Maroc (formation, banque) & ONU (maintien de la paix, ODD), UA (paix, ZLECAf), CEMAC (union monétaire), OMC (règles du commerce), Francophonie (langue, valeurs) \\
\hline
\textbf{Avantage} & Réactivité, adaptation aux besoins précis & Légitimité, mutualisation des coûts, normes universelles \\
\hline
\textbf{Limite} & Asymétrie de puissance, fragmentation & Lenteur, consensus difficile, « plus petit dénominateur commun » \\
\hline
\end{tabularx}
\end{center}

\begin{attention}[Le piège de l'exclusivité]
Le Cameroun ne « choisit pas » entre bilatéral et multilatéral : il \textbf{les combine}. Exemple : la formation d'ingénieurs camerounais en Allemagne (bilatéral) s'inscrit dans des programmes d'échanges universitaires (multilatéral) et respecte les normes de reconnaissance des diplômes portées par l'UNESCO (multilatéral universel). La diplomatie camerounaise pratique le \cle{multilatéralisme à géométrie variable} : active à l'UA pour la paix, à l'OMC pour le commerce, à la CEMAC pour la monnaie, bilatérale avec la Chine pour les routes, avec la France pour la défense et la culture.
\end{attention}

\courssub{Le rôle du droit international : traités, coutumes, résolutions}

Le droit international est la « grammaire » des relations internationales. Il transforme la force en droit et l'arbitraire en procédure.

\begin{propriete}[Sources du droit international (article 38 du Statut de la Cour internationale de Justice)]
\begin{enumerate}
\item \textbf{Conventions (traité, accord, convention, protocole, charte)} : engagements écrits explicites, consentement exprimé par signature puis ratification. \emph{Exemples :} Charte de l'ONU, Convention de Vienne sur le droit des traités (1969), Accord de Samoa (2023) entre l'Union européenne et les États d'Afrique, des Caraïbes et du Pacifique.
\item \textbf{Coutume internationale} : pratique générale et constante des États, acceptée comme droit (\emph{opinio juris}). \emph{Exemples :} immunité des chefs d'État en exercice, interdiction du génocide, liberté de la haute mer.
\item \textbf{Principes généraux de droit} : reconnus par les nations civilisées (bonne foi, équité, non-rétroactivité des lois).
\item \textbf{Décisions judiciaires et doctrine} : moyens subsidiaires d'interprétation (arrêts de la CIJ, travaux des juristes).
\end{enumerate}
\end{propriete}

\begin{exemplebox}[Hiérarchie et articulation : le cas des droits de l'homme au Cameroun]
Le Cameroun a ratifié la Charte africaine des droits de l'homme et des peuples (1989) et le Pacte international relatif aux droits civils et politiques (1984). Selon l'article 45 de la Constitution de 1996, les traités ratifiés ont, dès leur publication, une \textbf{autorité supérieure à celle des lois nationales}. La \textbf{coutume} (par exemple l'interdiction de la torture) lie le Cameroun indépendamment de toute ratification. Les \textbf{résolutions} du Conseil de sécurité de l'ONU adoptées sous le chapitre VII de la Charte s'imposent à tous les États membres.
\end{exemplebox}

\begin{methode}[Lire un traité : 5 repères rapides]
\begin{enumerate}
\item \textbf{Titre et préambule} : objet, parties, contexte politique.
\item \textbf{Champ d'application} : personnes, matières, territoires et durée concernés.
\item \textbf{Obligations substantielles} : droits et devoirs précis (verbes : « s'engage à », « veille à », « s'abstient de »).
\item \textbf{Mise en œuvre} : organes de suivi, rapports, conférence des parties, règlement des différends.
\item \textbf{Clauses finales} : signature, ratification, entrée en vigueur, réserves, dénonciation, révision, dépositaire, langues authentiques.
\end{enumerate}
\end{methode}

\courssub{Les indicateurs de mesure des relations internationales}

Comment « mesurer » l'intensité des relations d'un pays ? Les statisticiens et les diplomates utilisent des \cle{indicateurs quantitatifs} et \cle{qualitatifs}.

\begin{center}
\begin{tabularx}{\linewidth}{|l|X|X|}
\hline
\rowcolor{popblueL}
\textbf{Indicateur} & \textbf{Description} & \textbf{Exemple camerounais (ordres de grandeur récents)} \\
\hline
\textbf{Échanges commerciaux} & Valeur des importations et exportations, part des partenaires, solde de la balance commerciale & Exportations $\approx$ 3 000 milliards FCFA ; importations $\approx$ 4 000 milliards FCFA ; principaux clients : Pays-Bas, Chine, Inde, Italie, Espagne ; principaux fournisseurs : Chine, France, Inde, Nigeria — \emph{INS, Douanes} \\
\hline
\textbf{Flux d'IDE (investissements directs étrangers)} & Stocks et flux annuels, secteurs, pays d'origine & Flux annuels de quelques centaines de milliards FCFA (hydrocarbures, télécommunications, bois, agro-industrie) — \emph{BEAC, CNUCED} \\
\hline
\textbf{Accords signés / ratifiés} & Nombre, type (bilatéral / multilatéral), domaines & Accords bilatéraux (défense, éducation, santé, fiscalité) ; participation aux sommets multilatéraux (UA, ONU, Francophonie, Commonwealth, CEMAC, COP) — \emph{MINREX} \\
\hline
\textbf{Flux migratoires / mobilité} & Diaspora, étudiants, travailleurs, transferts d'argent & Diaspora estimée à plusieurs centaines de milliers de personnes ; transferts de la diaspora de l'ordre de plusieurs centaines de milliards FCFA par an — \emph{BEAC, Banque mondiale} \\
\hline
\textbf{Participation aux OI} & Contributions financières, personnel (casques bleus, experts), votes & Contingents de casques bleus dans les missions de l'ONU en Afrique centrale ; contributions aux budgets ONU/UA/CEMAC — \emph{MINREX, MINDEF} \\
\hline
\textbf{Coopération décentralisée} & Jumelages, projets ville-à-ville, financement local & Nombreux jumelages actifs (Douala, Yaoundé, Bafoussam avec des villes françaises, allemandes, chinoises, etc.) — \emph{CVUC, AIMF} \\
\hline
\end{tabularx}
\end{center}

\begin{popfigure}
\begin{tikzpicture}
\begin{axis}[
  width=\linewidth,
  height=9cm,
  ybar=2pt,
  bar width=10pt,
  ymin=0,
  ymax=1000,
  ylabel={Valeur des échanges (milliards FCFA)},
  symbolic x coords={Chine,France,Inde,Nigeria,Pays-Bas,Italie,Espagne},
  xtick=data,
  xticklabel style={rotate=45, anchor=east, font=\tiny},
  ymajorgrids=true,
  grid style=dashed,
  legend style={at={(0.5,-0.28)}, anchor=north, legend columns=2, font=\small},
  title style={font=\normalsize\bfseries},
  title={Principaux partenaires commerciaux du Cameroun (ordres de grandeur, INS/Douanes)},
  popbar1/.style={fill=popblue},
  popbar2/.style={fill=poporange},
]
\addplot[popbar1] coordinates {
  (Chine,900) (France,600) (Inde,450) (Nigeria,400) (Pays-Bas,250) (Italie,220) (Espagne,200)
};
\addlegendentry{Importations (Mrd FCFA)}

\addplot[popbar2] coordinates {
  (Chine,300) (France,200) (Inde,350) (Nigeria,150) (Pays-Bas,500) (Italie,350) (Espagne,400)
};
\addlegendentry{Exportations (Mrd FCFA)}
\end{axis}
\end{tikzpicture}
\legende{Échanges commerciaux du Cameroun par partenaire principal (ordres de grandeur). Barres bleues = importations ; barres orange = exportations. On observe un excédent commercial avec les Pays-Bas, l'Italie et l'Espagne (pétrole brut, cacao, bois, banane) et un déficit structurel vis-à-vis de la Chine et de la France (biens d'équipement, produits manufacturés). Source : INS, Douanes, BEAC.}
\end{popfigure}

\begin{aretenir}[Synthèse : ce qu'il faut retenir des Relations Internationales]
\begin{itemize}
\item Les \textbf{RI} sont l'ensemble des interactions normées entre acteurs étatiques et non étatiques.
\item L'\textbf{État} reste l'acteur central (souveraineté, siège à l'ONU), mais les \textbf{OI, ONG, EMN et collectivités territoriales} sont incontournables.
\item \textbf{Bilatéral} (2 parties, souplesse) et \textbf{multilatéral} ($\ge$ 3 parties, normes communes) sont complémentaires, non exclusifs.
\item Le \textbf{droit international} (traités, coutumes, résolutions) fournit le cadre juridique contraignant : \emph{pacta sunt servanda}.
\item On mesure l'intensité des RI par les \textbf{échanges commerciaux, les IDE, les accords, les mobilités humaines, la participation aux OI et la coopération décentralisée}.
\item Le Cameroun pratique une \textbf{diplomatie tous azimuts} : UA/CEMAC (régional), ONU/Francophonie/Commonwealth (universel), partenariats stratégiques bilatéraux (France, Chine, États-Unis, Maroc, Turquie, Russie, Brésil, Inde, etc.).
\end{itemize}
\end{aretenir}

\begin{exercice}[Application : analyse d'un accord bilatéral récent]
Le Cameroun et le Maroc ont signé un \textbf{accord-cadre de coopération dans le domaine de la formation professionnelle et de l'emploi}. Vous êtes chargé(e) de rédiger une note de synthèse pour le ministère en charge de la formation professionnelle.
\begin{enumerate}
\item Identifiez les \textbf{acteurs} (étatiques, non étatiques) impliqués directement ou indirectement.
\item Qualifiez la nature de l'instrument juridique (traité, accord simplifié, mémorandum ?) et son fondement (consentement).
\item Listez 3 \textbf{intérêts} probables pour le Cameroun et 3 pour le Maroc.
\item Proposez 2 \textbf{indicateurs de suivi} pour évaluer l'impact de cet accord à 3 ans.
\end{enumerate}
\end{exercice}

\begin{corrige}[Exercice n°1]
\begin{enumerate}
\item \textbf{Acteurs :} États (Cameroun et Maroc, via leurs ministères de la formation professionnelle et de l'emploi), ambassades, agences nationales pour l'emploi, centres de formation professionnelle publics et privés, secteur privé (patronats des deux pays), diasporas, collectivités territoriales.
\item \textbf{Instrument :} \emph{accord-cadre}, forme simplifiée (entrée en vigueur souvent dès la signature, sans ratification parlementaire), mais engagement politique fort. Fondement : \emph{pacta sunt servanda} et principe de la coopération Sud-Sud (Charte de l'ONU, Agenda 2063 de l'UA).
\item \textbf{Intérêts du Cameroun :} (1) accès à l'expertise marocaine en formation par alternance, (2) insertion des jeunes diplômés dans les filières porteuses (BTP, automobile, agro-industrie, numérique, tourisme), (3) transfert de méthodes de certification et de qualification. \textbf{Intérêts du Maroc :} (1) partage de son expertise en formation professionnelle, (2) accès à une main-d'œuvre qualifiée pour ses chantiers en Afrique, (3) rayonnement diplomatique et économique en Afrique centrale.
\item \textbf{Indicateurs à 3 ans :} (1) nombre de jeunes camerounais formés et certifiés dans le cadre de l'accord ; (2) taux d'insertion professionnelle des bénéficiaires six mois après la formation ; on peut ajouter le nombre de conventions signées entre écoles et centres des deux pays et le volume des investissements marocains dans la formation professionnelle au Cameroun.
\end{enumerate}
\end{corrige}

\coursec{Aspects et formes de la coopération internationale}

Après avoir défini les \cle{Relations Internationales} comme l'ensemble des interactions structurées entre acteurs étatiques et non étatiques sur la scène mondiale, nous nous penchons maintenant sur l'un de leurs moteurs principaux : la \cle{coopération internationale}. Si les relations internationales décrivent le « comment » les acteurs se parlent, se confrontent ou s'ignorent, la coopération en décrit le « faire ensemble » volontaire. C'est le passage de la logique de puissance à la logique de partenariat, même si, comme nous le verrons, les deux logiques s'entremêlent souvent.

\courssub{Définition et fondements de la coopération internationale}

\begin{definition}[Coopération internationale]
La \textbf{coopération internationale} est l'ensemble d'actions concertées entre acteurs (États, organisations internationales, collectivités territoriales, entreprises, ONG) pour atteindre des objectifs communs d'intérêt général, fondés sur les principes de souveraineté, d'égalité, de non-ingérence et de bénéfice mutuel.
\end{definition}

Contrairement à l'aide humanitaire d'urgence (réactive, ponctuelle, vitale), la coopération internationale s'inscrit dans la \textbf{durée} et vise le \textbf{développement structurel} (infrastructures, formation, réforme institutionnelle, transfert de technologie). Elle repose sur un \cle{contrat} (accord-cadre, convention, protocole d'accord) qui engage les parties sur des obligations réciproques : financement, expertise, mise à disposition de personnel, accès aux marchés, contreparties locales.

\begin{exemplebox}[Du concept à la réalité : l'accord-cadre Cameroun–Union Européenne]
L'Accord de Cotonou (2000) puis l'Accord de Samoa (2023) entre l'UE et les pays ACP (Afrique, Caraïbes, Pacifique) illustrent la structure type d'un cadre de coopération :
\begin{itemize}
    \item \textbf{Dialogue politique} : droits de l'homme, gouvernance démocratique, paix.
    \item \textbf{Coopération au développement} : secteurs prioritaires (éducation, santé, rural, infrastructures).
    \item \textbf{Commerce} : Accords de Partenariat Économique (APE) pour l'accès aux marchés.
    \item \textbf{Mise en œuvre} : Fonds Européen de Développement (FED), programmation pluriannuelle (7 ans), indicateurs de résultats.
\end{itemize}
Cet exemple montre que la coopération n'est pas un « don » gratuit : elle lie l'aide à des \cle{conditionnalités} (gouvernance, droits humains, gestion saine des finances publiques).
\end{exemplebox}

\courssub{Les grands secteurs (aspects) de la coopération}

La coopération se déploie dans des secteurs techniques précis, correspondant aux besoins du pays bénéficiaire et aux expertises du pays partenaire.

\begin{center}
\begin{tabularx}{\linewidth}{|l|X|X|}\hline
\rowcolor{popblueL}
\textbf{Secteur} & \textbf{Objectifs types} & \textbf{Exemples camerounais} \\ \hline
\textbf{Technique \& Scientifique} & Transfert de savoir-faire, formation de cadres, recherche appliquée, normalisation. & Projets \textit{IRAD/CIRAD} (variétés de manioc, cacao) ; formation d'ingénieurs à l'ENSPY via partenariats \textit{Écoles Centrales} ; appui à l'ANOR (normes). \\ \hline
\textbf{Culturel \& Éducatif} & Promotion des langues, patrimoine, mobilité étudiante, reconnaissance des diplômes. & Réseau des Alliances Françaises ; Campus France ; bourses d'excellence \textit{Eiffel}, \textit{Chine}, \textit{Turquie} ; jumelages universitaires (Univ. Yaoundé I – Sorbonne). \\ \hline
\textbf{Économique \& Financier} & Appui budgétaire, développement du secteur privé, commerce, investissement, PPP. & Appui budgétaire sectoriel (éducation/santé) Banque Mondiale/FMI ; guichet unique création d'entreprises (CFCE) ; Zone de Libre-Échange Continentale (ZLECAf). \\ \hline
\textbf{Militaire \& Sécuritaire} & Formation, équipement, renseignement, maintien de la paix, lutte terrorisme/piraterie. & Coopération \textit{Cameroun–USA} (formation BIR, drones), \textit{Cameroun–France} (écoles de guerre), \textit{Cameroun–Russie/Chine} (équipements) ; Force Multinationale Mixte (FMM) Lac Tchad. \\ \hline
\textbf{Gouvernance \& État de droit} & Réforme administration, justice, décentralisation, lutte corruption, statistiques. & Programme d'Appui à la Réforme de la Justice (PARJ) ; appui à la Cour des Comptes ; stratégie nationale de développement (SND30) financée par partenaires. \\ \hline
\end{tabularx}
\end{center}

\begin{attention}[Piège fréquent : confondre secteur et modalité]
Un même projet peut croiser plusieurs secteurs. Exemple : la construction d'un hôpital (infrastructure \textbf{économique}) équipé par un don (coopération \textbf{technique}) géré par du personnel formé (coopération \textbf{humaine}) dans un cadre de décentralisation (gouvernance). Au Probatoire, précisez toujours le secteur \textbf{dominant} du projet analysé.
\end{attention}

\courssub{Les formes juridiques et institutionnelles : bilatérale vs multilatérale}

La distinction fondamentale porte sur le \textbf{nombre de parties} signataires de l'accord-cadre.

\courssub{La coopération bilatérale : relation État à État}

C'est la forme la plus ancienne, la plus visible politiquement, souvent porteuse d'affinités historiques, linguistiques ou stratégiques.

\begin{propriete}[Caractéristiques de la coopération bilatérale]
\begin{itemize}
    \item \textbf{Négociation directe} : souplesse, adaptation au contexte national, alignement sur les priorités du bénéficiaire (principe de \cle{propriete} de l'aide).
    \item \textbf{Visibilité politique} : visites d'État, commissions mixtes, symbolique forte (ex. : don d'ambassade, hôpital de l'amitié).
    \item \textbf{Risque d'instrumentalisation} : l'aide peut servir d'effet de levier diplomatique (vote à l'ONU, alignement géopolitique, accès aux ressources).
\end{itemize}
\end{propriete}

\begin{exemplebox}[Deux modèles bilatéraux contrastés : France vs Chine]
\begin{center}
\begin{tabularx}{\linewidth}{|l|X|X|}\hline
\rowcolor{popblueL}
\textbf{Critère} & \textbf{Cameroun – France (Partenaire historique)} & \textbf{Cameroun – Chine (Partenaire émergent)} \\ \hline
\textbf{Cadre juridique} & Accords de coopération 1960, renégociés ; Cotonou/Samoa (via UE). & Forum sur la Coopération Sino-Africaine (FOCAC) ; Commission Mixte. \\ \hline
\textbf{Modalité dominante} & Dons (C2D – Contrat de Désendettement et Développement), expertise technique, bourses. & Prêts concessionnels (Exim Bank), infrastructures « clé en main », lignes de crédit. \\ \hline
\textbf{Secteurs phares} & Gouvernance, éducation, santé, agriculture, défense. & Infrastructures lourdes (barrages, routes, stades, bâtiments administratifs), mines, énergie. \\ \hline
\textbf{Conditionnalités} & Bonne gouvernance, droits de l'homme, transparence budgétaire. \emph{(Explicites)} & Non-ingérence, principe « Un seul Chine », viabilité économique du projet. \emph{(Implicites)} \\ \hline
\textbf{Exemple emblématique} & Programme d'Appui à la Décentralisation (PAD) ; Lycée Leclerc. & Barrage de Memve'ele ; Stade d'Olembé ; Route Kribi-Lolabé ; Assemblée Nationale. \\ \hline
\end{tabularx}
\end{center}
\end{exemplebox}

\courssub{La coopération multilatérale : l'action par les organisations internationales}

Elle passe par des institutions où le Cameroun est membre contributeur et bénéficiaire. Elle mutualise les ressources et impose des normes globales.

\begin{propriete}[La coopération multilatérale implique souvent des conditionnalités de bonne gouvernance]
Les institutions de Bretton Woods (Banque Mondiale, FMI) et les banques régionales (BAD, BEAC/CEMAC) conditionnent leurs décaissements à la mise en œuvre de \textbf{réformes structurelles} : assainissement finances publiques, libéralisation secteurs-clés, cadre des dépenses publiques, lutte contre la corruption, statistiques fiables. Ces conditionnalités visent l'efficacité de l'aide mais soulèvent des débats sur la \cle{souveraineté} économique.
\end{propriete}

\begin{center}
\begin{tabularx}{\linewidth}{|l|X|X|}\hline
\rowcolor{popblueL}
\textbf{Institution} & \textbf{Mandat / Outils} & \textbf{Interventions récentes au Cameroun} \\ \hline
\textbf{Banque Mondiale (IDA/IBRD)} & Réduction pauvreté, investissement capital humain, climat. Prêts/Don IDA (pays à revenu intermédiaire). & Projet \textit{PISA} (santé), \textit{PDRL} (développement rural), \textit{PAJER} (emploi jeunes), \textit{DPO} (appui budgétaire réformes). \\ \hline
\textbf{FMI} & Stabilité macroéconomique, viabilité dette, taux de change. Facilités (FEC, MEC, FCL). & Programme 2021-2024 (FEC/Mécanisme élargi de crédit) : 483 Mds FCFA ; réformes : TVA, subventions carburant, gouvernance SNH/SONARA. \\ \hline
\textbf{BAD (Banque Africaine de Développement)} & Infrastructures régionales, intégration, secteur privé, climat. & Corridor Douala–Ndjamena–Bangui ; Barrage de Nachtigal (cofinancement) ; Appui ZLECAf ; Projet d'eau Yaoundé. \\ \hline
\textbf{Agences ONU (PNUD, UNICEF, OMS, FAO, PAM)} & Normes mondiales, aide technique, urgence/développement, ODD. & PNUD : gouvernance locale, climat ; UNICEF : éducation/vaccination zones crises (NOSO, Extrême-Nord) ; OMS : riposte COVID/Choléra ; PAM : cantines scolaires, résilience. \\ \hline
\textbf{Fonds Vert / FEM} & Climat, biodiversité, désertification. & Projet \textit{REDD+} (forêts), adaptation agriculture climat-intelligente, gestion bassins versants. \\ \hline
\end{tabularx}
\end{center}

\begin{exoresolu}[Analyse d'un document : Extrait de conclusion de la revue Article IV FMI 2023 – Cameroun]
\textbf{Document :} « Les directeurs ont salué les progrès dans la mise en œuvre du programme appuyé par la FEC... Ils ont insisté sur la nécessité de poursuivre la mobilisation des recettes intérieures, d'améliorer la transparence du secteur pétrolier (SNH) et de protéger les dépenses sociales. Ils ont souligné les risques liés à l'insécurité et aux chocs climatiques. »

\textbf{Questions :}
1. Identifier le type de coopération et l'institution émettrice.
2. Citer trois conditionnalités explicites ou implicites.
3. Expliquer le lien entre « mobilisation recettes intérieures » et « protection dépenses sociales ».

\textbf{Solution détaillée :}
1. \textbf{Coopération multilatérale}, institution : \textbf{Fonds Monétaire International (FMI)}. C'est une surveillance multilatérale (Article IV) dans le cadre d'un programme de prêt (FEC).
2. \textbf{Conditionnalités :} (a) Mobilisation recettes intérieures (élargir assiette fiscale, réduire exonérations) ; (b) Transparence secteur pétrolier (comptes SNH, contrats) ; (c) Protection dépenses sociales (plancher dépenses santé/éducation) ; (d) Viabilité dette (plafond endettement non concessionnel).
3. \textbf{Lien :} Le Cameroun a un taux de pression fiscale faible (~13\% PIB). Pour financer les dépenses sociales (santé, éducation, filets sociaux) sans creuser le déficit ni s'endetter excessivement, il \textbf{faut} augmenter les recettes propres. La conditionnalité FMI lie donc \textbf{autonomie financière} et \textbf{protection sociale}.
\end{exoresolu}

\courssub{La coopération décentralisée : le pouvoir local au service du développement}

C'est une forme originale où les \textbf{collectivités territoriales décentralisées} (CTD : communes, régions) coopèrent directement avec leurs homologues étrangers, sans passer obligatoirement par l'État central. Elle est encadrée par la \textbf{Loi n°2004/017 du 22 juillet 2004} relative à l'orientation de la décentralisation (art. 30 : « Les communes et les régions peuvent conclure des conventions de jumelage... »).

\begin{savaistu}[Le programme « Villes jumelées » finance chaque année 15 projets locaux de développement urbain]
Le Ministère de l'Habitat et du Développement Urbain (MINHDU) et l'Association des Communes et Villes Unies du Cameroun (CVUC) pilotent un appel à projets annuel. Les jumelages actifs financent : assainissement, marchés municipaux, éclairage solaire, centres de formation professionnelle, gestion déchets, parcs urbains. Exemple : le jumelage \textbf{Douala–Lyon} a permis la création de l'Agence Locale de l'Énergie et du Climat (ALEC) de Douala et la réhabilitation du Marché Central.
\end{savaistu}

\begin{exemplebox}[Principaux jumelages décentralisés camerounais (échantillon)]
\begin{center}
\begin{tabularx}{\linewidth}{|l|l|X|}\hline
\rowcolor{popblueL}
\textbf{Commune / Région (Cameroun)} & \textbf{Partenaire étranger} & \textbf{Thématiques prioritaires} \\ \hline
\textbf{Douala} (Ville) & \textbf{Lyon} (France) & Développement urbain durable, énergie, climat, culture, jeunesse. \\ \hline
\textbf{Yaoundé} (Ville) & \textbf{Paris} (France) & Assainissement, mobilité urbaine, patrimoine, numérique. \\ \hline
\textbf{Bafoussam} (Commune) & \textbf{Reims} (France) & Eau potable, assainissement, formation professionnelle. \\ \hline
\textbf{Maroua} (Commune) & \textbf{Montreuil} (France) & Gestion déchets, agriculture urbaine, santé maternelle. \\ \hline
\textbf{Limbe} (Commune) & \textbf{Plymouth} (Royaume-Uni) & Environnement côtier, tourisme, éducation. \\ \hline
\textbf{Région du Centre} & \textbf{Île-de-France} (France) & Formation professionnelle, lycées techniques, santé. \\ \hline
\textbf{Kribi} (Commune) & \textbf{Qingdao} (Chine) & Port, pêche, infrastructures, tourisme. \\ \hline
\end{tabularx}
\end{center}
\end{exemplebox}

\begin{popfigure}
\begin{tikzpicture}[scale=0.85, every node/.style={font=\small}]
% Contour très simplifié du Cameroun (polygone approximatif)
\draw[popdark, thick, fill=popcream] 
(0,0) -- (8,0) -- (9,2) -- (8.5,5) -- (7,7) -- (4,7.5) -- (2,6) -- (0.5,4) -- (0,1) -- cycle;
% Villes camerounaises (points rouges)
\foreach \x/\y/\name in {
    3.5/2.5/Douala,
    4.5/4.5/Yaoundé,
    5.5/5.5/Bafoussam,
    7/5.5/Maroua,
    2.5/3.5/Limbé,
    5.5/2.5/Kribi,
    4.5/5.0/RégionCentre} {
    \fill[popred] (\x,\y) circle (3pt);
    \node[anchor=west, popdark] at (\x+0.15,\y) {\name};
}
% Flèches vers partenaires (positionnés hors carte pour lisibilité)
% Douala -> Lyon (Nord-Ouest)
\draw[popblue, thick, ->, >=Stealth] (3.5,2.5) .. controls (0,4) and (-2,6) .. (-3,7) node[left, popblue, align=right] {Lyon (France)\\(Jumelage 1987)};
% Yaoundé -> Paris (Nord)
\draw[popblue, thick, ->, >=Stealth] (4.5,4.5) .. controls (4.5,6) and (4.5,8) .. (4.5,9) node[above, popblue, align=center] {Paris (France)\\(Jumelage 1985)};
% Bafoussam -> Reims (Nord-Nord-Ouest)
\draw[popblue, thick, ->, >=Stealth] (5.5,5.5) .. controls (4,7) and (1,8) .. (0,9) node[left, popblue, align=right] {Reims (France)\\(Eau/Assainissement)};
% Maroua -> Montreuil (Nord-Ouest)
\draw[popblue, thick, ->, >=Stealth] (7,5.5) .. controls (5,7) and (2,8) .. (-1,9) node[left, popblue, align=right] {Montreuil (France)\\(Déchets/Santé)};
% Limbé -> Plymouth (Nord-Ouest lointain)
\draw[popblue, thick, ->, >=Stealth] (2.5,3.5) .. controls (0,5) and (-3,7) .. (-4,8) node[left, popblue, align=right] {Plymouth (UK)\\(Environnement)};
% Kribi -> Qingdao (Nord-Est symbolique, hors carte)
\draw[poporange, thick, ->, >=Stealth] (5.5,2.5) .. controls (8,3) and (10,4) .. (11,5) node[right, poporange, align=left] {Qingdao (Chine)\\(Port/Pêche)};
% Région Centre -> Île-de-France
\draw[popblue, thick, ->, >=Stealth] (4.5,5.0) .. controls (3,6.5) and (0,7.5) .. (-2,8.5) node[left, popblue, align=right] {Île-de-France (France)\\(Formation/Santé)};
% Légende
\node[anchor=south west, draw=popline, fill=white, rounded corners, inner sep=4pt] at (0.2,0.2) {
    \textbf{Légende :}\par
    \textcolor{popred}{●} Communes/Régions camerounaises\par
    \textcolor{popblue}{→} Jumelages avec la France / Europe\par
    \textcolor{poporange}{→} Jumelages avec la Chine / Asie\par
    \textit{Flèches : sens de la coopération (projets, expertise, financements)}
};
% Flèche Nord
\draw[popdark, ->, >=Stealth, thick] (8.5,7) -- (8.5,7.8) node[above] {N};
\end{tikzpicture}
\legende{Carte schématique des principaux jumelages décentralisés du Cameroun (échantillon 2024). Les flèches bleues indiquent les partenariats historiques avec la France et l'Europe ; la flèche orange le partenariat émergent avec la Chine. La coopération décentralisée permet un transfert direct de compétences techniques (gestion urbaine, eau, déchets) entre pairs élus.}
\end{popfigure}

\courssub{Les Partenariats Public-Privé (PPP) : un levier pour les infrastructures}

Le PPP n'est pas de l'aide publique au développement (APD). C'est un \textbf{contrat de long terme} entre une personne publique (État, collectivité) et un partenaire privé pour la conception, le financement, la construction et l'exploitation d'un ouvrage ou service public. Le partenaire privé supporte des risques (construction, exploitation, demande) et est rémunéré par des \textbf{redevances utilisateur} (péage, facture eau/électricité) et/ou des \textbf{paiements publics} (disponibilité, loyer).

\begin{attention}[Les PPP ne remplacent pas l'aide publique au développement (APD) mais la complètent]
\begin{itemize}
    \item \textbf{APD} : dons/prêts concessionnels (taux < 2\%, long différé), sectors sociaux (santé, éducation), non rentables marchandement. Objectif : réduction pauvreté, ODD.
    \item \textbf{PPP} : investissement privé rentable, secteurs « bankables » (énergie, transport, eau, numérique). Objectif : efficacité, transfert risque, levier financier.
    \item \textbf{Danger} : présenter un PPP comme « don » ou « aide » masque le coût réel pour le budget public (engagements hors dette, redevances sur 20-25 ans) et le risque de renchérissement du service pour l'usager.
\end{itemize}
\end{attention}

\begin{exemplebox}[Le barrage de Nachtigal : PPP de référence en Afrique centrale]
\begin{itemize}
    \item \textbf{Projet :} Aménagement hydroélectrique 420 MW (centrale + ligne de transport 50 km), coût total ~1 200 Mrd FCFA (1,8 Md €).
    \item \textbf{Structure PPP (Build-Own-Operate-Transfer - BOOT 35 ans) :}
    \begin{itemize}
        \item \textbf{Société de Projet (SPV) :} \textit{Nachtigal Hydro Power Company (NHPC)}.
        \item \textbf{Actionnaires :} EDF (France) 40\%, État du Cameroun 15\% (via SONATREL), IFC (Banque Mondiale) 40\%.
        \item \textbf{Financement :} 75\% dette (prêts BAD, AFD, Proparco, banques commerciales), 25\% fonds propres.
        \item \textbf{Contrat d'achat d'électricité (PPA) :} ENEO (distributeur) s'engage à acheter l'énergie à un tarif niveau (levelized cost) compétitif vs thermique.
    \end{itemize}
    \item \textbf{Apport APD complémentaire :} Don AFD 50 M€ pour études environnementales/sociales ; Garantie partielle de risque MIGA (Banque Mondiale) pour rassurer prêteurs privés.
    \item \textbf{Impact :} +30\% capacité installée nationale ; coût kWh divisé par 3 vs fuel ; réduction 2,5 Mt CO2/an ; 1 500 emplois construction.
\end{itemize}
\end{exemplebox}

\courssub{Conditionnalités, évaluation d'impact et alignement sur les ODD}

Toute coopération moderne (bilatérale, multilatérale, décentralisée, PPP) est soumise à une double exigence : \textbf{conditionnalité} à l'entrée (éligibilité) et \textbf{évaluation} à la sortie (résultats).

\begin{methode}[Construire un diagramme en secteurs (camembert) pour visualiser la répartition sectorielle de l'aide]
\textbf{Objectif :} Transformer un tableau de données brutes (montants par secteur) en graphique lisible pour l'oral ou l'écrit.
\begin{enumerate}
    \item \textbf{Collecter les données officielles} : Rapports annuels MINREX, MINEPAT, OCDE/DAC, ou rapports pays Banque Mondiale/FMI. Unités : millions USD ou milliards FCFA. Année de référence unique.
    \item \textbf{Calculer le total} : Somme des montants tous secteurs confondus.
    \item \textbf{Calculer le pourcentage par secteur} : $\frac{\text{Montant secteur}}{\text{Total}} \times 100$.
    \item \textbf{Convertir en degrés} : $\text{Pourcentage} \times 3,6$ (car $360^\circ / 100 = 3,6$).
    \item \textbf{Tracer le camembert} : Compas + rapporteur (papier) ou Excel / pgfplots (numérique). Ordonner les parts du plus grand au plus petit dans le sens horaire. Légender chaque part (secteur + \% + montant). Titre précis : « Répartition sectorielle de l'APD reçue par le Cameroun en 2022 (Source : OCDE/DAC) ».
    \item \textbf{Analyser} : Quel secteur domine ? Y a-t-il alignement sur la SND30 ? Où sont les lacunes ?
\end{enumerate}
\end{methode}

\begin{exemplebox}[En 2022, l'APD reçue par le Cameroun s'élevait à 1,2 milliard USD, dont 45 \% en coopération technique]
Source : OCDE/DAC (statistiques 2022, décaissements bruts).
\begin{center}
\begin{tabular}{|l|c|c|}\hline
\rowcolor{popblueL}
\textbf{Secteur (Code OCDE)} & \textbf{Montant (M USD)} & \textbf{\% Total} \\ \hline
Action humanitaire (urgence) & 180 & 15,0 \% \\ \hline
Gouvernance / Société civile & 156 & 13,0 \% \\ \hline
Éducation & 132 & 11,0 \% \\ \hline
Santé / Population & 120 & 10,0 \% \\ \hline
Eau / Assainissement & 96 & 8,0 \% \\ \hline
Autres secteurs (Agriculture, Énergie, Transport, Multisectoriel) & 516 & 43,0 \% \\ \hline
\textbf{Total APD} & \textbf{1 200} & \textbf{100 \%} \\ \hline
\end{tabular}
\end{center}
\textbf{Note :} La « coopération technique » (experts, formations, études, volontariat) est souvent \textbf{imputée transversalement} dans chaque secteur (ex. : expert santé compté dans « Santé »). Les 45\% correspondent à la part des dépenses de personnel/études dans le total décaissé.
\end{exemplebox}

\begin{popfigure}
\begin{tikzpicture}
\begin{axis}[
    ybar stacked,
    bar width=18pt,
    width=\linewidth,
    height=9cm,
    enlargelimits=0.15,
    legend style={at={(0.5,-0.25)}, anchor=north, legend columns=-1, font=\footnotesize, /tikz/every even column/.append style={column sep=5pt}},
    ylabel={Montant (Millions USD)},
    ylabel style={font=\small},
    xlabel={Années},
    xlabel style={font=\small},
    symbolic x coords={2018,2019,2020,2021,2022,2023},
    xtick=data,
    xticklabel style={font=\small, rotate=45, anchor=east},
    ymin=0, ymax=1400,
    ymajorgrids=true,
    grid style={popline},
    axis lines=left,
    tick style={draw=popline},
]
% Données simulées cohérentes avec tendances réelles (OCDE/DAC Cameroun)
% Humanitaire
\addplot[fill=popred!70] coordinates {(2018,120) (2019,140) (2020,200) (2021,180) (2022,180) (2023,160)};
% Gouvernance
\addplot[fill=popblue!70] coordinates {(2018,100) (2019,110) (2020,120) (2021,140) (2022,156) (2023,150)};
% Education
\addplot[fill=popgreen!70] coordinates {(2018,90) (2019,95) (2020,100) (2021,110) (2022,132) (2023,130)};
% Santé
\addplot[fill=poporange!70] coordinates {(2018,80) (2019,85) (2020,150) (2021,130) (2022,120) (2023,110)};
% Eau/Assainissement
\addplot[fill=poppurple!70] coordinates {(2018,60) (2019,65) (2020,70) (2021,80) (2022,96) (2023,90)};
% Autres (Infra, Agri, Énergie, Multi)
\addplot[fill=popgold!70] coordinates {(2018,450) (2019,480) (2020,400) (2021,500) (2022,516) (2023,550)};
\legend{Humanitaire, Gouvernance, Éducation, Santé, Eau/Assainissement, {Autres (Infra, Agri, Énergie)}}
\end{axis}
\end{tikzpicture}
\legende{Évolution des décaissements d'APD au Cameroun par grand secteur (2018–2023, estimations OCDE/DAC). On observe le pic humanitaire/santé en 2020 (COVID-19 + crises NOSO/Extrême-Nord), la montée structurelle de la gouvernance et de l'éducation, et la part dominante mais volatile des « Autres » (grands projets d'infrastructures, énergie, agriculture). Source : OCDE/DAC, tableaux statistiques pays, données 2024.}
\end{popfigure}

\begin{propriete}[Alignement sur les ODD : la grille de lecture universelle]
Depuis 2015, toute coopération (don, prêt, PPP, décentralisée) est évaluée à l'aune des \textbf{17 Objectifs de Développement Durable (ODD)}. Les indicateurs clés pour le Cameroun (SND30) :
\begin{itemize}
    \item \textbf{ODD 1 (Pauvreté) / ODD 2 (Faim) :} Taux de pauvreté, sécurité alimentaire, filets sociaux.
    \item \textbf{ODD 3 (Santé) / ODD 4 (Éducation) :} Mortalité maternelle/infantile, taux de scolarisation/achèvement, qualité.
    \item \textbf{ODD 5 (Genre) :} Taux de représentation femmes, violences basées genre, accès terre/crédit.
    \item \textbf{ODD 6 (Eau) / ODD 7 (Énergie) :} Accès eau potable/assainissement, taux d'électrification, part renouvelables.
    \item \textbf{ODD 8 (Emploi) / ODD 9 (Industrialisation) :} Taux chômage jeunes, part industrie dans PIB, infrastructures transport/numérique.
    \item \textbf{ODD 13 (Climat) / ODD 15 (Vie terrestre) :} Couverture forestière, émissions CO2, adaptation agriculture.
    \item \textbf{ODD 16 (Paix/Justice/Institutions) :} Indice corruption, efficacité justice, état civil, statistiques.
    \item \textbf{ODD 17 (Partenariats) :} Volume APD, investissements directs étrangers (IDE), transfert technologie, commerce.
\end{itemize}
L'évaluation d'impact (mid-term review, completion report) mesure la \textbf{contribution} (pas l'attribution unique) du projet à ces indicateurs.
\end{propriete}

\begin{exoresolu}[Mise en situation : Rédiger une fiche d'analyse d'un projet de coopération]
\textbf{Consigne :} À partir du projet « Appui à la gouvernance locale dans la Région de l'Extrême-Nord » (financé UE, mis en œuvre par GIZ, 5 M€, 2021-2025), remplissez la fiche type ci-dessous.

\textbf{Fiche d'analyse projet de coopération}

\begin{center}
\begin{tabularx}{\linewidth}{|l|X|}\hline
\rowcolor{popblueL}
\textbf{Rubrique} & \textbf{Contenu} \\ \hline
\textbf{1. Titre exact du projet} & Appui à la gouvernance locale et au développement résilient dans l'Extrême-Nord (AGO-DR). \\ \hline
\textbf{2. Type de coopération} & \textbf{Bilatérale déléguée} (UE → GIZ) + \textbf{Décentralisée} (Région Extrême-Nord, Communes bénéficiaires). \\ \hline
\textbf{3. Partenaires (Donneur / Exécutant / Bénéficiaires)} & Donneur : Union Européenne (Fonds Fiduciaire d'Urgence pour l'Afrique). Exécutant : GIZ (Coopération technique allemande). Bénéficiaires directs : 15 communes, Région Extrême-Nord, populations cibles (jeunes, femmes, PDI). \\ \hline
\textbf{4. Montant total \& Durée} & 5 millions EUR (≈ 3,3 Mds FCFA) sur 4 ans (2021-2025). \\ \hline
\textbf{5. Secteur(s) dominant(s) (Code OCDE)} & 15110 (Secteur public / Décentralisation) ; 16010 (Prévention conflits / Cohésion sociale). \\ \hline
\textbf{6. Objectifs globaux \& Spécifiques (OS)} & \textbf{Global :} Renforcer la résilience des communautés face aux crises (Boko Haram, climat). \textbf{OS1 :} Capacités de planification budgétaire participative des communes. \textbf{OS2 :} Infrastructures socio-économiques (marchés, forages, centres jeunesse). \textbf{OS3 :} Mécanismes locaux de prévention conflits / cohésion sociale. \\ \hline
\textbf{7. Conditionnalités / Critères éligibilité} & Communes signataires Charte gouvernance locale ; Plans de développement communaux (PDC) à jour ; Contrepartie foncière/locale ; Transparence passation marchés (procédures GIZ/UE). \\ \hline
\textbf{8. Indicateurs de résultats (alignés ODD)} & \textbf{ODD 16.6 :} \% communes avec budget voté/publié dans délais (cible 80\%). \textbf{ODD 16.7 :} Taux participation femmes/jeunes aux assemblées communales (cible 30\%). \textbf{ODD 1.4 :} Nb bénéficiaires accès services de base (eau, marché) via infrastructures (cible 50 000). \textbf{ODD 5.5 :} \% femmes dans bureaux conseils municipaux. \\ \hline
\textbf{9. Modalités de suivi-évaluation} & Comité de pilotage semestriel (Région, MINAT, UE, GIZ) ; Rapports trimestriels GIZ ; Revue à mi-parcours (an 2) ; Évaluation finale externe (an 4) ; Audit financier annuel. \\ \hline
\textbf{10. Risques majeurs \& Mesures d'atténuation} & \textbf{Risque 1 :} Insécurité (accès terrain) → \textit{Atténuation :} Approche « do no harm », staff local, flexibilité planning. \textbf{Risque 2 :} Faible appropriation élites locales → \textit{Atténuation :} Contractualisation (PDC), formation continue, reddition de comptes publique. \textbf{Risque 3 :} Durabilité post-projet → \textit{Atténuation :} Renforcement recettes propres communes, maintenance infrastructures intégrée budget communal. \\ \hline
\end{tabularx}
\end{center}
\end{exoresolu}

\courssub{Synthèse : une architecture coopérative multi-niveaux pour la SND30}

La coopération internationale du Cameroun n'est pas un bloc monolithique. C'est un \textbf{système emboîté} :
\begin{enumerate}
    \item \textbf{Niveau macro (État) :} Cadre stratégique (SND30), négociation accords-cadres, coordination via MINREX/MINEPAT, alignement ODD, gestion dette.
    \item \textbf{Niveau méso (Sectoriel) :} Programmes sectoriels (éducation, santé, agriculture, énergie) financés par partenaires bilatéraux (France, Chine, Allemagne, Japon, Corée, Turquie, Maroc...) et multilatéraux (BM, FMI, BAD, ONU, UE). \textbf{Tableau de bord sectoriel} (Joint Sector Reviews).
    \item \textbf{Niveau micro / territorial (Décentralisé) :} Jumelages villes-villes, régions-régions, projets concrets (marché, forage, lycée, déchets), appropriation locale, innovation sociale.
    \item \textbf{Niveau levier (Privé/PPP) :} Infrastructures lourdes (Nachtigal, autoroutes, port Kribi, numérique), apport de capitaux privés, transfert risque, exigence de rentabilité et cadres contractuels solides (loi PPP 2018, décret 2020).
\end{enumerate}

\begin{aretenir}[Points clés pour l'examen]
\begin{itemize}
    \item Distinguer \textbf{coopération} (structurelle, contractuelle, long terme) et \textbf{aide humanitaire} (urgence, vitale, court terme).
    \item Maîtriser les \textbf{5 secteurs} (technique, culturel, économique, militaire, gouvernance) et savoir classifier un projet.
    \item Opposer \textbf{bilatérale} (État-État, souple, politique, ex. France/Chine) et \textbf{multilatérale} (OI, conditionnalités gouvernance, normes, ex. BM/FMI/BAD/ONU).
    \item Connaître la \textbf{coopération décentralisée} (loi 2004, jumelages communes/régions, pair-à-pair, ex. Douala-Lyon) et son apport en gouvernance locale.
    \item Comprendre le \textbf{PPP} (contrat long terme, risque partagé, levier financier, non-substitut APD) et citer Nachtigal.
    \item Savoir lire un \textbf{tableau de décaissements APD par secteur/année} et en tirer des tendances (graphique à barres empilées).
    \item Lier toute action de coopération aux \textbf{ODD} et aux \textbf{indicateurs de résultats} (sorties, effets, impact).
\end{itemize}
\end{aretenir}

\begin{exercice}[Entraînement Probatoire – Type Bac]
\begin{enumerate}
    \item Définir la coopération internationale et la distinguer de l'aide humanitaire d'urgence (4 pts).
    \item Présenter sous forme de tableau les différences fondamentales entre coopération bilatérale et multilatérale en donnant un exemple d'institution ou de pays partenaire pour chaque forme (6 pts).
    \item Le barrage de Nachtigal est souvent cité comme modèle de PPP. Expliquer pourquoi ce projet n'est pas de l'APD et identifier les trois types de risques supportés par le partenaire privé (5 pts).
    \item À partir du graphique « Évolution des décaissements APD 2018-2023 » (figure ci-dessus), décrire l'évolution de l'aide humanitaire et proposer deux raisons contextuelles au pic observé en 2020 (5 pts).
    \item Citer trois principes qui régissent la coopération décentralisée au Cameroun et donner un exemple de réalisation concrète issue d'un jumelage (5 pts).
\end{enumerate}
\end{exercice}

\begin{corrige}[Exercice n°2 – Éléments de réponse]
\begin{enumerate}
    \item \textbf{Coopération internationale :} Actions concertées entre acteurs (États, OI, CTD, privé) pour objectifs communs développement, basées sur contrat, durée, réciprocité. \textbf{Aide humanitaire :} Assistance d'urgence (vivres, abri, soins) pour sauver vies/atténuer souffrances suite catastrophe/conflit, gratuite, impartiale, courte durée, non contractuelle de développement.
    \item \textbf{Tableau comparatif :}
    \begin{center}
    \begin{tabularx}{\linewidth}{|l|X|X|}\hline
    \rowcolor{popblueL}
    \textbf{Critère} & \textbf{Bilatérale} & \textbf{Multilatérale} \\ \hline
    Acteurs & 2 États (Donneur/Bénéficiaire) & État + Organisation Internationale (BM, FMI, BAD, ONU, UE) \\ \hline
    Négociation & Directe, souple, adaptée & Cadre institutionnel standardisé, procédures lourdes \\ \hline
    Conditionnalités & Politiques, stratégiques, souvent implicites & Techniques, gouvernance, macroéconomiques, explicites (matrice) \\ \hline
    Visibilité & Forte (diplomatie, image) & Dilguée (logo OI, fonds communs) \\ \hline
    Exemple & Cameroun–France (C2D), Cameroun–Chine (FOCAC) & Programme FEC FMI, Projet PISA Banque Mondiale, PNUD gouvernance \\ \hline
    \end{tabularx}
    \end{center}
    \item \textbf{Pas APD :} Investissement privé (EDF, IFC) + dette commerciale (75\%) remboursée par ventes électricité (PPA), pas don/prêt concessionnel. \textbf{Risques privés :} (1) Risque construction (retard, surcoût) ; (2) Risque hydrologique / production (débit fleuve) ; (3) Risque demande / contrepartie (solvabilité ENEO, paiement factures).
    \item \textbf{Évolution humanitaire :} 120 (2018) → 140 (2019) → \textbf{200 (2020)} → 180 (2021) → 180 (2022) → 160 (2023). Pic 2020 (+43\% vs 2019). \textbf{Raisons :} (1) Pandémie COVID-19 (besoins sanitaires, filets sociaux, appui budget) ; (2) Crises sécuritaires simultanées (NOSO : déplacés internes ; Extrême-Nord : Boko Haram, inondations) → afflux PDI/réfugiés.
    \item \textbf{Principes :} (1) Libre administration des CTD (Loi 2004/017 art. 30) ; (2) Réciprocité / bénéfice mutuel ; (3) Subsidiarité / proximité (projets locaux concrets) ; (4) Alignement sur plans locaux (PDC, PDR). \textbf{Exemple :} Jumelage Douala–Lyon → Création ALEC Douala (Agence Locale Énergie Climat), plan climat ville, formation agents, éclairage public solaire.
\end{enumerate}
\end{corrige}

\coursec{Dossier 1 : L'assistance humanitaire}

Après avoir analysé les \emph{aspects et les formes de la coopération internationale} (section 2), nous nous penchons sur l'une de ses manifestations les plus visibles et les plus urgentes : l'assistance humanitaire. Cette forme de coopération se distingue par son caractère d'urgence, son ancrage dans le droit international humanitaire et son objectif premier : sauver des vies et préserver la dignité humaine face aux crises.

\courssub{Définition et principes humanitaires}

\begin{definition}[Assistance humanitaire]
L'\cle{assistance humanitaire} est une aide d'urgence, impartiale et neutre, destinée à protéger la vie, la santé et la dignité des populations affectées par des conflits armés, des catastrophes naturelles ou des épidémies. Elle s'appuie sur le droit international humanitaire (Conventions de Genève de 1949 et leurs Protocoles additionnels) et sur les principes fondamentaux énoncés par la résolution 46/182 de l'Assemblée générale des Nations unies (1991).
\end{definition}

Les quatre principes directeurs, souvent résumés par l'acronyme \textbf{HNII}, structurent toute action humanitaire :

\begin{itemize}
    \item \textbf{Humanité} : prévenir et soulager les souffrances humaines, où qu'elles se trouvent, sans discrimination.
    \item \textbf{Neutralité} : ne pas prendre parti dans les hostilités ni dans les controverses d'ordre politique, racial, religieux ou idéologique.
    \item \textbf{Impartialité} : agir uniquement en fonction des besoins, en donnant la priorité aux cas les plus urgents.
    \item \textbf{Indépendance} : l'action humanitaire doit être autonome par rapport aux objectifs politiques, économiques, militaires ou autres des acteurs étatiques ou non étatiques.
\end{itemize}

\begin{propriete}[Respect de la neutralité]
Le respect de la neutralité impose à tout acteur humanitaire de ne pas prendre parti dans le conflit. Cela signifie qu'une ONG ne doit pas, par exemple, distribuer de l'aide exclusivement dans les zones contrôlées par l'une des parties belligérantes, sous peine de perdre sa légitimité et son accès aux populations vulnérables.
\end{propriete}

\begin{attention}[Piège fréquent : confusion aide humanitaire / aide au développement]
L'aide humanitaire répond à des besoins \emph{immédiats} (nourriture, abri, soins d'urgence) sur une courte période. L'aide au développement vise des changements \emph{structurels} et \emph{durables} (éducation, infrastructures, gouvernance) sur le long terme. Les deux sont complémentaires mais ne se substituent pas ; les mélanger brouille les responsabilités et peut compromettre l'accès humanitaire.
\end{attention}

\courssub{Principaux acteurs de l'assistance humanitaire}

L'architecture humanitaire repose sur une pluralité d'acteurs aux mandats complémentaires :

\begin{center}
\begin{tabularx}{\linewidth}{|l|X|X|}
\hline
\rowcolor{popblueL}
\textbf{Acteur} & \textbf{Mandat principal} & \textbf{Exemple d'intervention au Cameroun} \\ \hline
ONU – \cle{OCHA} (Bureau de la coordination des affaires humanitaires) & Coordination globale, plaidoyer, mobilisation des ressources & Coordination du Plan de réponse humanitaire (HRP) 2023‑2024 \\ \hline
ONU – \cle{HCR} (Haut Commissariat aux réfugiés) & Protection et assistance aux réfugiés, rapatriement volontaire & Gestion des camps de l'Est (Gado, Borgop, Mandjou) \\ \hline
ONU – \cle{PAM} (Programme alimentaire mondial) & Assistance alimentaire, nutrition, logistique & Distribution de vivres aux déplacés du Nord‑Ouest/Sud‑Ouest \\ \hline
\cle{CICR} / \cle{Croix‑Rouge} (Mouvement international) & Protection des victimes de conflits, visite des détenus, rétablissement des liens familiaux & Visites des centres de détention, appui aux familles séparées \\ \hline
\cle{ONG} internationales (MSF, ACF, NRC, IRC, etc.) & Soins médicaux, nutrition, eau‑hygiène‑assainissement (EHA), protection & MSF : cliniques mobiles dans le Nord‑Ouest ; ACF : programmes nutritionnels à l'Est \\ \hline
État hôte (Gouvernement du Cameroun) & Responsabilité primaire de protection, facilitation de l'accès, coordination nationale & Ministère de l'Action sociale (MINAS), Comité national de secours (CNS) \\ \hline
\end{tabularx}
\end{center}

\begin{savaistu}[Le saviez‑vous ?]
Le Cameroun accueille la plus grande population de réfugiés centrafricains d'Afrique centrale (environ \textbf{350 000 personnes} en 2023, source HCR). Le site de \textbf{Gado-Badzéré} (région de l'Est) est le plus grand site d'hébergement, avec environ \textbf{30 000 réfugiés}. D'autres sites majeurs incluent Borgop, Mandjou et Timangolo.
\end{savaistu}

\courssub{Cycle de l'aide humanitaire : de l'alerte au suivi‑évaluation}

L'efficacité de l'assistance repose sur un cycle structuré en cinq phases. Chaque phase conditionne la suivante ; une défaillance à une étape se répercute sur l'ensemble de la réponse.

\begin{popfigure}
\begin{tikzpicture}[
    node distance=1.8cm and 2.2cm,
    box/.style={rectangle, rounded corners, draw=popblue, thick, fill=popblueL, minimum width=2.8cm, minimum height=1.2cm, align=center, font=\small},
    arrow/.style={-Stealth, thick, popdark}
]
\node[box, align=center] (alerte){1. Alerte\\\& détection};
\node[box, right=of alerte, align=center] (eval){2. Évaluation\\\& analyse des besoins};
\node[box, right=of eval, align=center] (mob){3. Mobilisation\\\& planification};
\node[box, right=of mob, align=center] (dist){4. Distribution\\\& mise en œuvre};
\node[box, right=of dist, align=center] (suivi){5. Suivi‑évaluation\\\& redevabilité};

\draw[arrow] (alerte) -- (eval);
\draw[arrow] (eval) -- (mob);
\draw[arrow] (mob) -- (dist);
\draw[arrow] (dist) -- (suivi);
\draw[arrow, bend left=30] (suivi) to node[above, font=\tiny, text=popdark] {Apprentissage \& ajustement} (alerte);
\end{tikzpicture}
\legende{Cycle de l'aide humanitaire : cinq phases interdépendantes, de l'alerte au suivi‑évaluation. La flèche de retour illustre la boucle d'apprentissage continue.}
\end{popfigure}

\begin{methode}[Lecture d'un rapport humanitaire – fiche structurée]
Pour analyser un document (rapport OCHA, bulletin PAM, évaluation MSF), utilisez la fiche suivante :
\begin{enumerate}
    \item \textbf{Contexte} : nature de la crise, zone, population concernée, chronologie.
    \item \textbf{Acteurs} : qui intervient ? (ONU, ONG, État, communautés). Quels sont leurs mandats ?
    \item \textbf{Moyens} : ressources financières, logistiques, humaines mobilisées ; pipelines d'approvisionnement.
    \item \textbf{Résultats} : indicateurs atteints (tonnes distribuées, personnes soignées, taux de couverture).
    \item \textbf{Limites \& leçons} : accès, sécurité, financement, coordination, gaps identifiés.
\end{enumerate}
\end{methode}

\courssub{Cas camerounais : deux crises majeures}

\courssub{Crise du Nord‑Ouest / Sud‑Ouest (déplacés internes)}
Depuis 2017, la contestation sociopolitique dans les régions anglophones a engendré un conflit armé entre groupes séparatistes et forces de défense. Conséquences humanitaires (chiffres OCHA/HCR 2023) :
\begin{itemize}
    \item Plus de \textbf{700 000 déplacés internes} (IDP) à l'intérieur des deux régions et vers le Littoral, l'Ouest et le Centre.
    \item Destruction d'écoles, de centres de santé, de marchés ; perturbation des moyens d'existence (agriculture, commerce).
    \item Besoins prioritaires : protection (violences basées sur le genre, recrutement d'enfants), sécurité alimentaire, éducation, abris, EHA.
\end{itemize}

\courssub{Afflux de réfugiés centrafricains à l'Est}
La crise centrafricaine (2013‑présent) a poussé plus de \textbf{350 000 réfugiés} au Cameroun, dont la majorité dans la région de l'Est. Le HCR et le gouvernement gèrent une dizaine de sites (Gado, Borgop, Mandjou, Timangolo, etc.). Les défis : intégration locale, pression sur les ressources naturelles (eau, bois), coexistence avec communautés d'accueil, financement décroissant.

\begin{exemplebox}[Exemple chiffré – PAM 2023]
En 2023, le \cle{PAM} a distribué \textbf{45 000 tonnes} de vivres (céréales, légumineuses, huile, sel, aliments enrichis) à \textbf{1,2 million de personnes} au Cameroun (déplacés internes, réfugiés, populations hôtes vulnérables). Le coût total de l'opération s'élevait à environ \textbf{180 millions USD}, financé à 68~\% par les contributions volontaires des États donateurs.
\end{exemplebox}

\courssub{Données chiffrées 2020‑2024 : évolution des déplacés internes au Cameroun}

Le graphique ci‑dessous illustre l'évolution du nombre de personnes déplacées internes (IDP) enregistrées par OCHA entre 2017 et 2024. On observe une montée rapide jusqu'en 2020, un plateau en 2021‑2022, puis une légère baisse en 2023‑2024 grâce aux retours volontaires et aux programmes de réintégration, bien que les chiffres restent élevés.

\begin{popfigure}
\begin{tikzpicture}
\begin{axis}[
    ybar,
    bar width=0.55cm,
    width=\linewidth,
    height=9cm,
    enlarge x limits=0.15,
    ylabel={Nombre de déplacés internes (milliers)},
    xlabel={Année},
    symbolic x coords={2017,2018,2019,2020,2021,2022,2023,2024},
    xtick=data,
    nodes near coords,
    nodes near coords align={vertical},
    every node near coord/.append style={font=\tiny, rotate=90, anchor=west, yshift=2pt},
    ymin=0,
    ymax=800,
    ymajorgrids=true,
    grid style=dashed,
    tick label style={font=\small},
    label style={font=\small},
    legend style={at={(0.5,-0.15)}, anchor=north, legend columns=-1, font=\small},
    title={Évolution des déplacés internes au Cameroun (2017‑2024)},
    title style={font=\normalsize\bfseries, yshift=0.5cm}
]
\addplot[fill=popblue!70, draw=popblue] coordinates {
    (2017, 30)
    (2018, 120)
    (2019, 350)
    (2020, 680)
    (2021, 710)
    (2022, 720)
    (2023, 690)
    (2024, 660)
};
\addlegendentry{IDP (×1 000)}
\end{axis}
\end{tikzpicture}
\legende{Source : OCHA Cameroun, Rapports annuels 2017‑2024. Les valeurs sont arrondies au millier le plus proche.}
\end{popfigure}

\begin{aretenir}[Points clés à retenir]
\begin{itemize}
    \item L'assistance humanitaire repose sur quatre principes indissociables : humanité, neutralité, impartialité, indépendance.
    \item Elle mobilise une architecture multi‑acteurs (ONU, Croix‑Rouge, ONG, État) dont la coordination est assurée par OCHA.
    \item Le cycle de l'aide (alerte → évaluation → mobilisation → distribution → suivi‑évaluation) garantit la pertinence et la redevabilité.
    \item Au Cameroun, deux crises majeures (Nord‑Ouest/Sud‑Ouest et Est) génèrent des besoins massifs : ~700 000 IDP et ~350 000 réfugiés centrafricains (2023‑2024).
    \item Les données chiffrées (PAM 2023 : 45 000 t de vivres pour 1,2 M de bénéficiaires) illustrent l'ampleur logistique et financière de la réponse.
\end{itemize}
\end{aretenir}

\courssub{Mise en situation : rédaction d'une note diplomatique humanitaire}

\begin{exoresolu}[Exercice résolu : Note diplomatique]
\textbf{Énoncé :} Vous êtes attaché(e) au Ministère des Relations extérieures (MINREX). Rédigez une note diplomatique adressée au Représentant résident du PAM au Cameroun pour solliciter une augmentation de 15~\% de l'allocation alimentaire destinée aux déplacés du Nord‑Ouest pour le trimestre Q3 2024. La note doit respecter la forme diplomatique (en‑tête, référence, objet, corps, formule de politesse, signature) et s'appuyer sur des données chiffrées.

\textbf{Solution détaillée :}

\begin{verbatim}
RÉPUBLIQUE DU CAMEROUN
Paix – Travail – Patrie
-------------------
MINISTÈRE DES RELATIONS EXTÉRIEURES
Direction de la Coopération Humanitaire
N° 042/2024/MINREX/DCH

Yaoundé, le 15 juillet 2024

Objet : Demande d'augmentation de l'allocation alimentaire PAM
        pour les déplacés internes du Nord‑Ouest (Q3 2024)

Monsieur le Représentant résident,

Par la présente, j'ai l'honneur de porter à votre attention la situation
alimentaire préoccupante des personnes déplacées internes (PDI) de la
région du Nord‑Ouest, estimées à 320 000 personnes selon le dernier
rapport OCHA (mai 2024).

Les distributions actuelles (ration de 2 100 kcal/jour/personne) couvrent
environ 85 % des besoins nutritionnels minimums. L'évaluation conjointe
MINAS‑PAM‑ONG (juin 2024) révèle une détérioration de l'indice de
sécurité alimentaire (IPC phase 3 « Crise » pour 45 % des ménages
enquêtés).

Au vu de ces éléments, le Gouvernement de la République du Cameroun
sollicite une augmentation de 15 % du tonnage mensuel alloué (soit
environ 1 200 tonnes supplémentaires par mois) pour le troisième
trimestre 2024, afin de porter la ration à 2 400 kcal/jour/personne et
de réduire le gap nutritionnel.

Les financements complémentaires nécessaires (estimés à 4,8 millions USD)
seront inscrits dans le Plan de réponse humanitaire 2024 révisé, en
coordination avec OCHA et les donateurs bilatéraux.

Je vous saurais gré de bien vouloir transmettre cette demande aux
instances compétentes du PAM et de nous informer de la suite réservée.

En vous renouvelant l'assurance de ma haute considération.

Le Ministre des Relations extérieures
(Signature)
Le Directeur de la Coopération Humanitaire
\end{verbatim}
\end{exoresolu}

\begin{exercice}[Exercice d'application]
À partir du graphique de l'évolution des IDP (2017‑2024) et des données du PAM 2023, répondez aux questions suivantes :
\begin{enumerate}
    \item Calculez le taux de variation du nombre d'IDP entre 2020 et 2024.
    \item Identifiez deux facteurs possibles expliquant la baisse observée après 2022.
    \item Proposez deux indicateurs de suivi‑évaluation pertinents pour mesurer l'efficacité de la distribution alimentaire du PAM en 2024.
\end{enumerate}
\end{exercice}

\begin{corrige}[Exercice d'application]
\begin{enumerate}
    \item Taux de variation = $\frac{660-680}{680}\times100 = -2,94~\%$ (baisse d'environ 3~\%).
    \item Facteurs : (a) retours volontaires facilités par l'amélioration sécuritaire locale ; (b) programmes de réintégration socio‑économique (appui à l'agriculture, formation professionnelle) menés par le MINAS et les ONG.
    \item Indicateurs : (a) \emph{Taux de couverture alimentaire} = (bénéficiaires effectivement servis / bénéficiaires ciblés) $\times100$ ; (b) \emph{Score de diversité alimentaire des ménages (HDDS)} mesuré par enquête post‑distribution.
\end{enumerate}
\end{corrige}

\coursec{Les principes de la diplomatie camerounaise}

Après avoir analysé les mécanismes de l'assistance humanitaire — cette expression concrète de la solidarité internationale — nous comprenons que l'action extérieure d'un État ne se limite pas à l'urgence. Elle s'inscrit dans une stratégie de long terme, guidée par des principes directeurs. Le Cameroun, par sa position au carrefour de l'Afrique centrale et de l'Afrique de l'Ouest, a forgé une doctrine diplomatique originale, mêlant héritage historique, réalisme géopolitique et ambition de développement. Étudions les piliers qui structurent cette action.

\courssub{Les fondements de l'action extérieure du Cameroun}

La diplomatie camerounaise ne s'improvise pas ; elle repose sur des constantes qui guident le Ministère des Relations Extérieures (MINREX) et l'ensemble du corps diplomatique. Ces principes, conformes à la Charte des Nations Unies et à l'Acte constitutif de l'Union Africaine, constituent la \cle{grille de lecture} de tout engagement international du pays.

\begin{definition}[Diplomatie]
Ensemble des méthodes, procédures et actions par lesquelles un État conduit ses relations extérieures avec d'autres États ou organisations internationales, dans le but de défendre ses intérêts, de promouvoir la paix et de favoriser la coopération.
\end{definition}

\begin{definition}[Relations internationales (rappel)]
Ensemble des relations (politiques, économiques, culturelles, militaires) que les États et les organisations internationales entretiennent entre eux. Elles peuvent être \cle{bilatérales} (entre deux États) ou \cle{multilatérales} (au sein d'organisations regroupant plusieurs États : ONU, UA, CEMAC).
\end{definition}

\begin{propriete}[Primauté du multilatéralisme africain]
La diplomatie camerounaise privilégie le cadre multilatéral, notamment au sein de l'Union Africaine (UA), de la CEMAC (Communauté Économique et Monétaire de l'Afrique Centrale) et de la CEEAC (Communauté Économique des États de l'Afrique Centrale). Le Cameroun considère que l'intégration régionale est une condition essentielle de sa sécurité et de son développement.
\end{propriete}

\textbf{Principe 1 : Respect de la souveraineté et de l'intégrité territoriale.}\par
C'est la pierre angulaire de la doctrine camerounaise. Ayant lui-même connu un conflit frontalier (affaire de la péninsule de Bakassi, réglée par l'arrêt de la Cour Internationale de Justice de 2002 et par l'accord de Greentree de 2006), le Cameroun élève le respect des frontières héritées de la colonisation (principe de l'\cle{uti possidetis juris}) en règle intangible. Il refuse toute ingérence militaire extérieure non sollicitée et toute remise en cause des frontières des États voisins.

\begin{exemplebox}[Application : la péninsule de Bakassi]
Face à l'arrêt de la Cour Internationale de Justice (CIJ) de 2002 attribuant la péninsule de Bakassi au Cameroun, Yaoundé a privilégié la voie juridique et diplomatique (accord de Greentree, 2006, sous l'égide de l'ONU) plutôt que la force. Le retrait nigérian s'est achevé en 2008 (transfert définitif de l'administration en 2013). Cette gestion pacifique a valu au pays une reconnaissance internationale de sa maturité étatique.
\end{exemplebox}

\textbf{Principe 2 : Non-ingérence dans les affaires intérieures des États.}\par
Ce principe, corollaire de la souveraineté, interdit à un État d'intervenir dans le domaine réservé d'un autre (choix politiques, système électoral, gestion des ressources). Attention cependant : \cle{non-ingérence ne signifie pas indifférence}.

\begin{attention}[Piège fréquent : confondre non-ingérence et indifférence]
Le Cameroun n'hésite pas à proposer ses \cle{bons offices} ou à accepter un mandat de médiation lorsque la stabilité d'un voisin menace la sécurité régionale. L'ingérence est un acte unilatéral imposé à un État ; la médiation est une démarche concertée, acceptée par les parties en conflit. Dans une copie, ne jamais écrire que la non-ingérence interdit toute action du Cameroun chez ses voisins.
\end{attention}

\textbf{Principe 3 : Recherche du règlement pacifique des différends.}\par
L'article 33 de la Charte de l'ONU énumère les moyens pacifiques de règlement des conflits. Le Cameroun les utilise de manière graduée :
\begin{itemize}
    \item \textbf{Négociation directe :} pour les contentieux frontaliers mineurs ou la gestion des ressources partagées (fleuves transfrontaliers, barrages).
    \item \textbf{Médiation et bons offices :} rôle de facilitateur neutre accepté par les parties (crises en Afrique centrale).
    \item \textbf{Arbitrage et justice internationale :} recours à la CIJ pour les contentieux juridiques complexes (affaire de Bakassi avec le Nigeria).
\end{itemize}

\textbf{Principe 4 : Promotion de la coopération régionale et continentale.}\par
Le Cameroun est membre actif des organisations sous-régionales et continentales. Cette présence institutionnelle ancre le pays comme un acteur diplomatique central de l'Afrique centrale.

\begin{aretenir}[Organisations auxquelles le Cameroun participe activement]
\begin{itemize}
    \item \textbf{CEMAC} (siège à Bangui) : marché commun, monnaie unique (le franc CFA de la zone BEAC), libre circulation des personnes et des biens.
    \item \textbf{CEEAC} (siège à Libreville) : sécurité collective, intégration physique (routes, énergie).
    \item \textbf{UA} (siège à Addis-Abeba) : participation aux efforts de paix du continent.
    \item \textbf{ONU} : le Cameroun est membre fondateur depuis son adhésion en 1960.
    \item \textbf{Commonwealth} : le Cameroun y a adhéré en 1995, ce qui reflète son bilinguisme et son double héritage anglophone et francophone.
\end{itemize}
\end{aretenir}

\textbf{Principe 5 : Diplomatie économique au service du développement.}\par
Depuis la \cle{Stratégie Nationale de Développement 2020-2030 (SND30)}, l'action diplomatique est explicitement orientée vers l'attraction des Investissements Directs Étrangers (IDE), l'ouverture des marchés aux produits camerounais (cacao, café, bois, coton, banane, aluminium) et le transfert de technologies. Les ambassadeurs reçoivent désormais des lettres de mission comportant des objectifs économiques.

\begin{exemplebox}[Illustration : la diplomatie économique en action]
Les grands projets structurants du Cameroun illustrent cette diplomatie économique : le barrage hydroélectrique de Nachtigal (financement international associant notamment la France et la Banque mondiale), le port en eau profonde de Kribi (construit avec des partenaires chinois et français), l'autoroute Yaoundé--Douala. Chaque projet résulte de négociations diplomatiques et financières menées par l'État avec des partenaires variés : le Cameroun diversifie ses partenaires (Union européenne, Chine, Turquie, pays du Golfe) pour ne dépendre d'aucun.
\end{exemplebox}

\courssub{Illustrations par des actes diplomatiques}

La théorie trouve sa preuve dans l'action. Deux dossiers illustrent la mise en œuvre combinée de ces principes.

\begin{experience}[Étude de cas guidée : l'engagement du Cameroun pour la paix en Centrafrique]
\textbf{Contexte :} la République centrafricaine (RCA), voisine du Cameroun à l'Est, traverse de graves crises sécuritaires. L'insécurité en RCA affecte directement le Cameroun : afflux de réfugiés dans les régions de l'Est et de l'Adamaoua, interruption du corridor commercial Douala--Bangui (voie d'approvisionnement essentielle de la RCA, pays enclavé).
\textbf{Démarche camerounaise (analyse pas à pas) :}
\begin{enumerate}
    \item \textbf{Légitimité multilatérale (principe 4) :} le Cameroun agit dans le cadre des organisations sous-régionales (CEEAC) et aux côtés de la MINUSCA (Mission des Nations Unies en RCA), et non de manière unilatérale.
    \item \textbf{Non-ingérence respectée (principe 2) :} l'action camerounaise (bons offices, facilitation du dialogue) se fait avec l'accord des autorités centrafricaines.
    \item \textbf{Appui logistique et humanitaire (principe 4 et dossier 1) :} le corridor Douala--Bangui sert de base arrière à l'aide humanitaire et à la MINUSCA ; le Cameroun accueille des centaines de milliers de réfugiés centrafricains.
    \item \textbf{Résultat recherché :} la stabilisation de la RCA sécurise l'Est camerounais, permet le retour progressif des réfugiés et la reprise du commerce transfrontalier.
\end{enumerate}
\textbf{Interprétation :} ce dossier valide le modèle camerounais : \emph{diplomatie de proximité + légitimité institutionnelle + capacité logistique}. La sécurité du voisin est comprise comme une partie de la sécurité nationale.
\end{experience}

\begin{exoresolu}[Analyse d'un accord de coopération : le partenariat Cameroun--Union européenne]
\textbf{Énoncé :} le Cameroun entretient avec l'Union européenne des accords de coopération commerciale (Accord de Partenariat Économique, en vigueur de manière provisoire depuis 2014) et sectorielle (pêche, forêts avec l'accord APV-FLEGT sur la légalité du bois). Analysez ce partenariat au prisme des cinq principes de la diplomatie camerounaise.
\tcblower
\textbf{Solution détaillée :}
\begin{enumerate}
    \item \textbf{Souveraineté (principe 1) :} le Cameroun négocie en État souverain l'accès à ses ressources (bois, poisson) et à son marché ; il peut accepter, amender ou refuser des clauses.
    \item \textbf{Non-ingérence (principe 2) :} l'UE ne dicte pas la politique forestière ou halieutique camerounaise ; les appuis techniques (lutte contre l'exploitation illégale du bois, pêche illicite non déclarée et non réglementée) répondent à des demandes camerounaises.
    \item \textbf{Règlement pacifique (principe 3) :} les différends commerciaux éventuels se règlent par négociation dans les instances prévues par les accords, sans rapport de force.
    \item \textbf{Coopération régionale (principe 4) :} le Cameroun défend ses positions en cohérence avec la CEMAC, dont il est l'économie la plus importante.
    \item \textbf{Diplomatie économique (principe 5) :} retombées concrètes : accès préférentiel des produits camerounais (banane, cacao, bois certifié) au marché européen, financements de projets (routes, énergie), appuis budgétaires.
\end{enumerate}
\textbf{Conclusion :} ce partenariat illustre une diplomatie économique où le droit souverain sur les ressources est monétisé et réinvesti dans le développement national, tout en diversifiant les partenaires du pays.
\end{exoresolu}

\courssub{Outils et mise en œuvre : le réseau diplomatique et la note diplomatique}

Les principes ne valent que par les hommes et les procédures qui les portent. Le Cameroun dispose d'un réseau d'ambassades et de consulats sur tous les continents, ainsi que de missions permanentes auprès des grandes organisations internationales (ONU à New York et Genève, UA à Addis-Abeba, UNESCO à Paris, UE à Bruxelles). Ce maillage est l'outil opérationnel de la doctrine.

\begin{popfigure}
\begin{tikzpicture}[
    scale=0.9,
    every node/.style={font=\small},
    event/.style={circle, fill=popblue, inner sep=1.5pt},
    milestone/.style={rectangle, rounded corners, fill=poporangeL, draw=poporange, text=popdark, minimum height=1cm, text width=3.2cm, align=center, inner sep=2pt},
    arrow/.style={->, >=Stealth, thick, popblue}
]
\draw[popline, thick] (0,0) -- (14.5,0);
\foreach \x/\year in {0/1960, 3.5/1980, 7/2000, 10.5/2010, 14/2020} {
    \draw (\x,-0.15) -- (\x,0.15) node[below=2pt] {\year};
}
\node[event] (e1) at (0,0) {};
\node[milestone, anchor=south] (m1) at (0,0.5) {Indépendance et adhésion à l'ONU (1960)};
\draw[arrow] (e1) -- (m1.south);
\node[event] (e2) at (0.6,0) {};
\node[milestone, anchor=north] (m2) at (0.6,-0.5) {Réunification (1961)};
\draw[arrow] (e2) -- (m2.north);
\node[event] (e3) at (2.1,0) {};
\node[milestone, anchor=south] (m3) at (2.1,0.5) {République unie du Cameroun (1972)};
\draw[arrow] (e3) -- (m3.south);
\node[event] (e4) at (5.6,0) {};
\node[milestone, anchor=north] (m4) at (5.6,-0.5) {Adhésion au Commonwealth (1995)};
\draw[arrow] (e4) -- (m4.north);
\node[event] (e5) at (7.4,0) {};
\node[milestone, anchor=south] (m5) at (7.4,0.5) {Arrêt CIJ sur Bakassi (2002) ; membre du Conseil de sécurité (2002-2003)};
\draw[arrow] (e5) -- (m5.south);
\node[event] (e6) at (8.8,0) {};
\node[milestone, anchor=north] (m6) at (8.8,-0.5) {Accord de Greentree (2006)};
\draw[arrow] (e6) -- (m6.north);
\node[event] (e7) at (12.6,0) {};
\node[milestone, anchor=south] (m7) at (12.6,0.5) {Grand Dialogue National (2019)};
\draw[arrow] (e7) -- (m7.south);
\node[event] (e8) at (13.8,0) {};
\node[milestone, anchor=north] (m8) at (13.8,-0.5) {SND30 : diplomatie économique prioritaire (2020-2030)};
\draw[arrow] (e8) -- (m8.north);
\end{tikzpicture}
\legende{Chronologie simplifiée des grandes étapes diplomatiques du Cameroun (1960-2020). Sources : MINREX, Nations Unies.}
\end{popfigure}

\begin{popfigure}
\begin{tikzpicture}[scale=0.75, every node/.style={font=\small}]
% Cadre Afrique schématique
\draw[fill=popgreenL, draw=popline, rounded corners] (0,0) rectangle (10,8);
\node[anchor=north west, font=\small\itshape, text=popmuted] at (0.2,7.9) {Afrique (croquis très simplifié)};
% Cameroun
\fill[poporange] (3.2,3.4) circle (0.18);
\node[anchor=west, font=\small\bfseries] at (3.4,3.4) {Cameroun (Yaoundé)};
% Partenaires régionaux
\foreach \x/\y/\nom in {1.8/3.6/Nigeria, 4.6/3.2/Tchad, 4.8/2.2/RCA, 3.0/1.8/Gabon-Congo, 2.2/2.6/Guinée équat.} {
    \fill[popblue] (\x,\y) circle (0.1);
    \node[anchor=west, font=\scriptsize] at (\x+0.12,\y) {\nom};
}
% Liens régionaux
\foreach \x/\y in {1.8/3.6, 4.6/3.2, 4.8/2.2, 3.0/1.8, 2.2/2.6} {
    \draw[popblue, thick, <->, >=Stealth] (3.2,3.4) -- (\x,\y);
}
% Partenaires mondiaux (hors cadre, flèches)
\node[rectangle, rounded corners, fill=poppurpleL, draw=poppurple, text width=2.6cm, align=center] (ue) at (8.6,6.8) {Union européenne\\(coopération, commerce)};
\node[rectangle, rounded corners, fill=poppurpleL, draw=poppurple, text width=2.6cm, align=center] (ch) at (8.6,4.6) {Chine\\(infrastructures, IDE)};
\node[rectangle, rounded corners, fill=poppurpleL, draw=poppurple, text width=2.6cm, align=center] (onu) at (8.6,2.4) {ONU / UA\\(paix, multilatéralisme)};
\node[rectangle, rounded corners, fill=poppurpleL, draw=poppurple, text width=2.6cm, align=center] (ss) at (8.6,0.6) {Sud-Sud\\(Brésil, Inde, Turquie)};
\foreach \c in {ue, ch, onu, ss} {
    \draw[poppurple, thick, <->, >=Stealth] (3.2,3.4) -- (\c.west);
}
% Légende
\node[anchor=south west, align=left, font=\scriptsize, fill=white, draw=popline, rounded corners, inner sep=3pt] at (0.2,0.2) {\textcolor{popblue}{$\bullet$} Voisins (coopération régionale CEMAC/CEEAC)\\ \textcolor{poppurple}{$\leftrightarrow$} Partenaires mondiaux diversifiés};
\end{tikzpicture}
\legende{Croquis schématique du réseau de partenariats du Cameroun : ancrage régional (voisins immédiats) et diversification mondiale des partenaires. Schéma indicatif, sans précision cartographique.}
\end{popfigure}

La maîtrise de l'outil diplomatique passe par la rédaction d'actes formels. La \cle{note diplomatique} (ou \emph{note verbale}) est l'instrument de base de la correspondance entre États.

\begin{methode}[Rédiger une note diplomatique (structure standard)]
\textbf{Objectif :} transmettre une position officielle, une protestation, une demande d'information ou une réponse, sans caractère personnel.
\begin{enumerate}
    \item \textbf{En-tête :} République du Cameroun -- Ministère des Relations Extérieures -- Direction concernée. Numéro de référence, lieu et date.
    \item \textbf{Destinataire :} Ambassade du pays concerné à Yaoundé, ou Ministère des Affaires étrangères du pays (par voie diplomatique).
    \item \textbf{Objet :} clair, précis, en majuscules.
    \item \textbf{Formule d'entrée :} \emph{« Le Ministère des Relations Extérieures présente ses compliments à l'Ambassade de ... et a l'honneur de lui faire savoir que... »}
    \item \textbf{Contexte et rappel des faits :} exposé neutre et factuel (résolutions, accords antérieurs, événements).
    \item \textbf{Position et arguments :} fondement juridique (traité, coutume, principe), intérêt national, cohérence avec les engagements régionaux (UA, CEMAC).
    \item \textbf{Demande ou proposition :} action attendue, avec délai si pertinent.
    \item \textbf{Formule de politesse finale :} \emph{« Le Ministère des Relations Extérieures saisit cette occasion pour renouveler à l'Ambassade de ... les assurances de sa haute considération. »}
    \item \textbf{Signature :} le Ministre ou son représentant habilité, avec cachet officiel.
\end{enumerate}
\end{methode}

\begin{exoresolu}[Rédaction guidée : note verbale de protestation (cas fictif d'entraînement)]
\textbf{Énoncé :} \emph{(situation fictive)} un aéronef étranger non autorisé a survolé une localité frontalière du Cameroun. Rédigez la note verbale adressée à l'ambassade du pays concerné.
\tcblower
\textbf{Solution :}
\begin{quote}
\textsc{République du Cameroun}\\
\textsc{Ministère des Relations Extérieures}\\
N° 0042/MINREX/SG \hfill Yaoundé, le 16 mars 2024

\textbf{À Son Excellence Monsieur l'Ambassadeur de la République X à Yaoundé}

\textbf{OBJET : PROTESTATION SUITE À LA VIOLATION DE L'ESPACE AÉRIEN CAMEROUNAIS.}

Excellence,

Le Ministère des Relations Extérieures présente ses compliments à l'Ambassade de la République X et a l'honneur de lui faire part de sa vive préoccupation suite au survol non autorisé de l'espace aérien camerounais, constaté le 15 mars 2024 au-dessus d'une localité frontalière du territoire national.

Cet acte constitue une violation de la souveraineté et de l'intégrité territoriale du Cameroun, principes consacrés par la Charte des Nations Unies et l'Acte constitutif de l'Union Africaine. Il contrevient en outre aux engagements de bon voisinage souscrits dans le cadre sous-régional.

Au regard de la gravité de cet incident, le Gouvernement de la République du Cameroun :
\begin{enumerate}
    \item condamne fermement cette intrusion ;
    \item demande des explications officielles dans les meilleurs délais ;
    \item exige l'engagement formel que de tels actes ne se renouvelleront pas ;
    \item se réserve le droit de saisir les instances régionales compétentes.
\end{enumerate}

Le Ministère des Relations Extérieures saisit cette occasion pour renouveler à l'Ambassade de la République X les assurances de sa haute considération.

\textit{(Signature et cachet)}\\
\textbf{Le Ministre des Relations Extérieures}
\end{quote}
\textbf{Analyse :} structure respectée (en-tête, objet, formules consacrées) ; ton ferme mais courtois ; ancrage juridique (Charte ONU, UA) ; demandes précises ; mention des recours multilatéraux (principe 4).
\end{exoresolu}

\courssub{Synthèse : une diplomatie d'équilibre au service de l'émergence}

\begin{center}
\begin{tabularx}{\linewidth}{|l|X|X|}\hline
\rowcolor{popblueL}
\textbf{Principe} & \textbf{Mécanisme opérationnel} & \textbf{Illustration} \\ \hline
1. Souveraineté et intégrité territoriale & Voie juridique (CIJ), sécurisation des frontières & Règlement pacifique de Bakassi (2002-2006) \\ \hline
2. Non-ingérence (active) & Bons offices, médiation acceptée, refus des ingérences & Facilitation du dialogue en Afrique centrale \\ \hline
3. Règlement pacifique des différends & Négociation, médiation, arbitrage, CIJ & Contentieux réglés sans recours à la force \\ \hline
4. Coopération régionale et continentale & CEMAC, CEEAC, UA, opérations de paix & Corridor Douala--Bangui, accueil de réfugiés \\ \hline
5. Diplomatie économique & Attraction des IDE, accords commerciaux, diversification des partenaires & Nachtigal, port de Kribi, APE avec l'UE \\ \hline
\end{tabularx}
\end{center}

\begin{savaistu}[Le Cameroun au Conseil de sécurité de l'ONU]
Le Cameroun a siégé comme membre non permanent du Conseil de sécurité des Nations Unies pour la période 2002-2003 (il l'avait déjà été en 1974-1975). Cette responsabilité, la plus élevée qu'un État puisse exercer en matière de paix mondiale en dehors des membres permanents, témoigne de la confiance de la communauté internationale et a consolidé l'expertise multilatérale du corps diplomatique camerounais.
\end{savaistu}

\begin{exercice}[Entraînement au Probatoire / Baccalauréat technique -- sujet type]
\textbf{Sujet :} « La diplomatie camerounaise est souvent décrite comme une diplomatie de l'équilibre. » À partir des cinq principes étudiés, analysez cette affirmation en vous appuyant sur deux exemples concrets.
\textbf{Consignes :}
\begin{enumerate}
    \item Définissez en introduction la « diplomatie de l'équilibre » dans le contexte camerounais.
    \item Structurez le développement en deux parties : 1) l'équilibre entre souveraineté et ouverture (principes 1, 2 et 5) ; 2) l'équilibre entre ancrage régional et rayonnement mondial (principes 3 et 4).
    \item Appuyez-vous obligatoirement sur le règlement de Bakassi et sur l'engagement du Cameroun en Centrafrique.
    \item Concluez sur la cohérence de cette doctrine face aux défis actuels (sécurité, émergence économique).
\end{enumerate}
\end{exercice}

\begin{corrige}[Exercice n°1 -- éléments de corrigé attendus]
\textbf{Introduction :} la diplomatie de l'équilibre est la capacité d'un État à préserver son autonomie de décision (souveraineté) tout en s'engageant dans des partenariats multiples (Occident, Chine, Afrique) sans alignement exclusif sur un seul bloc.
\textbf{Partie 1 -- Équilibre entre souveraineté et ouverture :}
\begin{itemize}
    \item Souveraineté affirmée (principe 1) : règlement de Bakassi par le droit, sans céder ni recourir à la guerre.
    \item Non-ingérence active (principe 2) : la médiation en Centrafrique est une intervention \emph{acceptée}, donc légitime, et non une ingérence.
    \item Diplomatie économique diversifiée (principe 5) : partenariats simultanés avec l'UE (APE), la Chine (port de Kribi), des bailleurs multilatéraux (Nachtigal) : aucun partenaire exclusif.
\end{itemize}
\textbf{Partie 2 -- Équilibre entre ancrage régional et rayonnement mondial :}
\begin{itemize}
    \item Ancrage régional (principe 4) : rôle moteur dans la CEMAC et la CEEAC, corridor Douala--Bangui, accueil des réfugiés centrafricains.
    \item Rayonnement mondial (principe 3) : membre du Conseil de sécurité (2002-2003), membre du Commonwealth (1995), recours à la CIJ.
    \item Boîte à outils complète du règlement pacifique : négociation, médiation, justice internationale.
\end{itemize}
\textbf{Conclusion :} l'équilibre n'est pas une immobilité mais une stratégie d'adaptation. Le défi des années 2020-2030 est de transformer cet équilibre diplomatique en levier concret d'industrialisation (SND30) et de sécurité (Boko Haram dans l'Extrême-Nord, crises en Afrique centrale).
\end{corrige}

\coursec{Synthèse et mise en situation : rédaction d'une note diplomatique}

\courssub{Transition et revue croisée des acquis}

Après avoir analysé en détail les \cle{principes directeurs de la diplomatie camerounaise} (non-ingérence, égalité souveraine des États, règlement pacifique des différends, coopération internationale, panafricanisme), nous disposons désormais de la « boîte à outils » conceptuelle pour agir en diplomate. Les sections précédentes ont posé les fondations : les \cle{Relations Internationales} comme cadre global d'interaction entre acteurs étatiques et non étatiques ; la \cle{coopération internationale} (bilatérale, multilatérale, décentralisée) comme modalité concrète de l'action ; et l'\cle{assistance humanitaire} comme expression de la solidarité internationale et du droit international humanitaire.

\begin{aretenir}[Revue croisée des quatre thèmes]
\begin{itemize}
    \item \textbf{Relations Internationales :} cadre juridique (Charte des Nations Unies, Acte constitutif de l'Union Africaine) et acteurs (États, organisations internationales, ONG, entreprises transnationales).
    \item \textbf{Coopération internationale :} formes (technique, financière, culturelle, militaire) et instruments (accords-cadres, conventions spécifiques, protocoles d'accord).
    \item \textbf{Assistance humanitaire :} principes (humanité, neutralité, impartialité, indépendance), acteurs (OCHA, CICR, PAM, HCR, ONG), phases de la gestion de crise.
    \item \textbf{Diplomatie camerounaise :} principes constitutionnels (préambule et article 45 de la Constitution de 1996), diplomatie de voisinage, diplomatie économique, diplomatie culturelle, rôle de médiateur (exemples : règlement du différend de Bakassi, médiation dans la crise centrafricaine).
\end{itemize}
\end{aretenir}

\courssub{Analyse d'un dossier composite : le différend frontalier de la « Zone de Mbalang »}

Pour mettre en œuvre nos acquis, nous simulons une situation de crise frontalière \textbf{fictive mais réaliste}, inspirée des contentieux réels que le Cameroun a connus (Bakassi, frontière avec le Nigeria, bassin du Lac Tchad). Vous êtes attaché(e) au \textbf{Département des Affaires Juridiques et Consulaires} du \textbf{Ministère des Relations Extérieures (MINREX)}.

\begin{attention}[Nature du dossier]
La « République de Loména » et la « Zone de Mbalang » sont des \textbf{cas fictifs d'entraînement}. En revanche, les textes juridiques cités (Arrêt de la CIJ de 2002, Accord de Greentree de 2006, Convention de Vienne de 1969, Résolution 46/182 de l'ONU) sont \textbf{réels} et doivent être connus pour l'examen.
\end{attention}

\begin{popfigure}
\centering
\begin{tikzpicture}[node distance=1.1cm and 1.4cm, font=\small]
    \tikzset{
        doc/.style={rectangle, draw=popblue, thick, fill=popblue!5, rounded corners, minimum width=4.2cm, minimum height=1.9cm, align=center},
        arrow/.style={-Stealth, thick, popdark},
        tag/.style={font=\tiny\itshape, text=popmuted, anchor=north west}
    }
    \node[doc, align=center] (doc1){\textbf{Communiqué de presse MINREX}\\[2pt]\textit{« Le Cameroun condamne}\\\textit{l'incursion... »}};
    \node[doc, right=of doc1, align=center] (doc2){\textbf{Extraits juridiques}\\[2pt]\textit{Arrêt CIJ 2002,}\\\textit{Accord de Greentree 2006,}\\\textit{Accord de démarcation 2018}};
    \node[doc, below=of doc1, align=center] (doc3){\textbf{Données humanitaires OCHA}\\[2pt]\textit{12 450 déplacés,}\\\textit{2,3 millions USD de besoins,}\\\textit{3 centres de santé fermés}};
    \node[doc, right=of doc3, align=center] (doc4){\textbf{Rapport de l'État-Major}\\[2pt]\textit{Mouvements de troupes,}\\\textit{violation de la ligne}\\\textit{de démarcation}};
    \node[doc, fill=popgold!10, draw=popgold, below=2.6cm of doc1, xshift=2.8cm, minimum width=6cm, minimum height=2.6cm, align=center] (synth){
        \textbf{NOTE DE POSITION MINREX}\\[2pt]
        \textit{Objet : différend Zone Mbalang}\\[3pt]
        \textbf{Principes mobilisés :}\\
        intégrité territoriale -- règlement pacifique\\
        droit humanitaire -- bonne foi
    };
    \draw[arrow] (doc1.south) -- (synth.north west);
    \draw[arrow] (doc2.south) -- (synth.north east);
    \draw[arrow] (doc3.east) -- (synth.west);
    \draw[arrow] (doc4.west) -- (synth.east);
    \node[tag] at (doc1.north west) {Source officielle};
    \node[tag] at (doc2.north west) {Base légale};
    \node[tag] at (doc3.north west) {Contexte humain};
    \node[tag] at (doc4.north west) {Faits de terrain};
\end{tikzpicture}
\legende{Figure 5.1 : Architecture du dossier composite pour la rédaction de la note diplomatique. Quatre pièces de nature différente (politique, juridique, humanitaire, militaire) convergent vers un document unique : la note de position.}
\end{popfigure}

\begin{exemplebox}[Contenu simulé du dossier « Mbalang »]
\textbf{1. Communiqué de presse MINREX (extrait) :} « Le Gouvernement de la République du Cameroun proteste énergiquement contre la progression, le 12 mars 2025, d'éléments des forces de défense du pays voisin, la République de Loména, au nord du pilier frontalier n°47 (Zone de Mbalang), en violation de l'Accord de démarcation de 2018. Le Cameroun réaffirme son attachement à l'intangibilité des frontières héritées de la colonisation (\emph{uti possidetis juris}). »

\textbf{2. Extraits de textes juridiques :}
\begin{itemize}
    \item \emph{Arrêt de la CIJ, Cameroun c. Nigeria (10 octobre 2002) :} consacre le principe de l'intangibilité des frontières héritées de la colonisation (\emph{uti possidetis juris}) et attribue la péninsule de Bakassi au Cameroun.
    \item \emph{Accord de Greentree (12 juin 2006) :} met en œuvre l'arrêt de la CIJ ; organise le retrait des troupes nigérianes et le transfert d'autorité à Bakassi.
    \item \emph{Accord de démarcation camerouno-loménien (2018, fictif) :} matérialisation de la frontière au sol (piliers n°1 à n°52) et création d'une zone démilitarisée de 5 km.
\end{itemize}

\textbf{3. Données humanitaires (OCHA / HCR Cameroun, rapport de situation, mars 2025, données fictives) :}
\begin{itemize}
    \item \textbf{12 450} personnes déplacées internes (PDI) dans les arrondissements frontaliers de la Région du Nord.
    \item \textbf{3 centres de santé sur 5} fermés pour insécurité ; rupture de stock de vaccins.
    \item \textbf{Besoins financiers urgents : 2,3 millions USD} (abris, eau, protection, nutrition).
    \item Présence signalée d'éléments armés non étatiques (bandits transfrontaliers) profitant du vide sécuritaire.
\end{itemize}

\textbf{4. Renseignement militaire (État-Major interarmées) :} observation de 2 pelotons (environ 60 hommes) avec appui de mortiers, positionnés à 2 km au nord du pilier n°47 depuis le 10 mars 2025. Violation de la zone démilitarisée de 5 km prévue par l'Accord de 2018.
\end{exemplebox}

\courssub{Identification des principes diplomatiques mobilisables}

Face à ce dossier, l'agent diplomatique doit \cle{qualifier juridiquement les faits} pour sélectionner les principes applicables : c'est l'opération qui consiste à rattacher chaque fait constaté à une règle de droit.

\begin{propriete}[Grille de qualification des faits -- Zone Mbalang]
\begin{center}
\begin{tabularx}{\linewidth}{|X|X|X|}\hline
\rowcolor{popblueL}\textbf{Fait constaté (dossier)} & \textbf{Principe camerounais mobilisé} & \textbf{Base juridique / référence} \\ \hline
Incursion militaire au nord du pilier n°47 & Intégrité territoriale ; intangibilité des frontières & Préambule et art. 45 de la Constitution ; \emph{uti possidetis juris} (CIJ, 2002) \\ \hline
Violation de l'Accord de 2018 (zone démilitarisée) & Respect des engagements internationaux (\emph{pacta sunt servanda}) & Convention de Vienne de 1969, art. 26 \\ \hline
Déplacement de 12 450 civils ; fermeture de centres de santé & Protection des populations ; droit international humanitaire & Conventions de Genève de 1949 ; principes humanitaires \\ \hline
Présence de groupes armés non étatiques & Coopération sécuritaire ; lutte contre l'insécurité transfrontalière & Mécanismes CEMAC et UA ; résolutions du Conseil de sécurité \\ \hline
Besoin de 2,3 millions USD d'assistance & Solidarité internationale ; diplomatie humanitaire & Résolution 46/182 de l'Assemblée générale de l'ONU \\ \hline
Refus du dialogue par Loména (hypothèse) & Règlement pacifique des différends ; médiation ; bons offices & Charte de l'ONU, chap. VI ; Acte constitutif de l'UA \\ \hline
\end{tabularx}
\end{center}
\end{propriete}

\begin{attention}[Piège à éviter : la confusion des genres]
Ne mélangez pas le \emph{ton politique} (discours public, fermeté : « condamne énergiquement ») et le \emph{ton juridique et diplomatique} de la note verbale (précision factuelle, références textuelles, réserve des droits, proposition de procédure). La note verbale est un \textbf{acte juridique formel} ; le communiqué de presse est un \textbf{acte de communication politique}.
\end{attention}

\courssub{Méthode guidée : rédaction de la note de position du MINREX}

La \cle{note verbale} (ou note de position) est l'instrument standard de la correspondance diplomatique formelle entre ministères des Affaires étrangères. Rédigée à la troisième personne, elle engage la responsabilité de l'État.

\begin{methode}[Rédaction d'une note diplomatique en 5 étapes]
\begin{enumerate}
    \item \textbf{Lire tous les documents} : identifier la nature de chaque pièce (juridique, factuelle, chiffrée, politique). Séparer le certain du probable.
    \item \textbf{Extraire les faits clés} : dater, localiser, quantifier (coordonnées, effectifs, montants, dates de traités). Construire une chronologie incontestable.
    \item \textbf{Relier chaque fait à un principe} : utiliser la grille de qualification ci-dessus. Chaque paragraphe « Exposé des faits » doit avoir son pendant « Fondement juridique ».
    \item \textbf{Structurer la note} : respecter le canevas : en-tête $\rightarrow$ références $\rightarrow$ exposé des faits $\rightarrow$ fondements juridiques $\rightarrow$ position et revendications $\rightarrow$ proposition de suite $\rightarrow$ formules de politesse $\rightarrow$ signature.
    \item \textbf{Relire et valider} : vérifier l'exactitude des numéros de traités, des dates, des montants. Contrôler le style administratif (impersonnel, présent de l'indicatif, « le Ministère », « le Gouvernement »). Faire valider par la hiérarchie (Directeur $\rightarrow$ Secrétaire Général $\rightarrow$ Ministre).
\end{enumerate}
\end{methode}

\begin{popfigure}
\centering
\begin{tikzpicture}[font=\scriptsize, scale=0.62, transform shape]
    \draw[popline, line width=1pt] (0,0) rectangle (16, 22);
    \node[align=center, font=\bfseries\scriptsize, text=popblue] at (8, 21.2) {RÉPUBLIQUE DU CAMEROUN \\ Paix -- Travail -- Patrie \\[2pt] MINISTÈRE DES RELATIONS EXTÉRIEURES};
    \draw[popblue, thick] (1, 20.4) -- (15, 20.4);
    \node[align=center] at (8, 19.7) {SECRÉTARIAT GÉNÉRAL \\ DÉPARTEMENT DES AFFAIRES JURIDIQUES ET CONSULAIRES};
    \draw[popmuted] (1, 19.1) rectangle (15, 18.1);
    \node[align=left, anchor=north west] at (1.2, 18.9) {
        \textbf{NOTE VERBALE N° \underline{\hspace{2.5cm}} /MINREX/SG/DAJC} \\[3pt]
        \textbf{Objet :} différend frontalier Zone de Mbalang -- protestation et proposition de règlement. \\[3pt]
        \textbf{Références :} (1) Accord de démarcation 2018 ; (2) Arrêt CIJ 2002 ; (3) Convention de Vienne 1969 ; (4) Rapports OCHA mars 2025.
    };
    \node[anchor=north west, text width=14.5cm, align=justify] at (1.2, 17.8) {Le Ministère des Relations Extérieures de la République du Cameroun présente ses compliments à l'Ambassade de la République de Loména à Yaoundé et a l'honneur de l'informer de ce qui suit :};
    \node[anchor=north west, font=\scriptsize\bfseries, text=popdark] at (1.2, 16.6) {\underline{I. EXPOSÉ DES FAITS}};
    \node[anchor=north west, text width=14.5cm, align=justify] at (1.2, 16.3) {
        1. Le \textbf{12 mars 2025}, des éléments des Forces de Défense de la République de Loména ont franchi la ligne de démarcation matérialisée par le \textbf{pilier frontalier n°47}, pénétrant sur une profondeur de \textbf{2,5 km} dans le territoire camerounais (Zone de Mbalang).\\
        2. Cette incursion viole expressément l'\textbf{article 4 de l'Accord de démarcation signé à Yaoundé le 15 juin 2018}, qui institue une zone démilitarisée de 5 km de part et d'autre de la ligne frontière.\\
        3. Cette situation a provoqué le déplacement de \textbf{12 450 civils} (données OCHA/HCR du 20 mars 2025) et la fermeture de \textbf{3 centres de santé sur 5}, créant une urgence humanitaire estimée à \textbf{2,3 millions USD}.
    };
    \node[anchor=north west, font=\scriptsize\bfseries, text=popdark] at (1.2, 13.6) {\underline{II. FONDEMENTS JURIDIQUES}};
    \node[anchor=north west, text width=14.5cm, align=justify] at (1.2, 13.3) {
        1. L'intangibilité des frontières (\emph{uti possidetis juris}) consacrée par l'\textbf{Arrêt de la CIJ du 10 octobre 2002} (Cameroun c. Nigeria) et réaffirmée par la \textbf{Constitution camerounaise}.\\
        2. Le principe \emph{\textbf{pacta sunt servanda}} (Convention de Vienne sur le droit des traités, 1969, art. 26) oblige la partie loménienne à respecter l'Accord de 2018.\\
        3. L'obligation de protéger les populations civiles (Conventions de Genève de 1949, art. 3 commun ; Résolution 46/182 de l'ONU).
    };
    \node[anchor=north west, font=\scriptsize\bfseries, text=popdark] at (1.2, 11.0) {\underline{III. POSITION DU GOUVERNEMENT CAMEROUNAIS}};
    \node[anchor=north west, text width=14.5cm, align=justify] at (1.2, 10.7) {
        1. \textbf{Condamne} cette violation de l'intégrité territoriale.\\
        2. \textbf{Exige} le retrait immédiat et inconditionnel des forces loméniennes au nord du pilier n°47 avant le \textbf{30 avril 2025}.\\
        3. \textbf{Demande} l'ouverture d'un corridor humanitaire sécurisé pour l'acheminement de l'aide (2,3 millions USD) aux 12 450 déplacés.\\
        4. \textbf{Réserve} le droit de saisir le Conseil de Sécurité de l'ONU et la Cour Internationale de Justice à défaut de règlement pacifique.
    };
    \node[anchor=north west, font=\scriptsize\bfseries, text=popdark] at (1.2, 8.4) {\underline{IV. PROPOSITION DE SUITE}};
    \node[anchor=north west, text width=14.5cm, align=justify] at (1.2, 8.1) {
        Le Ministère propose la tenue d'une \textbf{réunion de la Commission Mixte de Démarcation} sous l'égide de l'Union Africaine dans un délai de \textbf{15 jours} à Yaoundé, pour : (a) vérifier le retrait ; (b) sécuriser le corridor humanitaire ; (c) achever le bornage résiduel.
    };
    \node[anchor=north west, text width=14.5cm, align=justify] at (1.2, 6.2) {
        Le Ministère des Relations Extérieures saisit cette occasion pour renouveler à l'Ambassade de la République de Loména l'assurance de sa haute considération.
    };
    \node[anchor=north west] at (1.2, 5.0) {Fait à Yaoundé, le \underline{\hspace{2.5cm}} 2025};
    \node[anchor=north west] at (1.2, 4.2) {\textbf{Le Ministre des Relations Extérieures}};
    \node[anchor=north west] at (1.2, 3.4) {\textit{Par délégation :} \textbf{Le Secrétaire Général}};
    \node[anchor=north west] at (1.2, 2.4) {\underline{\hspace{4cm}}};
    \draw[poporange, thick, decorate, decoration={brace, amplitude=5pt}] (15.2, 19.1) -- (15.2, 18.1) node[midway, right=4pt, align=left, font=\tiny, text=poporange] {Bloc \textbf{références}\\(traçabilité)};
    \draw[poporange, thick, decorate, decoration={brace, amplitude=5pt}] (15.2, 16.6) -- (15.2, 8.4) node[midway, right=4pt, align=left, font=\tiny, text=poporange] {Corps : \textbf{faits $\rightarrow$ droit}\\\textbf{$\rightarrow$ position}};
    \draw[poporange, thick, decorate, decoration={brace, amplitude=5pt}] (15.2, 6.2) -- (15.2, 5.0) node[midway, right=4pt, align=left, font=\tiny, text=poporange] {Formules \textbf{codifiées}\\(obligatoires)};
\end{tikzpicture}
\legende{Figure 5.2 : Modèle annoté de note verbale MINREX (Zone Mbalang). Les zones à compléter sont soulignées. Observer la progression logique : faits $\rightarrow$ droit $\rightarrow$ position $\rightarrow$ suite.}
\end{popfigure}

\begin{exoresolu}[Rédaction guidée : paragraphe « Position » intégrant des données chiffrées]
\textbf{Consigne :} rédiger le point 3 de la section III (corridor humanitaire) en style administratif, en intégrant les données du dossier.

\textbf{Ébauche d'élève (à corriger) :} « Il faut laisser passer les camions pour les 12 450 déplacés qui ont faim. Il faut 2,3 millions de dollars. On demande un corridor. »
\tcblower
\textbf{Correction en style diplomatique :}
\begin{quote}
3. \textbf{Demande} l'établissement immédiat d'un corridor humanitaire sécurisé et sans entrave, sous l'égide du CICR et de l'OCHA, afin de permettre l'acheminement de l'assistance d'urgence aux \textbf{12 450 personnes déplacées internes} recensées dans les arrondissements frontaliers, pour un besoin financier identifié de \textbf{2,3 millions de dollars des États-Unis}, conformément aux principes d'humanité et d'impartialité énoncés par la Résolution 46/182 de l'Assemblée Générale des Nations Unies.
\end{quote}
\textbf{Analyse de la correction :} usage du présent de l'indicatif (demande) ; vocabulaire précis (\emph{sans entrave, sous l'égide, acheminement, assistance d'urgence}) ; données chiffrées intégrées comme preuves (\textbf{12 450} ; \textbf{2,3 millions USD}) ; renvoi à la source juridique (Résolution 46/182). Le style familier (« il faut », « on », « ont faim ») est entièrement supprimé.
\end{exoresolu}

\courssub{Auto-évaluation : grille de critères pour la note diplomatique}

Avant de rendre votre copie (ou de transmettre la note à la hiérarchie), vous devez l'auditer. Voici une grille d'évaluation inspirée des pratiques de formation des attachés des Affaires étrangères.

\begin{center}
\begin{tabularx}{\linewidth}{|l|X|X|}\hline
\rowcolor{popblueL}\textbf{Critère} & \textbf{Indicateurs de réussite (niveau expert)} & \textbf{Vérification dans le texte} \\ \hline
\textbf{1. Cohérence interne} & Logique faits $\rightarrow$ droit $\rightarrow$ position impeccable, sans contradiction. & Les faits du \S I fondent-ils le droit du \S II ? \\ \hline
\textbf{2. Respect des principes} & Au moins 4 principes cités (intégrité territoriale, \emph{pacta sunt servanda}, DIH, règlement pacifique) et correctement appliqués. & Art. 45 Const. ; art. 26 Conv. Vienne 1969 ; Rés. 46/182 ; chap. VI Charte ONU. \\ \hline
\textbf{3. Argumentation chiffrée} & Au moins \textbf{2 données quantitatives exactes} intégrées comme preuves. & Ex. : 12 450 PDI ; 2,3 M USD ; 2,5 km ; pilier n°47 ; 30 avril 2025. \\ \hline
\textbf{4. Style administratif} & Impersonnel, présent de l'indicatif, formules codifiées, vocabulaire précis. & Pas de « je », « on », « il faut ». « Le Ministère exige », « présente ses compliments ». \\ \hline
\textbf{5. Structure formelle} & En-tête, références, objet, exposé, fondements, position, suite, formules, signature, date, lieu. & Aucune section manquante ; numérotation continue. \\ \hline
\textbf{6. Pertinence de la suite} & Proposition concrète, datée, institutionnelle (Commission mixte, UA, CIJ), réaliste. & À éviter : menace de guerre, sanctions unilatérales, ultimatum sans issue. \\ \hline
\end{tabularx}
\end{center}

\begin{aretenir}[Check-list finale avant signature]
\begin{itemize}
    \item Les \textbf{repères géographiques} (pilier n°47) et \textbf{chronologiques} (12 mars 2025) sont-ils exacts ?
    \item Les \textbf{montants} (2,3 M USD) et \textbf{effectifs} (12 450 PDI) sont-ils sourcés (OCHA, HCR) ?
    \item Les \textbf{références des traités} (Accord 2018 art. 4 ; CIJ 2002 ; Vienne 1969 art. 26) sont-elles complètes ?
    \item La \textbf{formule de politesse finale} est-elle la bonne ? (\emph{« l'assurance de sa haute considération »} pour une Ambassade ; \emph{« la considération distinguée »} pour un ministère étranger).
    \item Le \textbf{choix de l'instrument} (note verbale, lettre d'État, aide-mémoire) est-il justifié ? Ici : note verbale, acte formel standard entre ministères.
\end{itemize}
\end{aretenir}

\begin{savaistu}[Les archives diplomatiques du Cameroun]
Les notes verbales, aides-mémoires, lettres d'État et procès-verbaux de commissions mixtes produits par le MINREX depuis l'indépendance sont conservés aux \textbf{Archives Nationales du Cameroun} à Yaoundé. Ils constituent des sources primaires indispensables pour l'histoire des relations internationales du Cameroun. C'est notamment grâce à ces documents (procès-verbaux de la Commission Mixte Cameroun-Nigeria, cartes et procès-verbaux de bornage) que l'Arrêt de la CIJ sur Bakassi a pu être exécuté, aboutissant au transfert complet de la péninsule au Cameroun en 2013. Comme dans la plupart des États, la consultation des archives diplomatiques récentes est soumise à des délais de déclassification.
\end{savaistu}

\courssub{Exercices d'application et corrigés types}

\begin{exercice}[Entraînement Bac -- Sujet 01 : rédaction partielle]
À partir du dossier « Mbalang » (figure 5.1), rédigez \textbf{uniquement la section II (Fondements juridiques)} de la note verbale. Vous devez y faire apparaître explicitement : le principe \emph{uti possidetis juris}, le principe \emph{pacta sunt servanda} et l'obligation de protection des civils (DIH). Longueur : 15 lignes maximum. Style administratif impératif.
\end{exercice}

\begin{corrige}[Exercice 01 -- Proposition de rédaction]
\underline{II. FONDEMENTS JURIDIQUES}

1. Le Gouvernement de la République du Cameroun rappelle que la frontière séparant les deux États est consacrée par le principe \emph{uti possidetis juris}, tel que confirmé par l'\textbf{Arrêt de la Cour Internationale de Justice du 10 octobre 2002} (affaire Cameroun c. Nigeria) et réaffirmé par la \textbf{Constitution de la République du Cameroun du 18 janvier 1996}.

2. L'\textbf{Accord de démarcation signé à Yaoundé le 15 juin 2018}, en son article 4, institue une zone démilitarisée de 5 km. Conformément au principe \emph{\textbf{pacta sunt servanda}} (Convention de Vienne sur le droit des traités, 1969, article 26), la République de Loména est tenue de l'exécuter de bonne foi.

3. L'incursion ayant causé le déplacement de \textbf{12 450 civils} et la fermeture de structures sanitaires, le Cameroun invoque l'\textbf{article 3 commun aux Conventions de Genève de 1949} et la \textbf{Résolution 46/182 de l'Assemblée Générale des Nations Unies}, qui imposent le respect et la protection des personnes ne participant pas aux hostilités ainsi que l'accès humanitaire impartial.
\end{corrige}

\begin{exercice}[Entraînement Bac -- Sujet 02 : analyse critique]
Une note verbale rédigée par un stagiaire contient la phrase suivante : \emph{« Nous, le Ministère, on est très fâchés parce que vos soldats sont rentrés chez nous sans permission. Il faut qu'ils partent vite sinon on va à l'ONU. »}
Identifiez \textbf{4 erreurs majeures} (forme et fond) par rapport aux standards de la diplomatie camerounaise et proposez la reformulation correcte pour chaque erreur.
\end{exercice}

\begin{corrige}[Exercice 02 -- Grille d'analyse]
\begin{center}
\begin{tabularx}{\linewidth}{|X|X|X|}\hline
\rowcolor{popblueL}\textbf{Erreur identifiée} & \textbf{Règle non respectée} & \textbf{Reformulation correcte} \\ \hline
1. « Nous, le Ministère, on est très fâchés » & Style familier, subjectivité, émotion. & \emph{« Le Ministère des Relations Extérieures condamne fermement... »} \\ \hline
2. « Vos soldats sont rentrés chez nous » & Imprécision juridique ; aucune référence géographique ni textuelle. & \emph{« Des éléments des Forces de Défense loméniennes ont franchi la ligne de démarcation au niveau du pilier n°47... »} \\ \hline
3. « Sans permission » & Absence de qualification juridique de la violation. & \emph{« En violation de l'article 4 de l'Accord de démarcation de 2018 et du principe de l'intangibilité des frontières. »} \\ \hline
4. « Sinon on va à l'ONU » & Menace non diplomatique ; imprécision procédurale (quel organe ? quelle procédure ?). & \emph{« À défaut de retrait dans les délais, le Cameroun se réserve le droit de saisir le Conseil de Sécurité des Nations Unies conformément au Chapitre VI de la Charte. »} \\ \hline
\end{tabularx}
\end{center}
\end{corrige}

\begin{exercice}[Entraînement Bac -- Sujet 03 : mise en situation humanitaire]
Le MINREX doit répondre à une demande d'appui logistique du PAM (Programme Alimentaire Mondial) pour l'acheminement de 500 tonnes de vivres vers les 12 450 PDI de Mbalang. Le corridor humanitaire proposé dans la note verbale (exercice 01) est accepté par Loména. Rédigez la \textbf{note verbale de suivi} (objet : mise en œuvre du corridor humanitaire Mbalang) en 3 points : 1. confirmation de l'accord ; 2. modalités techniques (points de passage, escorte) ; 3. demande de rapport de mission du PAM.
\end{exercice}

\begin{corrige}[Exercice 03 -- Structure attendue]
\textbf{Objet :} mise en œuvre du corridor humanitaire Zone Mbalang -- coordination logistique avec le PAM.

\textbf{Références :} note verbale n°XXX du [date] ; accord de l'Ambassade de Loména du [date] ; demande du PAM Yaoundé n°YYY.

Le Ministère des Relations Extérieures présente ses compliments à l'Ambassade de la République de Loména et a l'honneur de confirmer :

1. \textbf{Accord de principe :} le Gouvernement camerounais confirme l'ouverture du corridor humanitaire pour une durée de 30 jours renouvelable, conformément aux engagements pris.

2. \textbf{Modalités techniques :} point d'entrée : poste frontalier désigné d'un commun accord ; escorte assurée par les forces de défense camerounaises sur le tronçon menant au site des déplacés ; déchargement au site de Mbalang Centre ; application des procédures douanières allégées prévues pour les secours humanitaires.

3. \textbf{Suivi :} le Ministère demande au PAM de transmettre un \textbf{rapport de mission hebdomadaire} (tonnage livré : objectif 500 tonnes ; bénéficiaires atteints : 12 450 PDI ; incidents éventuels) à la division compétente du MINREX, avec copie au MINAS (Affaires sociales) et au MINSANTE (Santé publique).

Le Ministère saisit cette occasion pour renouveler à l'Ambassade l'assurance de sa haute considération.
\end{corrige}

\begin{attention}[Consignes d'examen (Probatoire / Baccalauréat technique)]
\begin{itemize}
    \item \textbf{Lisez le dossier deux fois :} première lecture globale, deuxième lecture « chasse aux chiffres et aux dates » (surlignez-les).
    \item \textbf{Ne recopiez pas le dossier :} la note verbale \textbf{synthétise} et \textbf{qualifie juridiquement}. Elle n'est pas un résumé.
    \item \textbf{Chiffres obligatoires :} une note sans données quantitatives (dates, montants, effectifs, numéros d'articles) est une note vide.
    \item \textbf{Formules de politesse :} leur omission ou leur erreur coûte des points faciles. Apprenez-les par cœur (modèle de la figure 5.2).
    \item \textbf{Gestion du temps :} 30 min de lecture et d'analyse -- 60 min de rédaction -- 15 min de relecture avec la grille d'auto-évaluation.
\end{itemize}
\end{attention}

\newpage

\coursec{Activité d'intégration}

\coursec{Leçon 02 : Les principes de la diplomatie camerounaise}

\courssub{Objectifs de l'activité}

À l'issue de cette activité d'intégration, l'élève sera capable de :
\begin{itemize}
\item mobiliser ses connaissances sur les \cle{Relations Internationales}, les formes de la \cle{coopération internationale} et les \cle{principes de la diplomatie camerounaise} pour analyser une crise réelle ;
\item lire et exploiter des documents de natures différentes (carte, tableau statistique, textes juridiques, communiqué de presse) ;
\item calculer des indicateurs simples (part en pourcentage, taux de variation) et les interpréter ;
\item rédiger et justifier une position officielle sous la forme d'une \cle{note diplomatique} respectant le modèle en usage ;
\item évaluer le travail d'autrui à l'aide d'une grille critériée (évaluation par les pairs).
\end{itemize}

\courssub{Situation}

\textbf{Communiqué du Ministère des Pêches du Nigeria (extrait, mars 2024)} : « Le Gouvernement fédéral du Nigeria rappelle que les eaux situées au large de la péninsule de Bakassi relèvent de sa juridiction historique et annonce le déploiement de patrouilles navales pour y protéger ses pêcheurs. Toute activité de pêche étrangère dans cette zone devra être autorisée par Abuja. »

Quelques jours plus tard, trois pirogues artisanales camerounaises parties de Limbe sont arraisonnées par une vedette nigériane dans la zone dite « B » du Golfe de Guinée. Les pêcheurs de Campo et de Kribi s'inquiètent ; les médias s'enflamment. Le Ministre des Relations Extérieures (MINREX) convoque en urgence une cellule de crise. \textbf{Vous êtes les experts de cette cellule.} Vous devez préparer la position officielle du Cameroun.

\textbf{Document 1 : croquis des zones économiques exclusives (ZEE) du Golfe de Guinée.}

\begin{popfigure}
\begin{center}
\begin{tikzpicture}[scale=0.95, font=\small]
% Mer (fond)
\fill[popblueL!50] (0,0) rectangle (10,6.5);
% Côte (terre) - polygone simplifié
\fill[popgoldL] (0,4.2) -- (2,4.6) -- (4,5.2) -- (6,5.8) -- (8,6.2) -- (10,6.5) -- (10,6.5) -- (0,6.5) -- cycle;
% ZEE Cameroun (lignes vertes en pointillés)
\draw[very thick, popgreen, densely dashed] (2.2,4.5) -- (2.6,2.4) -- (4.4,1.6) -- (5.2,4.9);
% ZEE Nigeria (lignes orange en pointillés)
\draw[very thick, poporange, densely dashed] (0,4.2) -- (1.2,2.8) -- (2.6,2.4);
% ZEE Guinée équatoriale (lignes violettes en pointillés)
\draw[very thick, poppurple, densely dashed] (5.2,4.9) -- (4.4,1.6) -- (6.8,0.8) -- (8,6.2);
% Zone litigieuse B (hachurée rose)
\fill[popink!30, pattern=north east lines, pattern color=popink] (1.2,2.8) -- (2.6,2.4) -- (2.2,4.5) -- (0.8,4.3) -- cycle;
\node[popink, font=\small\bfseries] at (1.6,3.4) {Zone B};
% Villes (points noirs)
\fill[popdark] (1.0,4.35) circle (0.07); \node[above, font=\scriptsize] at (1.0,4.4) {Calabar (Nig.)};
\fill[popdark] (2.4,4.55) circle (0.07); \node[above, font=\scriptsize] at (2.4,4.6) {Limbe};
\fill[popdark] (3.6,4.9) circle (0.07); \node[above, font=\scriptsize] at (3.6,4.95) {Douala};
\fill[popdark] (5.0,5.35) circle (0.07); \node[above, font=\scriptsize] at (5.0,5.4) {Kribi};
\fill[popdark] (6.6,5.9) circle (0.07); \node[above, font=\scriptsize] at (6.6,5.95) {Bata (G. éq.)};
% Légendes ZEE
\node[popgreen, font=\small\bfseries] at (3.7,3.0) {ZEE Cameroun};
\node[poporange, font=\small\bfseries] at (0.9,1.8) {ZEE Nigeria};
\node[poppurple, font=\small\bfseries] at (6.6,2.6) {ZEE Guinée};
\node[poppurple, font=\small\bfseries] at (6.6,2.2) {équatoriale};
% Flèche Nord
\draw[-{Stealth[length=3mm]}, thick, popdark] (9.4,5.2) -- (9.4,6.1); \node[above, font=\scriptsize] at (9.4,6.1) {N};
% Échelle graphique
\draw[thick, popdark] (0.3,0.4) -- (1.8,0.4);
\draw[thick, popdark] (0.3,0.35) -- (0.3,0.45);
\draw[thick, popdark] (1.8,0.35) -- (1.8,0.45);
\node[below, font=\scriptsize] at (1.05,0.3) {0 \quad 100 km};
\end{tikzpicture}
\end{center}
\legende{Croquis schématique (non contractuel) des ZEE du Golfe de Guinée : la zone B, revendiquée par le Nigeria, est située au large de Bakassi et de Limbe.}
\end{popfigure}

\textbf{Document 2 : captures de pêche maritime camerounaise (en milliers de tonnes).}

\begin{center}
\begin{tabularx}{\linewidth}{|l|X|X|X|X|}
\hline
\rowcolor{popblueL} \textbf{Année} & \textbf{Zone A (Sud)} & \textbf{Zone B (litigieuse)} & \textbf{Zone C (Nord)} & \textbf{Total} \\
\hline
2020 & 96 & 28 & 41 & 165 \\
\hline
2021 & 99 & 31 & 44 & 174 \\
\hline
2022 & 103 & 33 & 46 & 182 \\
\hline
2023 & 107 & 36 & 48 & 191 \\
\hline
\end{tabularx}
\end{center}

\textbf{Document 3 : extraits de textes juridiques.}
\begin{itemize}
\item \emph{Convention des Nations Unies sur le droit de la mer (Montego Bay, 1982)} : « L'État côtier exerce, dans la zone économique exclusive, des droits souverains aux fins d'exploration et d'exploitation des ressources naturelles » (art. 56) ; « le règlement des différends se fait par des moyens pacifiques » (art. 279).
\item \emph{Accord de partenariat de pêche UE--Cameroun} : il encadre l'accès des navires européens aux eaux camerounaises contre contrepartie financière, dans le respect de la souveraineté camerounaise.
\end{itemize}

\textbf{Document 4 : rappel des principes de la diplomatie camerounaise} : indépendance et souveraineté nationale ; non-ingérence dans les affaires intérieures des États ; règlement pacifique des différends ; coopération internationale ; respect des traités et engagements internationaux.

\courssub{Consignes}

Travail en groupes de 4 élèves. Durée : 2 séances. Chaque groupe remet un dossier comprenant le croquis, les calculs et la note de position.

\begin{enumerate}
\item \textbf{Localiser.} Reproduire le croquis du Document 1, y placer la zone B, les villes de Limbe, Douala et Kribi, et hachurer la zone revendiquée par le Nigeria. Indiquer l'orientation et l'échelle.
\item \textbf{Calculer.} À partir du Document 2 : a) calculer, pour chaque année, la part (en \%) des captures de la zone B dans le total national ; b) calculer le taux de variation des captures totales entre 2020 et 2023 ; c) conclure en deux phrases sur l'enjeu économique de la zone B pour le Cameroun.
\item \textbf{Analyser.} Identifier, à partir des Documents 3 et 4 : a) le principe juridique qui fonde les droits du Cameroun sur la zone B ; b) deux principes de la diplomatie camerounaise directement applicables à cette crise ; c) la forme de coopération internationale qui pourrait aider à sortir de la crise (coopération bilatérale ou multilatérale ? judiciaire ?).
\item \textbf{Rédiger.} Rédiger une \cle{note diplomatique} d'environ 300 mots adressée au Ministère des Affaires étrangères du Nigeria, selon le modèle officiel : en-tête (MINREX, référence, date), formule protocolaire d'ouverture, rappel des faits, rappel du droit, position du Cameroun, proposition de sortie de crise, formule protocolaire de clôture.
\item \textbf{Évaluer.} Lors de la mise en commun, chaque groupe évalue la note d'un autre groupe à l'aide de la grille fournie (exactitude des faits, exactitude du droit, respect du format, qualité de l'argumentation, correction de la langue — chaque critère noté sur 4).
\end{enumerate}

\courssub{Ressources à mobiliser}

\begin{aretenir}[Outils de l'activité]
\begin{itemize}
\item \textbf{Relations Internationales} : ensemble des relations (politiques, économiques, culturelles, militaires) entre les acteurs de la scène mondiale (États, OIG, ONG, firmes).
\item \textbf{Formes de coopération} : bilatérale (entre deux États) ou multilatérale (au sein d'organisations comme l'ONU, la CEEAC, l'Union africaine) ; elle peut être politique, économique, judiciaire, humanitaire.
\item \textbf{Principes de la diplomatie camerounaise} : souveraineté, non-ingérence, règlement pacifique des différends, coopération, respect des engagements internationaux.
\item \textbf{Calculs} : part en pourcentage $p = \dfrac{\text{partie}}{\text{total}} \times 100$ ; taux de variation $t = \dfrac{V_f - V_i}{V_i} \times 100$.
\end{itemize}
\end{aretenir}

\begin{savaistu}[Dossier 1 : L'assistance humanitaire dans le contexte de la crise]
La crise halieutique en zone B a une dimension humanitaire immédiate : l arraisonnement de pirogues artisanales met en danger la vie de pêcheurs civils, menace la sécurité alimentaire des populations côtières (Limbe, Kribi, Campo) et fragilise l'économie locale. L'\cle{assistance humanitaire} — forme de coopération internationale visant à protéger la vie, la dignité et les moyens d'existence des populations vulnérables — impose aux deux États de garantir la sécurité des pêcheurs, de faciliter leur libération et de préserver l'accès aux ressources halieutiques essentielles. Le Cameroun, en demandant la libération immédiate des équipages, agit conformément au droit international humanitaire et aux principes de solidarité africaine.
\end{savaistu}

\begin{attention}[Pièges fréquents]
Ne confonds pas \emph{mer territoriale} (12 milles marins, souveraineté pleine) et \emph{ZEE} (200 milles marins, droits souverains sur les ressources). Dans la note diplomatique, reste ferme sur le droit mais courtois dans la forme : une note n'est ni un ultimatum ni une supplique. N'oublie jamais les formules protocolaires d'ouverture et de clôture. Vérifie l'orthographe des noms propres (Montego Bay, CEEAC, MINREX).
\end{attention}

\courssub{Démarche et corrigé commenté}

\begin{exoresolu}[Corrigé commenté de la simulation]
\textbf{Étape 1 — Localisation.} La zone B se situe au nord-ouest du Golfe de Guinée, entre la péninsule de Bakassi et Limbe, à cheval sur la limite revendiquée par le Nigeria. Sur le croquis reproduit, elle est hachurée, la flèche du nord pointe vers le haut et l'échelle indique 100 km pour le segment repère. \emph{Commentaire : le croquis doit rester schématique mais cohérent ; il ne prétend pas à une précision cartographique.}

\tcblower

\textbf{Étape 2 — Calculs.} Parts de la zone B dans le total national :
\begin{align*}
2020 &: \frac{28}{165}\times 100 \approx 17{,}0\,\% &
2021 &: \frac{31}{174}\times 100 \approx 17{,}8\,\% \\
2022 &: \frac{33}{182}\times 100 \approx 18{,}1\,\% &
2023 &: \frac{36}{191}\times 100 \approx 18{,}8\,\%
\end{align*}
Taux de variation des captures totales entre 2020 et 2023 :
\[ t = \frac{191 - 165}{165} \times 100 \approx 15{,}8\,\%. \]
\emph{Conclusion : la zone B fournit près d'un cinquième des captures nationales et sa part augmente chaque année ; la perdre priverait le Cameroun d'environ 36 000 tonnes par an et fragiliserait des milliers de pêcheurs de Limbe à Kribi. L'enjeu économique et social est donc majeur.}

\textbf{Étape 3 — Analyse juridique et diplomatique.}
\begin{enumerate}
\item Le fondement juridique est l'article 56 de la Convention de Montego Bay (1982) : l'État côtier exerce des droits souverains sur les ressources de sa ZEE. Le Cameroun, État côtier du Golfe de Guinée et partie à la Convention, exerce ces droits sur la zone B ; l'accord de pêche UE--Cameroun confirme d'ailleurs la reconnaissance internationale de cette juridiction.
\item Principes applicables : (a) le \cle{règlement pacifique des différends} (art. 279 de la Convention et principe cardinal de la diplomatie camerounaise) : nul ne doit recourir à la force, l'arraisonnement des pirogues est donc à dénoncer ; (b) la \cle{souveraineté et l'indépendance nationale}, qui interdisent au Nigeria d'imposer unilatéralement sa juridiction ; on peut ajouter le \cle{respect des traités} et la \cle{coopération internationale}.
\item Forme de coopération adaptée : d'abord une \cle{coopération bilatérale} (commission mixte Cameroun--Nigeria, négociation directe), puis, en cas d'échec, une voie \cle{multilatérale} : médiation de la CEEAC ou de l'Union africaine, voire recours judiciaire devant la Cour internationale de Justice — comme lors de l'affaire de Bakassi tranchée en 2002, dont le Cameroun a respecté l'arrêt, preuve de sa tradition de légalité internationale.
\end{enumerate}

\textbf{Étape 4 — Modèle de note diplomatique (environ 300 mots).}

\emph{Ministère des Relations Extérieures — République du Cameroun\\
Note verbale n° 2024/041/MINREX — Yaoundé, le 15 mars 2024}

Le Ministère des Relations Extérieures de la République du Cameroun présente ses compliments au Ministère des Affaires étrangères de la République fédérale du Nigeria et a l'honneur de porter à sa connaissance ce qui suit.

Le 10 mars 2024, trois pirogues artisanales camerounaises immatriculées à Limbe ont été arraisonnées par une vedette nigériane dans la zone dite « B » du Golfe de Guinée, au motif que cette zone relèverait de la juridiction nigériane.

Le Gouvernement camerounais rappelle que, conformément à l'article 56 de la Convention des Nations Unies sur le droit de la mer du 10 décembre 1982, ratifiée par les deux États, l'État côtier exerce des droits souverains d'exploration et d'exploitation des ressources dans sa zone économique exclusive. La zone B est traditionnellement exploitée par les pêcheurs camerounais ; elle représente en 2023 près de 18,8 \% des captures nationales, soit environ 36 000 tonnes, source de subsistance de plusieurs milliers de familles.

Fidèle aux principes qui fondent sa diplomatie — souveraineté, non-ingérence, règlement pacifique des différends et respect des engagements internationaux —, le Cameroun déplore le recours à des mesures de fait unilatérales et demande la libération immédiate des pêcheurs et de leurs embarcations.

Dans un esprit de bon voisinage et de coopération, le Gouvernement camerounais propose la convocation, dans un délai de trente jours, de la commission mixte Cameroun--Nigeria afin d'examiner la délimitation des espaces maritimes concernés. À défaut d'aboutissement, les deux parties pourraient s'en remettre conjointement à la médiation de la CEEAC ou à une instance judiciaire internationale, conformément à l'article 279 de la Convention.

Le Ministère des Relations Extérieures de la République du Cameroun saisit cette occasion pour renouveler au Ministère des Affaires étrangères de la République fédérale du Nigeria les assurances de sa très haute considération.

\emph{Commentaire : la note respecte le modèle officiel (en-tête, référence, formules protocolaires), mobilise les faits chiffrés calculés à l'étape 2, fonde la position sur le droit et propose une sortie de crise pacifique : c'est l'application concrète des principes de la diplomatie camerounaise.}

\textbf{Étape 5 — Évaluation par les pairs.} Grille proposée (chaque critère sur 4) : exactitude des faits et des chiffres ; exactitude du droit cité ; respect du format de la note ; qualité de l'argumentation ; correction de la langue. Total sur 20. Un groupe obtenant moins de 2/4 à un critère doit reformuler la partie correspondante.
\end{exoresolu}

\courssub{Bilan de l'activité}

\begin{aretenir}[Bilan et compétences mobilisées]
Cette simulation a permis de mettre en évidence trois acquis essentiels. D'abord, les \cle{Relations Internationales} ne sont pas une réalité abstraite : un différend de pêche engage à la fois des États, des textes juridiques, des organisations régionales et des populations concrètes, celles des pêcheurs de Limbe ou de Kribi. Ensuite, les \cle{principes de la diplomatie camerounaise} — souveraineté, non-ingérence, règlement pacifique, coopération, respect des traités — constituent une véritable boîte à outils pour construire une position nationale ferme mais pacifique, à l'image du règlement historique de l'affaire de Bakassi. Enfin, l'élève a pratiqué une démarche complète d'expert : lire des documents variés, cartographier un litige, calculer et interpréter des indicateurs, puis rédiger et justifier un texte officiel argumenté. Ces compétences — analyse documentaire, calcul, argumentation écrite et évaluation critique — sont directement mobilisables à l'examen du Probatoire et du Baccalauréat technique, notamment dans les épreuves de rédaction et de commentaire de documents en ECM.
\end{aretenir}

\newpage

\coursec{Méthodes et exercices résolus}

\coursec{Leçon 02 : Les principes de la diplomatie camerounaise}

\courssub{Le concept de Relations Internationales}

\begin{definition}[Relations Internationales]
Les \cle{Relations Internationales} désignent l'ensemble des rapports (politiques, économiques, culturels, militaires, humanitaires) qui s'établissent entre les acteurs de la scène mondiale — États, organisations internationales (ONU, UA, CEEAC), ONG, entreprises transnationales, populations — au-delà des frontières nationales. Ces rapports peuvent être \cle{pacifiques} (coopération, échanges, traités) ou \cle{conflictuels} (tensions, sanctions, guerres).
\end{definition}

\begin{methode}[Analyser une situation avec le concept de Relations Internationales]
\begin{enumerate}
\item \textbf{Identifier les acteurs} de la situation : États, organisations internationales (ONU, UA, CEEAC), ONG, entreprises transnationales, populations.
\item \textbf{Repérer la nature du lien} qui unit ces acteurs : politique, économique, culturel, militaire, humanitaire.
\item \textbf{Classer la relation} : \cle{bilatérale} (entre deux États) ou \cle{multilatérale} (entre plusieurs États ou au sein d'une organisation).
\item \textbf{Mobiliser la définition} : les Relations Internationales sont l'ensemble des rapports (pacifiques ou conflictuels) qui s'établissent entre les acteurs de la scène mondiale, au-delà des frontières.
\item \textbf{Conclure en justifiant} : montrer en quoi la situation étudiée illustre la définition, en citant précisément les éléments du texte ou de la situation.
\end{enumerate}
\textbf{Réflexe de rédaction} : toujours articuler la copie en trois temps : définition du concept, application aux faits, conclusion justifiée.
\end{methode}

\begin{exoresolu}[Visite officielle et coopération]
\textbf{Énoncé.} En 2023, le Président de la République du Cameroun reçoit à Yaoundé une délégation de chefs d'État de la CEEAC pour discuter de la sécurité dans le bassin du Lac Tchad et de la lutte contre Boko Haram. Un accord de coopération militaire et économique est signé à l'issue de la rencontre.
\begin{enumerate}
\item Définis l'expression « Relations Internationales ».
\item Montre, à partir de cette situation, que le Cameroun entretient des relations internationales, en précisant leurs formes.
\end{enumerate}
\tcblower
\textbf{Solution détaillée.}
\begin{enumerate}
\item \textbf{Définition.} Les Relations Internationales désignent l'ensemble des rapports politiques, économiques, culturels, militaires et humanitaires qui s'établissent entre les acteurs de la scène mondiale (États, organisations internationales, ONG), au-delà des frontières nationales. Ces rapports peuvent être pacifiques (coopération, échanges) ou conflictuels (tensions, guerres).
\item \textbf{Application à la situation.}
\begin{itemize}
\item \emph{Acteurs identifiés} : le Cameroun (État) et les États membres de la CEEAC (organisation sous-régionale). Il s'agit donc d'acteurs étatiques et d'une organisation internationale.
\item \emph{Nature des liens} : la rencontre porte sur la sécurité (lien militaire et politique : lutte contre Boko Haram) et débouche sur un accord de coopération économique (lien économique).
\item \emph{Forme de la relation} : la relation est \cle{multilatérale}, car elle associe plusieurs États réunis au sein d'une organisation sous-régionale (la CEEAC), et non seulement deux États.
\item \emph{Caractère pacifique} : il s'agit d'une coopération (dialogue, signature d'un accord), donc d'une relation pacifique, même si son objet est la lutte contre un conflit.
\end{itemize}
\item \textbf{Conclusion.} Cette situation illustre pleinement le concept de Relations Internationales : le Cameroun, en dialoguant avec ses partenaires de la CEEAC et en signant un accord de coopération militaire et économique, entretient des rapports pacifiques et multilatéraux avec d'autres acteurs de la scène internationale, conformément à la définition.
\end{enumerate}
\end{exoresolu}

\courssub{Aspects et formes de la coopération internationale}

\begin{definition}[Coopération internationale : aspects et formes]
La \cle{coopération internationale} est l'ensemble des actions concertées entre acteurs pour atteindre des objectifs communs. Elle se caractérise par deux entrées complémentaires :
\begin{itemize}
\item Les \cle{aspects} (domaines d'intervention) : politique, économique, culturel, scientifique, militaire, humanitaire.
\item Les \cle{formes} (modalités) : \cle{bilatérale} (deux partenaires) ou \cle{multilatérale} (plusieurs partenaires ou organisation internationale) ; \cle{Nord-Sud} (pays développé vers pays en développement) ou \cle{Sud-Sud} (entre pays en développement).
\end{itemize}
\end{definition}

\begin{aretenir}[Tableau récapitulatif : Aspects et Formes]
\begin{center}
\begin{tabularx}{\linewidth}{|l|X|X|}\hline
\rowcolor{popblueL} \textbf{Critère} & \textbf{Catégories} & \textbf{Indicateurs de repérage} \\ \hline
\textbf{Aspect (Domaine)} & Politique, Économique, Culturel, Scientifique, Militaire, Humanitaire & Objet de l'accord : élections, commerce, bourses, recherche, défense, secours \\ \hline
\textbf{Forme (Nombre)} & Bilatérale (2) / Multilatérale ($\ge$ 3 ou OI) & Compter les États signataires ou l'appartenance à une organisation \\ \hline
\textbf{Forme (Niveau)} & Nord-Sud / Sud-Sud & Niveau de développement des partenaires (PIB/hab, IDH) \\ \hline
\end{tabularx}
\end{center}
\end{aretenir}

\begin{methode}[Classer un exemple de coopération]
\begin{enumerate}
\item \textbf{Rappeler les repères} : les \cle{aspects} de la coopération sont ses domaines (politique, économique, culturel, scientifique, militaire, humanitaire) ; les \cle{formes} sont ses modalités (coopération bilatérale ou multilatérale ; coopération Nord-Sud ou Sud-Sud).
\item \textbf{Repérer le domaine} concerné dans l'exemple : s'agit-il d'échanges commerciaux (économique), d'échanges scolaires (culturel), d'aide d'urgence (humanitaire) ?
\item \textbf{Compter les partenaires} : deux États = bilatérale ; plusieurs États ou une organisation internationale = multilatérale.
\item \textbf{Identifier le sens de la coopération} : entre pays développés et pays en développement = Nord-Sud ; entre pays en développement = Sud-Sud.
\item \textbf{Rédiger la réponse} en nommant précisément l'aspect ET la forme, avec la justification tirée de l'énoncé.
\end{enumerate}
\end{methode}

\begin{exoresolu}[Classer des exemples de coopération]
\textbf{Énoncé.} Pour chacune des situations suivantes, précise l'aspect et la forme de la coopération internationale qu'elle illustre, en justifiant ta réponse.
\begin{enumerate}
\item La Chine finance et construit le stade de football de Limbé au Cameroun.
\item Le Cameroun, le Tchad et le Niger mettent en commun leurs forces armées dans la Force multinationale mixte (FMM) contre Boko Haram.
\item La France accorde des bourses à des étudiants camerounais pour étudier à Lyon.
\end{enumerate}
\tcblower
\textbf{Solution détaillée.}
\begin{enumerate}
\item \textbf{Construction du stade de Limbé par la Chine.}
\begin{itemize}
\item \emph{Aspect} : économique (financement et réalisation d'infrastructures).
\item \emph{Forme} : coopération \cle{bilatérale} (deux États : la Chine et le Cameroun) et coopération \cle{Sud-Sud} (deux pays en développement ou émergents).
\item \emph{Justification} : l'énoncé ne mentionne que deux partenaires et un transfert financier et technique entre eux.
\end{itemize}
\item \textbf{Force multinationale mixte (FMM).}
\begin{itemize}
\item \emph{Aspect} : militaire (mise en commun de forces armées).
\item \emph{Forme} : coopération \cle{multilatérale} (au moins quatre États : Cameroun, Tchad, Niger, Nigeria, Bénin) et coopération \cle{Sud-Sud} (États africains en développement).
\item \emph{Justification} : la pluralité des États engagés dans une structure commune caractérise le multilatéralisme.
\end{itemize}
\item \textbf{Bourses d'études en France.}
\begin{itemize}
\item \emph{Aspect} : culturel et scientifique (formation, échanges universitaires).
\item \emph{Forme} : coopération \cle{bilatérale} (France--Cameroun) et coopération \cle{Nord-Sud} (pays développé vers pays en développement).
\item \emph{Justification} : deux États seulement sont concernés, et l'aide va d'un pays du Nord vers un pays du Sud.
\end{itemize}
\end{enumerate}
\textbf{Conclusion.} Une même coopération se décrit toujours par deux entrées complémentaires : son aspect (le domaine) et sa forme (la modalité). La justification doit s'appuyer sur les éléments précis de l'énoncé.
\end{exoresolu}

\courssub{Dossier 1 : L'assistance humanitaire}

\begin{definition}[Assistance humanitaire]
L'\cle{assistance humanitaire} est l'ensemble des secours d'urgence (vivres, soins médicaux, abris, protection, eau potable) apportés aux populations victimes de catastrophes naturelles, de conflits armés ou de crises complexes, sans discrimination. Elle repose sur quatre \cle{principes fondamentaux} :
\begin{itemize}
\item \cle{Humanité} : sauver des vies, soulager les souffrances.
\item \cle{Impartialité} : aide proportionnée aux besoins, sans distinction de nationalité, race, religion, opinion.
\item \cle{Neutralité} : ne pas prendre parti dans les hostilités ou les controverses politiques.
\item \cle{Indépendance} : autonomie des objectifs humanitaires par rapport aux objectifs politiques, économiques ou militaires.
\end{itemize}
\end{definition}

\begin{methode}[Exploiter des documents sur l'assistance humanitaire]
\begin{enumerate}
\item \textbf{Définir la notion} : l'assistance humanitaire est l'ensemble des aides d'urgence apportées aux populations victimes de catastrophes, conflits ou crises, dans le respect des principes d'humanité, d'impartialité, de neutralité et d'indépendance.
\item \textbf{Identifier dans les documents} : la crise (cause, lieu, victimes), les acteurs de l'aide (États, ONU et ses agences : \cle{HCR} (réfugiés), \cle{PAM} (alimentation), \cle{OMS} (santé), ONG : Croix-Rouge, MSF), les moyens mobilisés.
\item \textbf{Dégager les principes} de l'action humanitaire mis en œuvre ou menacés.
\item \textbf{Repérer les difficultés} éventuelles : accès aux zones de conflit, insuffisance des financements, insécurité des humanitaires, coordination.
\item \textbf{Rédiger une synthèse structurée} : introduction (présentation du dossier), développement (idées croisées des documents), conclusion (appréciation personnelle justifiée).
\end{enumerate}
\end{methode}

\begin{exoresolu}[L'aide aux déplacés internes de l'Extrême-Nord]
\textbf{Énoncé.} Document : « Dans l'Extrême-Nord du Cameroun, les attaques de Boko Haram et les inondations ont contraint des centaines de milliers de personnes à fuir leurs villages. Le HCR et le PAM, avec l'appui du gouvernement camerounais et d'ONG, distribuent des vivres, installent des abris et ouvrent des centres de santé dans les sites de déplacés. Les humanitaires dénoncent toutefois l'insuffisance des fonds et l'insécurité qui limite l'accès à certaines localités. »
\begin{enumerate}
\item Définis l'assistance humanitaire.
\item À partir du document, identifie les acteurs, les actions menées et les difficultés rencontrées.
\end{enumerate}
\tcblower
\textbf{Solution détaillée.}
\begin{enumerate}
\item \textbf{Définition.} L'assistance humanitaire est l'ensemble des secours d'urgence (nourriture, soins médicaux, abris, protection) apportés aux populations victimes de catastrophes naturelles, de conflits armés ou de crises, dans le respect des principes d'humanité, d'impartialité, de neutralité et d'indépendance.
\item \textbf{Exploitation du document.}
\begin{itemize}
\item \emph{La crise} : double origine — conflit armé (attaques de Boko Haram) et catastrophe naturelle (inondations) — provoquant des déplacements massifs de populations (déplacés internes) dans l'Extrême-Nord.
\item \emph{Les acteurs} : les agences des Nations unies (HCR pour la protection/déplacés, PAM pour l'alimentation), le gouvernement camerounais (autorité nationale, coordination) et des ONG (opérationnels de terrain). On note la complémentarité entre acteurs internationaux et acteurs nationaux.
\item \emph{Les actions menées} : distribution de vivres (aide alimentaire/PAM), installation d'abris (hébergement d'urgence/HCR/ONG), ouverture de centres de santé (soins médicaux/OMS/ONG). Ces actions correspondent aux besoins vitaux : manger, se loger, se soigner.
\item \emph{Les difficultés} : insuffisance des financements (appels de fonds sous-financés) et insécurité limitant l'accès humanitaire (principe d'humanité menacé), ce qui empêche l'aide d'atteindre toutes les victimes.
\end{itemize}
\item \textbf{Conclusion.} Ce document montre que l'assistance humanitaire au Cameroun repose sur une coopération multilatérale entre l'État, les agences onusiennes et les ONG, mais qu'elle se heurte à des limites financières et sécuritaires qui appellent une solidarité internationale renforcée et le respect du droit international humanitaire.
\end{enumerate}
\end{exoresolu}

\courssub{Les principes de la diplomatie camerounaise}

\begin{definition}[Diplomatie]
La \cle{diplomatie} est l'ensemble des moyens pacifiques (négociation, dialogue, médiation, recours aux juridictions internationales, signature de traités, action des représentations diplomatiques) par lesquels un État conduit sa politique étrangère, défend ses intérêts et protège ses ressortissants à l'étranger. L'instrument principal au Cameroun est le \cle{MINREX} (Ministère des Relations Extérieures) et son réseau d'ambassades.
\end{definition}

\begin{aretenir}[Les six principes directeurs de la diplomatie camerounaise]
\begin{enumerate}
\item \cle{Souveraineté nationale et indépendance} : le Cameroun décide librement de sa politique intérieure et extérieure.
\item \cle{Non-ingérence} : respect des affaires intérieures des autres États ; refus de toute intervention étrangère dans ses propres affaires.
\item \cle{Règlement pacifique des conflits} : primauté du dialogue, de la négociation, de la médiation, de l'arbitrage ou de la voie judiciaire (CIJ) sur la force.
\item \cle{Coopération internationale} : recherche du développement partagé par le dialogue, les échanges et la solidarité.
\item \cle{Panafricanisme et unité africaine} : engagement actif dans l'UA, la CEEAC, la CEMAC pour l'intégration et la paix en Afrique.
\item \cle{Respect du droit international et des engagements souscrits} : application des traités ratifiés, des résolutions de l'ONU, des arrêts de la CIJ.
\end{enumerate}
\end{aretenir}

\begin{methode}[Justifier une position diplomatique du Cameroun]
\begin{enumerate}
\item \textbf{Définir la diplomatie} : moyens pacifiques par lesquels un État conduit sa politique étrangère.
\item \textbf{Citer les principes de la diplomatie camerounaise} (liste ci-dessus, formulation exacte du cours).
\item \textbf{Relier chaque principe à un fait} de la situation étudiée (exemple : négociation plutôt que guerre = règlement pacifique des conflits).
\item \textbf{Mentionner les instruments} : MINREX, ambassades, traités, adhésion aux organisations internationales (ONU, UA, CEEAC, CEMAC, Commonwealth, OIF).
\item \textbf{Conclure} en montrant la cohérence entre le principe invoqué et l'action menée.
\end{enumerate}
\end{methode}

\begin{exoresolu}[Le règlement du différend de Bakassi]
\textbf{Énoncé.} Après des décennies de tensions avec le Nigeria au sujet de la presqu'île de Bakassi, le Cameroun a saisi la Cour internationale de Justice (CIJ), qui a rendu son arrêt en 2002. Le Cameroun et le Nigeria ont ensuite signé en 2006 les Accords de Greentree, sous l'égide de l'ONU, permettant un transfert pacifique de la souveraineté sur la zone.
\begin{enumerate}
\item Rappelle ce qu'est la diplomatie.
\item Identifie, dans ce règlement, deux principes de la diplomatie camerounaise et justifie ta réponse.
\end{enumerate}
\tcblower
\textbf{Solution détaillée.}
\begin{enumerate}
\item \textbf{Définition.} La diplomatie est l'ensemble des moyens pacifiques — négociation, dialogue, recours au droit, signature de traités, action des ambassades — par lesquels un État conduit ses relations avec les autres États et défend ses intérêts sur la scène internationale.
\item \textbf{Principes mis en œuvre.}
\begin{itemize}
\item \emph{Principe 1 : le règlement pacifique des conflits.} Au lieu de recourir à la force armée, le Cameroun a choisi la voie judiciaire en saisissant la CIJ, puis la voie de la négociation en signant les Accords de Greentree. Le différend territorial s'est réglé sans guerre : c'est l'application directe du principe du règlement pacifique des conflits.
\item \emph{Principe 2 : le respect du droit international et des engagements souscrits.} Le Cameroun a accepté la procédure devant la CIJ et a appliqué l'arrêt de 2002 ainsi que les Accords de 2006 signés sous l'égide de l'ONU. Respecter les décisions de justice internationales et les traités signés traduit l'attachement du Cameroun au droit international.
\item \emph{(Principe complémentaire possible : la souveraineté nationale, car toute la démarche visait à faire reconnaître pacifiquement les droits du Cameroun sur son territoire.)}
\end{itemize}
\item \textbf{Conclusion.} L'affaire de Bakassi est un modèle de diplomatie camerounaise : en privilégiant le droit, la négociation et la coopération avec l'ONU, le Cameroun a défendu sa souveraineté tout en préservant la paix avec son voisin nigérian, conformément à ses principes diplomatiques.
\end{enumerate}
\end{exoresolu}

\courssub{Synthèse et mise en situation : rédaction d'une note diplomatique}

\begin{methode}[Rédiger une note diplomatique]
\begin{enumerate}
\item \textbf{Respecter la structure} : en-tête (expéditeur : ministère ou ambassade ; destinataire ; date et lieu ; objet), formule de courtoisie d'ouverture, corps du message, formule de politesse de clôture, signature.
\item \textbf{Soigner l'objet} : il doit annoncer clairement le sujet (protestation, proposition de coopération, remerciement, demande, information).
\item \textbf{Rédiger le corps} en trois temps : rappel des faits (contexte), position de l'État expéditeur appuyée sur un principe diplomatique, proposition ou demande concrète.
\item \textbf{Adopter un ton} formel, courtois et mesuré, même en cas de désaccord : la diplomatie privilégie le dialogue et la forme.
\item \textbf{Relire} : vérifier la cohérence entre les faits, le principe invoqué et la demande formulée ; vérifier les formules de politesse codifiées.
\end{enumerate}
\end{methode}

\begin{exoresolu}[Note diplomatique de proposition de coopération humanitaire]
\textbf{Énoncé.} À la suite d'inondations qui ont frappé une région frontalière du Cameroun et d'un État voisin, tu es chargé(e) de rédiger, pour le compte du MINREX, une note diplomatique adressée à l'ambassade de cet État voisin, proposant une coordination de l'assistance humanitaire aux populations sinistrées. Rédige cette note en respectant la structure et le ton diplomatiques.
\tcblower
\textbf{Solution détaillée (corrigé type).}

\textbf{Ministère des Relations Extérieures de la République du Cameroun}\\
Yaoundé, le 15 octobre 2025\\
\textbf{À l'attention de l'Ambassade de l'État voisin à Yaoundé}\\
\textbf{Objet : Proposition de coordination de l'assistance humanitaire aux populations sinistrées des inondations.}

Le Ministère des Relations Extérieures de la République du Cameroun présente ses compliments à l'Ambassade et a l'honneur de porter à sa connaissance ce qui suit.

Les inondations récentes ont causé d'importantes pertes humaines et matérielles dans les localités frontalières de nos deux pays, provoquant le déplacement de nombreuses familles qui ont besoin de vivres, d'abris et de soins d'urgence.

Fidèle aux principes de sa diplomatie — la coopération internationale, la solidarité entre États et le règlement concerté des problèmes communs —, le Gouvernement camerounais propose la mise en place d'un mécanisme conjoint de coordination de l'assistance humanitaire : partage des informations sur les zones sinistrées, acheminement commun des secours avec l'appui des agences des Nations unies (PAM, HCR) et sécurisation des couloirs humanitaires.

Le Ministère saurait gré à l'Ambassade de bien vouloir transmettre cette proposition à ses autorités et lui faire connaître, dans les meilleurs délais, les suites qu'elles entendent y réserver.

Le Ministère des Relations Extérieures de la République du Cameroun saisit cette occasion pour renouveler à l'Ambassade les assurances de sa haute considération.

\textbf{Commentaires sur la méthode} : l'en-tête identifie l'expéditeur, le destinataire, la date et l'objet ; le corps suit le plan faits -- position fondée sur un principe -- proposition concrète ; les formules de courtoisie ouvrent et ferment la note ; le ton reste formel et constructif, conforme à l'esprit du dialogue diplomatique.
\end{exoresolu}

\courssub{Méthode 6 : Rédiger et justifier une conclusion argumentée au Bac}

\begin{methode}[Construire une réponse justifiée en ECM]
\begin{enumerate}
\item \textbf{Lire la consigne en entier} et souligner les verbes d'action : « définis », « identifie », « montre », « justifie » — chacun exige un type de réponse différent.
\item \textbf{Rédiger des phrases complètes} : jamais de réponse en un seul mot ; toute réponse contient le terme du cours et son explication.
\item \textbf{Justifier systématiquement} : après chaque affirmation, ajouter « car... », « en effet... », « ce qui se traduit par... » en s'appuyant sur un élément précis de l'énoncé ou du document.
\item \textbf{Structurer} : une idée par paragraphe, avec les connecteurs logiques (d'abord, ensuite, enfin, par conséquent).
\item \textbf{Conclure} par une phrase de synthèse qui répond exactement à la question posée, sans introduire d'idée nouvelle.
\end{enumerate}
\end{methode}

\begin{exoresolu}[Question de synthèse type Bac]
\textbf{Énoncé.} « Le Cameroun est membre de l'ONU, de l'Union africaine, de la CEEAC, de la CEMAC, du Commonwealth et de l'Organisation internationale de la Francophonie. Il accueille régulièrement des sommets sous-régionaux et participe aux opérations de maintien de la paix en Afrique centrale. » Montre, en un paragraphe argumenté et justifié, que l'adhésion du Cameroun aux organisations internationales traduit les principes de sa diplomatie.
\tcblower
\textbf{Solution détaillée.}

\emph{Étape 1 — Analyse de la consigne} : le verbe « montre » exige une démonstration appuyée sur le texte, et « justifié » impose de relier chaque fait à un principe.

\emph{Étape 2 — Rédaction du paragraphe (corrigé type).}

L'adhésion du Cameroun à de nombreuses organisations internationales traduit concrètement les principes de sa diplomatie. D'abord, en siégeant à l'ONU et à l'Union africaine, le Cameroun applique le principe de la coopération internationale, car il accepte de discuter et d'agir avec les autres États plutôt que de rester isolé. Ensuite, sa présence dans les organisations sous-régionales (CEEAC, CEMAC) et sa participation aux opérations de maintien de la paix en Afrique centrale illustrent le principe du règlement pacifique des conflits : en effet, contribuer à la sécurité collective, c'est préférer la prévention et la négociation à la guerre. Enfin, l'appartenance au Commonwealth et à la Francophonie, héritée de son histoire bilingue, montre que le Cameroun défend ses intérêts nationaux tout en restant ouvert au monde, dans le respect du droit international. Par conséquent, chaque adhésion n'est pas un simple symbole : elle met en pratique les principes de coopération, de paix et de souveraineté qui fondent la diplomatie camerounaise.

\emph{Étape 3 — Vérification} : le paragraphe cite les faits du texte (ONU, UA, CEEAC, CEMAC, maintien de la paix), nomme les principes (coopération internationale, règlement pacifique des conflits, souveraineté), les relie par des justifications (« car », « en effet ») et se termine par une phrase de synthèse qui répond à la question.
\end{exoresolu}

\begin{attention}[Les 5 erreurs qui coûtent des points au Bac]
\begin{enumerate}
\item \textbf{Confondre aspect et forme de la coopération.} Dire « la coopération est économique » ne répond qu'à la moitié de la question : il faut toujours préciser aussi la forme (bilatérale ou multilatérale, Nord-Sud ou Sud-Sud). Compter les partenaires dans l'énoncé avant de répondre.
\item \textbf{Répondre sans justifier.} Affirmer « c'est une relation multilatérale » sans dire pourquoi fait perdre la moitié des points : toute réponse doit s'appuyer sur un élément précis du texte ou de la situation (« car trois États sont concernés »).
\item \textbf{Inventer des principes diplomatiques.} Les principes de la diplomatie camerounaise sont limités et précis (souveraineté, non-ingérence, règlement pacifique des conflits, coopération, panafricanisme, respect du droit international) : n'en cite que ceux du cours, avec la formulation exacte.
\item \textbf{Négliger la forme de la note diplomatique.} Oublier l'en-tête, l'objet ou les formules de courtoisie est sanctionné même si le fond est bon : la structure (expéditeur, destinataire, date, objet, corps, politesse, signature) est exigée.
\item \textbf{Recopier le document sans l'exploiter.} Dans un dossier documentaire (assistance humanitaire), recopier des phrases du texte sans les analyser ne rapporte rien : il faut identifier les acteurs, les actions et les difficultés, puis les relier à la définition et aux principes du cours.
\end{enumerate}
\end{attention}

\newpage

\coursec{Exercices}

\coursec{Les principes de la diplomatie camerounaise}

\courssub{Je vérifie mes connaissances}

\begin{exercice}[QCM : Principes directeurs]
Parmi les propositions suivantes, laquelle \textbf{n'est pas} un principe fondamental de la diplomatie camerounaise tel qu'énoncé dans la Constitution et la pratique diplomatique ?
\begin{enumerate}
    \item Le respect de la souveraineté nationale et de l'intégrité territoriale.
    \item Le non-recours à la force et le règlement pacifique des différends.
    \item L'ingérence humanitaire comme droit d'intervention unilatérale.
    \item La coopération régionale et sous-régionale (CEMAC, UA, Commonwealth, OIF).
\end{enumerate}
\end{exercice}

\begin{exercice}[Vrai ou Faux justifié]
Indiquez si les affirmations suivantes sont vraies (V) ou fausses (F) et justifiez brièvement votre réponse en vous appuyant sur le cours.
\begin{enumerate}
    \item Les Relations Internationales se limitent aux seuls rapports entre États souverains.
    \item L'assistance humanitaire est une forme de coopération internationale désintéressée qui ne doit pas servir d'instrument de pression politique.
    \item Le Cameroun a renoncé à tout recours à la force pour régler le différend de Bakassi, privilégiant la voie juridique (CIJ) et diplomatique.
    \item La diplomatie économique vise exclusivement l'attraction des Investissements Directs Étrangers (IDE) sans se soucier du transfert de technologies.
\end{enumerate}
\end{exercice}

\begin{exercice}[Définitions et concepts clés]
Donnez une définition précise et concise (2 à 3 lignes maximum) pour chacun des termes suivants :
\begin{enumerate}
    \item Relations Internationales.
    \item Coopération bilatérale vs coopération multilatérale.
    \item Assistance humanitaire (citer les trois principes directeurs : humanité, neutralité, impartialité).
    \item Diplomatie préventive.
\end{enumerate}
\end{exercice}

\begin{exercice}[Calcul et interprétation rapide]
Le tableau ci-dessous présente l'évolution des flux d'Aide Publique au Développement (APD) reçus par le Cameroun de deux partenaires majeurs (en milliards de FCFA).

\begin{center}
\begin{tabularx}{\linewidth}{|l|X|X|}\hline
\textbf{Année} & \textbf{Partenaire A (Zone Euro)} & \textbf{Partenaire B (Pays émergent)} \\ \hline
2020 & 120 & 45 \\ \hline
2021 & 132 & 54 \\ \hline
2022 & 125 & 68 \\ \hline
\end{tabularx}
\end{center}

\begin{enumerate}
    \item Calculez le taux de variation de l'APD du Partenaire A entre 2020 et 2021.
    \item Calculez le taux de croissance moyen annuel de l'APD du Partenaire B sur la période 2020-2022 (utilisez la formule du taux de croissance moyen annuel : $TCMA = \left(\frac{V_{final}}{V_{initial}}\right)^{\frac{1}{n}} - 1$).
    \item Quelle observation diplomatique pouvez-vous faire sur la diversification des partenariats ?
\end{enumerate}
\end{exercice}

\courssub{Je m'entraîne}

\begin{exercice}[Analyse de situation : Crise humanitaire dans l'Extrême-Nord]
\textbf{Contexte :} Suite aux attaques de groupes armés non étatiques et aux inondations récurrentes, la région de l'Extrême-Nord accueille plus de 400 000 déplacés internes et réfugiés (données HCR/OCHA 2023). Le Gouvernement camerounais, avec l'appui du Système des Nations Unies et d'ONG internationales, met en œuvre le Plan de Réponse Humanitaire (PRH).

\textbf{Questions :}
\begin{enumerate}
    \item Identifiez les trois piliers de l'action humanitaire (selon la résolution 46/182 de l'AGNU) et expliquez comment le principe de \emph{neutralité} s'applique concrètement pour les acteurs humanitaires opérant dans une zone de conflit actif.
    \item En quoi la coordination de cette assistance par le MINAT (Ministère de l'Administration Territoriale) et le MINREX (Ministère des Relations Extérieures) illustre-t-elle le principe de \emph{souveraineté nationale} dans la gestion de l'aide ?
    \item Citez deux formes de coopération internationale mobilisées ici (ex : coopération onusienne, coopération décentralisée, coopération Sud-Sud).
\end{enumerate}
\end{exercice}

\begin{exercice}[Étude de cas : Le processus de Bakassi]
\textbf{Document :} Extrait du Communiqué final du Sommet de Greentree (New York, 12 juin 2006) : « \emph{Les Parties réaffirment leur engagement à respecter l'arrêt de la Cour Internationale de Justice du 10 octobre 2002... et conviennent de mettre en place une Commission Mixte de Suivi pour assurer le transfert pacifique de l'autorité dans la péninsule de Bakassi.} »

\begin{enumerate}
    \item Quel principe du droit international public (et de la diplomatie camerounaise) a prévalu au règlement de ce différend frontalier ? Citez l'organe judiciaire saisi.
    \item Expliquez le rôle de la \emph{diplomatie préventive} et de la \emph{médiation} (sous l'égide du Secrétaire Général de l'ONU) dans l'évitement d'un conflit armé entre le Cameroun et le Nigeria.
    \item Pourquoi le retrait progressif et le transfert d'autorité (phase verte, phase bleue) constituent-ils une application du principe de \emph{règlement pacifique des différends} ?
\end{enumerate}
\end{exercice}

\begin{exercice}[Schéma : Architecture de la coopération camerounaise]
Complétez le schéma ci-dessous en plaçant les éléments suivants dans les cases correspondantes : 
\textit{Union Africaine (UA), CEMAC, Commonwealth, OIF, France, Chine, USA, Allemagne, Coopération décentralisée (Villes jumelées), ONG internationales, Système des Nations Unies (PNUD, UNICEF, FAO), Banque Mondiale, BAD.}

\begin{popfigure}
\begin{center}
\begin{tikzpicture}[node distance=2.2cm and 2.5cm, every node/.style={font=\small, align=center}, >=Stealth]
% Centre
\node[draw, rectangle, rounded corners, fill=popblueL, minimum width=3.5cm, minimum height=1.2cm] (centre) {\textbf{COOPÉRATION CAMEROUNAISE}};

% Bilatérale (Gauche)
\node[draw, rectangle, rounded corners, fill=popgreenL, left=of centre] (bilat) {\textbf{COOPÉRATION BILATÉRALE}};
\node[draw, rectangle, dashed, fill=popcream, below left=of bilat, xshift=-0.5cm, align=center] (bnord){Partenaires traditionnels (Nord)\\\textbf{[À COMPLÉTER]}};
\node[draw, rectangle, dashed, fill=popcream, below right=of bnord, xshift=0.5cm, align=center] (bsud){Partenaires émergents (Sud)\\\textbf{[À COMPLÉTER]}};

% Multilatérale (Droite)
\node[draw, rectangle, rounded corners, fill=poporangeL, right=of centre] (multi) {\textbf{COOPÉRATION MULTILATÉRALE}};
\node[draw, rectangle, dashed, fill=popcream, above right=of multi, xshift=0.5cm, align=center] (mreg){Intégration Régionale/Sous-rég.\\[2pt]\textbf{[À COMPLÉTER]}};
\node[draw, rectangle, dashed, fill=popcream, below right=of multi, xshift=0.5cm, align=center] (mint){Organisations Internationales\\[2pt]\textbf{[À COMPLÉTER]}};
\node[draw, rectangle, dashed, fill=popcream, right=of multi, xshift=0.5cm, align=center] (mfin){Institutions Financières\\[2pt]\textbf{[À COMPLÉTER]}};

% Autres formes (Bas)
\node[draw, rectangle, rounded corners, fill=poppinkL, below=of centre, yshift=-0.5cm] (autre) {\textbf{AUTRES FORMES}};
\node[draw, rectangle, dashed, fill=popcream, left=of autre, xshift=-0.5cm, align=center] (dec){Coop. Décentralisée\\[2pt]\textbf{[À COMPLÉTER]}};
\node[draw, rectangle, dashed, fill=popcream, right=of autre, xshift=0.5cm, align=center] (ong){Société civile / ONG\\[2pt]\textbf{[À COMPLÉTER]}};

% Flèches
\draw[popblue, thick] (centre) -- (bilat);
\draw[popblue, thick] (centre) -- (multi);
\draw[popblue, thick] (centre) -- (autre);
\draw[popmuted] (bilat) -- (bnord);
\draw[popmuted] (bilat) -- (bsud);
\draw[popmuted] (multi) -- (mreg);
\draw[popmuted] (multi) -- (mint);
\draw[popmuted] (multi) -- (mfin);
\draw[popmuted] (autre) -- (dec);
\draw[popmuted] (autre) -- (ong);
\end{tikzpicture}
\end{center}
\legende{Architecture de la coopération internationale du Cameroun (schéma à compléter).}
\end{popfigure}
\end{exercice}

\begin{exercice}[Rédaction administrative : Note verbale]
Vous êtes attaché(e) à la Direction des Organisations Internationales au MINREX. Rédigez une \textbf{Note Verbale} (formule type) adressée à la Représentation Résidente du PNUD à Yaoundé pour :
\begin{itemize}
    \item Remercier pour l'appui technique au processus d'élaboration de la Stratégie Nationale de Développement 2030 (SND30).
    \item Solliciter l'appui du PNUD pour l'organisation d'une table ronde des bailleurs de fonds (données : date prévisionnelle, lieu, thème « Investir dans le capital humain »).
    \item Rappeler l'engagement du Cameroun à atteindre les ODD (Objectifs de Développement Durable).
\end{itemize}
Respectez les conventions diplomatiques (en-tête, formules d'appel et de clôture, cachet).
\end{exercice}

\courssub{Je me prépare au Bac}

\begin{exercice}[Sujet Type Bac : Diplomatie humanitaire et principes directeurs (12 pts)]
\textbf{Document 1 :} Extrait de la Loi n°96/06 du 18 janvier 1996 portant révision de la Constitution (Préambule) : « \emph{Le Peuple du Cameroun... affirme son attachement aux libertés fondamentales inscrites dans la Déclaration universelle des droits de l'homme, la Charte des Nations Unies et la Charte africaine des droits de l'homme et des peuples... Le Peuple du Cameroun proclame que la Nation protège les minorités et préserve les droits des peuples autochtones...} »

\textbf{Document 2 :} Carte schématique de la région de l'Extrême-Nord (fournie à l'examen) montrant : la zone du Lac Tchad, les axes routiers Maroua-Kousseri, Mora-Kousseri, les localités d'accueil de déplacés (Minawao, Zamai, Mémé), les zones d'insécurité (frontière nigériane, îles du Lac Tchad), les bureaux de coordination OCHA/UNHCR à Maroua.

\textbf{Document 3 :} Tableau statistique : « Évolution du financement du Plan de Réponse Humanitaire (PRH) Extrême-Nord 2020-2023 » (Besoins en M USD vs Fonds reçus en M USD vs Taux de couverture \%).
\begin{center}
\begin{tabular}{|c|c|c|c|}\hline
\textbf{Année} & \textbf{Besoins (M \$US)} & \textbf{Fonds reçus (M \$US)} & \textbf{Couverture (\%)} \\ \hline
2020 & 320 & 144 & 45 \\ \hline
2021 & 350 & 175 & 50 \\ \hline
2022 & 380 & 152 & 40 \\ \hline
2023 & 410 & 205 & 50 \\ \hline
\end{tabular}
\end{center}

\textbf{Questions :}
\begin{enumerate}
    \item \textbf{(2 pts)} À partir du Document 1, citez deux principes constitutionnels qui fondent l'action humanitaire de l'État camerounais envers les populations vulnérables de l'Extrême-Nord.
    \item \textbf{(3 pts)} En vous appuyant sur le Document 2 (carte) et vos connaissances, expliquez comment le principe de \emph{non-ingérence} dans les affaires intérieures s'articule avec l'acceptation par le Cameroun de l'assistance humanitaire internationale sur son territoire. Donnez un exemple concret de coordination (ex : cluster protection, cluster abris).
    \item \textbf{(3 pts)} À partir du Document 3, calculez le montant moyen annuel des besoins non couverts sur la période 2020-2023. Commentez l'évolution du taux de couverture et proposez deux leviers de la \emph{diplomatie économique et humanitaire} que le MINREX pourrait actionner pour améliorer le financement.
    \item \textbf{(4 pts)} Rédigez un paragraphe argumenté (15-20 lignes) montrant en quoi la gestion de la crise de l'Extrême-Nord illustre la mise en œuvre cohérente des principes de la diplomatie camerounaise : \emph{souveraineté, solidarité africaine, partenariat multilatéral, règlement pacifique et protection des droits humains}.
\end{enumerate}
\end{exercice}

\begin{exercice}[Sujet Type Bac : Diplomatie économique et médiation régionale (10 pts)]
\textbf{Contexte :} Le Cameroun souhaite accélérer sa transformation structurelle (SND30) en attirant des investissements dans le corridor Kribi-Brazzaville (port en eau profonde, chemin de fer, zone industrielle). Un différend persiste sur la délimitation maritime avec la Guinée Équatoriale, freinant l'exploitation conjointe d'hydrocarbures transfrontaliers.

\textbf{Données :} Flux d'IDE entrants (Milliards FCFA) : 2019 : 320 ; 2020 : 280 ; 2021 : 350 ; 2022 : 410 ; 2023 : 390. Principaux pourvoyeurs : France, Chine, Maroc, Belgique, Canada.

\textbf{Questions :}
\begin{enumerate}
    \item \textbf{(2 pts)} Calculez le Taux de Croissance Moyen Annuel (TCMA) des IDE sur la période 2019-2023. Interprétez le résultat au regard de la diplomatie économique.
    \item \textbf{(3 pts)} Le Ministre des Relations Extérieures mandate l'Ambassadeur du Cameroun à Malabo pour proposer une \emph{Commission Mixte de Développement Transfrontalier} inspirée du modèle de la Commission Mixte Cameroun-Nigeria (CMCN). Expliquez en quoi cette démarche relève de la \emph{diplomatie préventive} et du \emph{règlement pacifique des différends} (Charte de l'UA, Acte constitutif de la CEMAC).
    \item \textbf{(3 pts)} Identifiez trois avantages comparatifs que le Cameroun peut valoriser dans sa \emph{diplomatie de lobbying} auprès des investisseurs du Golfe et d'Asie pour le corridor Kribi-Brazzaville (infrastructures, ressources, marchés).
    \item \textbf{(2 pts)} Rédigez l'objet et les trois premiers paragraphes (exposé des motifs, proposition concrète, appel à l'action) d'une \textbf{Lettre Officielle} du MINREX à son homologue de Guinée Équatoriale pour la tenue d'un Sommet bilatéral « Paix, Sécurité et Développement partagé ».
\end{enumerate}
\end{exercice}

\begin{exercice}[Sujet Type Bac : Étude documentaire synthétique - La place du Cameroun dans la gouvernance mondiale (8 pts)]
\textbf{Document A :} Discours du Président de la République à la 78e AGNU (New York, septembre 2023) : « \emph{...La réforme du Conseil de Sécurité n'est plus une option, c'est une impératif de justice... L'Afrique réclame deux sièges permanents avec droit de veto... Le Cameroun réaffirme son engagement pour un multilatéralisme rénové, efficace et représentatif...} »

\textbf{Document B :} Participation camerounaise aux Opérations de Maintien de la Paix (OMP) de l'ONU (2023) : 850 policiers et militaires déployés (MINUSCA, MINUSMA - retrait en cours, MONUSCO). Rang : 1er contributeur africain francophone, top 20 mondial.

\textbf{Document C :} Candidature du Cameroun au Conseil des Droits de l'Homme (CDH) pour la période 2025-2027. Priorités affichées : Droit au développement, Protection des minorités, Lutte contre les discours de haine, Éducation à la citoyenneté.

\textbf{Questions :}
\begin{enumerate}
    \item \textbf{(2 pts)} À partir du Document A, expliquez la position diplomatique du Cameroun sur la réforme de la gouvernance mondiale. Quel groupe de négociation (ex : Groupe africain, C10, G77+Chine) porte cette revendication ?
    \item \textbf{(2 pts)} En vous appuyant sur le Document B, montrez comment la contribution aux OMP traduit le principe de \emph{solidarité internationale} et de \emph{responsabilité de protéger (R2P)}. Citez une contrainte majeure pour le contingent camerounais (ex : équipement, formation, mandat robuste).
    \item \textbf{(2 pts)} Le Document C mentionne la « Lutte contre les discours de haine ». Reliez cette priorité à la situation intérieure (crise anglophone, Extrême-Nord) et à l'engagement du Cameroun vis-à-vis du \emph{Pacte mondial pour des migrations sûres, ordonnées et régulières} (Marrakech 2018).
    \item \textbf{(2 pts)} Synthèse : En 10 lignes maximum, démontrez que l'action du Cameroun sur la scène internationale (Documents A, B, C) est guidée par une logique de \emph{diplomatie d'influence} au service de l'intérêt national et de la paix collective.
\end{enumerate}
\end{exercice}

\newpage

\coursec{Corrigés des exercices}

\begin{corrige}[Exercice 1 — QCM : Principes directeurs]
\textbf{Réponse : 3.} L'ingérence humanitaire comme droit d'intervention unilatérale n'est pas un principe de la diplomatie camerounaise.

\textbf{Justification :} La diplomatie camerounaise repose sur le respect de la souveraineté des États, la non-ingérence dans les affaires intérieures, le règlement pacifique des différends et la coopération internationale. Le principe de non-ingérence, inscrit dans la Charte des Nations Unies (article 2, paragraphe 7) et dans l'Acte constitutif de l'Union Africaine, exclut toute intervention unilatérale sur le territoire d'un État sans son consentement ou sans mandat du Conseil de Sécurité. L'assistance humanitaire, elle, n'est légitime que si elle respecte les principes d'humanité, de neutralité et d'impartialité, avec l'accord de l'État d'accueil.

\textbf{Point de vigilance :} ne pas confondre « assistance humanitaire » (consensuelle, encadrée) et « ingérence humanitaire » (intervention unilatérale, contraire au droit international classique).
\end{corrige}

\begin{corrige}[Exercice 2 — Vrai ou Faux justifié]
\begin{enumerate}
    \item \textbf{Faux.} Les Relations Internationales dépassent les seuls rapports interétatiques : elles incluent les organisations internationales (ONU, UA), les ONG, les firmes multinationales, les collectivités territoriales (coopération décentralisée) et même les individus (diasporas).
    \item \textbf{Vrai.} L'assistance humanitaire obéit aux principes d'humanité, de neutralité et d'impartialité (résolution 46/182 de l'AGNU) : elle vise à soulager la souffrance sans discrimination et ne doit pas être instrumentalisée à des fins politiques, militaires ou économiques.
    \item \textbf{Vrai.} Dans le différend de Bakassi, le Cameroun a choisi la voie juridique en saisissant la Cour Internationale de Justice (arrêt du 10 octobre 2002), puis la voie diplomatique (accords de Greentree, 12 juin 2006), illustrant le principe du règlement pacifique des différends.
    \item \textbf{Faux.} La diplomatie économique ne vise pas seulement l'attraction des IDE : elle poursuit aussi le transfert de technologies, la formation de la main-d'œuvre, l'ouverture de marchés pour les exportations camerounaises et l'appui à la transformation structurelle (SND30).
\end{enumerate}
\end{corrige}

\begin{corrige}[Exercice 3 — Définitions et concepts clés]
\begin{enumerate}
    \item \textbf{Relations Internationales :} ensemble des relations politiques, économiques, diplomatiques, culturelles et militaires qui s'établissent entre les acteurs de la scène mondiale (États, organisations internationales, ONG, firmes).
    \item \textbf{Coopération bilatérale :} coopération entre deux États (ex. Cameroun-France). \textbf{Coopération multilatérale :} coopération réunissant plusieurs États au sein d'organisations internationales (ex. Cameroun-ONU, Cameroun-CEMAC).
    \item \textbf{Assistance humanitaire :} aide d'urgence (vivres, soins, abris, protection) apportée aux populations victimes de crises, guidée par trois principes : \emph{humanité} (soulager la souffrance), \emph{neutralité} (ne pas prendre parti dans le conflit), \emph{impartialité} (aider selon les seuls besoins, sans discrimination).
    \item \textbf{Diplomatie préventive :} ensemble des actions diplomatiques (médiation, bons offices, négociation) menées pour empêcher l'éclatement ou l'escalade d'un conflit avant qu'il ne devienne armé.
\end{enumerate}
\end{corrige}

\begin{corrige}[Exercice 4 — Calcul et interprétation rapide]
\begin{enumerate}
    \item \textbf{Taux de variation du Partenaire A (2020-2021) :}
    \[ t = \frac{132 - 120}{120} \times 100 = \frac{12}{120} \times 100 = 10\,\% \]
    L'APD du Partenaire A a augmenté de \textbf{10\,\%} entre 2020 et 2021.
    \item \textbf{TCMA du Partenaire B (2020-2022), } $n = 2$ ans :
    \[ TCMA = \left(\frac{68}{45}\right)^{\frac{1}{2}} - 1 = \sqrt{\{1,5111\}} - 1 \approx \{1,2293\} - 1 = \{0,2293\} \]
    Soit un taux de croissance moyen annuel d'environ \textbf{22,9\,\%}.
    \item \textbf{Observation diplomatique :} l'APD du partenaire traditionnel (Zone Euro) stagne (125 en 2022, proche de 120 en 2020), tandis que celle du partenaire émergent progresse fortement (+51\,\% sur deux ans). Cela traduit la \emph{diversification des partenariats} du Cameroun, qui ne dépend plus exclusivement de ses partenaires traditionnels du Nord et développe la coopération Sud-Sud.
\end{enumerate}
\textbf{Point de vigilance :} dans le TCMA, $n$ est le nombre d'années écoulées (ici 2, et non 3).
\end{corrige}

\begin{corrige}[Exercice 5 — Analyse de situation : Crise humanitaire dans l'Extrême-Nord]
\begin{enumerate}
    \item \textbf{Les trois piliers (résolution 46/182) :} l'\emph{humanité} (la souffrance humaine doit être soulagée partout où elle se trouve), la \emph{neutralité} (les acteurs humanitaires ne prennent parti pour aucun camp dans un conflit) et l'\emph{impartialité} (l'aide est fournie selon les seuls besoins, sans discrimination de nationalité, d'ethnie, de religion ou d'opinion). Concrètement, dans une zone de conflit actif, la neutralité impose aux ONG et agences onusiennes de ne pas armer, financer ou soutenir publiquement un camp, de dialoguer avec toutes les parties pour garantir l'accès aux populations, et de ne pas escorter militairement leurs convois au profit d'un belligérant.
    \item \textbf{Souveraineté nationale :} la coordination par le MINAT et le MINREX montre que l'aide internationale ne s'impose pas au Cameroun : c'est l'État qui autorise, encadre et coordonne l'action des partenaires sur son territoire (enregistrement des ONG, validation des plans de réponse, contrôle des accès). L'assistance humanitaire s'exerce donc \emph{avec le consentement} de l'État d'accueil, dans le respect de sa souveraineté.
    \item \textbf{Formes de coopération mobilisées :} la \emph{coopération multilatérale onusienne} (HCR, OCHA, PNUD, UNICEF) et la \emph{coopération avec la société civile internationale} (ONG internationales). On peut aussi citer la coopération bilatérale (bailleurs États) et la coopération sous-régionale (gestion des réfugiés nigérians et tchadiens).
\end{enumerate}
\end{corrige}

\begin{corrige}[Exercice 6 — Étude de cas : Le processus de Bakassi]
\begin{enumerate}
    \item Le principe qui a prévalu est le \emph{règlement pacifique des différends} par la voie juridictionnelle, conformément à la Charte des Nations Unies (article 33). L'organe judiciaire saisi est la \textbf{Cour Internationale de Justice} (CIJ), qui a rendu son arrêt le 10 octobre 2002.
    \item La \emph{diplomatie préventive} s'est exercée à travers les bons offices du Secrétaire Général de l'ONU (Kofi Annan), qui a organisé les sommets de Greentree et facilité le dialogue direct entre les chefs d'État camerounais et nigérian. La \emph{médiation} a permis de transformer un contentieux potentiellement explosif en processus négocié : création de la Commission Mixte Cameroun-Nigeria (CMCN), établissement d'un calendrier de retrait, mécanismes de suivi. Ainsi, l'escalade militaire a été évitée.
    \item Le retrait progressif et le transfert d'autorité par phases (phase verte : retrait des forces nigérianes et administration transitoire ; phase bleue : exercice progressif de la souveraineté camerounaise) illustrent le règlement pacifique car la souveraineté est transférée \emph{sans recours à la force}, par étapes négociées, contrôlées par la CMCN et garanties par la communauté internationale (témoins : ONU, Allemagne, USA, France, Royaume-Uni).
\end{enumerate}
\textbf{Point de vigilance :} bien distinguer la phase juridictionnelle (arrêt CIJ de 2002) de la phase diplomatique de mise en œuvre (Greentree 2006, CMCN).
\end{corrige}

\begin{corrige}[Exercice 7 — Schéma : Architecture de la coopération camerounaise]
\textbf{Complétion attendue du schéma :}
\begin{itemize}
    \item \textbf{Coopération bilatérale — Partenaires traditionnels (Nord) :} France, USA, Allemagne.
    \item \textbf{Coopération bilatérale — Partenaires émergents (Sud) :} Chine.
    \item \textbf{Coopération multilatérale — Intégration régionale et sous-régionale :} Union Africaine (UA), CEMAC.
    \item \textbf{Coopération multilatérale — Organisations internationales :} Commonwealth, OIF, Système des Nations Unies (PNUD, UNICEF, FAO).
    \item \textbf{Coopération multilatérale — Institutions financières :} Banque Mondiale, BAD (Banque Africaine de Développement).
    \item \textbf{Autres formes — Coopération décentralisée :} Villes jumelées (collectivités territoriales).
    \item \textbf{Autres formes — Société civile / ONG :} ONG internationales.
\end{itemize}
\textbf{Point de vigilance :} le Commonwealth et l'OIF sont des organisations multilatérales (et non des partenaires bilatéraux), même si le Cameroun y entretient des relations privilégiées avec certains membres. La Chine relève ici de la coopération Sud-Sud (pays émergent).
\end{corrige}

\begin{corrige}[Exercice 8 — Rédaction administrative : Note verbale]
\textbf{Éléments attendus (corrigé type) :}

\textit{MINISTÈRE DES RELATIONS EXTÉRIEURES — Direction des Organisations Internationales — Yaoundé, le [date] — N°\ldots/MINREX/DOI}

La Direction des Organisations Internationales du Ministère des Relations Extérieures de la République du Cameroun présente ses compliments à la Représentation Résidente du Programme des Nations Unies pour le Développement (PNUD) à Yaoundé et a l'honneur de la remercier chaleureusement pour l'appui technique appréciable apporté au processus d'élaboration de la Stratégie Nationale de Développement 2030 (SND30).

Dans le prolongement de cette fructueuse collaboration, le Ministère sollicite l'accompagnement du PNUD pour l'organisation d'une table ronde des bailleurs de fonds, prévue le [date] à [lieu], sur le thème « Investir dans le capital humain ». Cette rencontre vise à mobiliser les ressources nécessaires à la mise en œuvre des priorités de la SND30.

Le Ministère rappelle à cette occasion l'engagement résolu du Gouvernement camerounais en faveur de l'atteinte des Objectifs de Développement Durable (ODD), et compte sur l'expertise reconnue du PNUD pour le succès de cette initiative.

La Direction des Organisations Internationales du Ministère des Relations Extérieures saisit cette occasion pour renouveler à la Représentation Résidente du PNUD les assurances de sa haute considération. \textit{(Cachet et signature)}

\textbf{Points de vigilance :} la note verbale s'écrit à la troisième personne, sans « je » ni « nous » direct ; elle s'ouvre par « présente ses compliments » et se clôt par « renouvelle\ldots les assurances de sa haute considération » ; elle porte un numéro d'enregistrement, un cachet et une signature.
\end{corrige}

\begin{corrige}[Exercice 9 — Sujet Type Bac : Diplomatie humanitaire et principes directeurs (12 pts)]
\textbf{Barème indicatif :} Q1 : 2 pts ; Q2 : 3 pts ; Q3 : 3 pts ; Q4 : 4 pts.

\begin{enumerate}
    \item \textbf{(2 pts)} Deux principes constitutionnels : l'\emph{attachement aux libertés fondamentales} inscrites dans la Déclaration universelle des droits de l'homme, la Charte des Nations Unies et la Charte africaine des droits de l'homme et des peuples (1 pt) ; la \emph{protection des minorités} et la préservation des droits des peuples autochtones (1 pt). Ces principes obligent l'État à protéger et assister les populations vulnérables, y compris les déplacés internes de l'Extrême-Nord.
    \item \textbf{(3 pts)} Le principe de non-ingérence interdit toute intervention étrangère dans les affaires intérieures du Cameroun ; mais l'assistance humanitaire n'est pas une ingérence car elle est \emph{sollicitée et acceptée} par l'État, qui en garde la coordination (MINAT, MINREX) (1,5 pt). L'État conserve donc sa souveraineté tout en bénéficiant de la solidarité internationale. Exemple de coordination : le \emph{cluster protection} (coordonné par le HCR avec les services de l'État pour l'enregistrement des déplacés) ou le \emph{cluster abris} (construction de shelters à Minawao sous supervision des autorités administratives) (1,5 pt).
    \item \textbf{(3 pts)} Besoins non couverts = Besoins $-$ Fonds reçus :
    \begin{align*}
    2020 &: 320 - 144 = 176 \\
    2021 &: 350 - 175 = 175 \\
    2022 &: 380 - 152 = 228 \\
    2023 &: 410 - 205 = 205
    \end{align*}
    Moyenne annuelle : $\dfrac{176 + 175 + 228 + 205}{4} = \dfrac{784}{4} = 196$ M \$US (1 pt).
    Commentaire : le taux de couverture reste insuffisant et instable (45\,\%, 50\,\%, 40\,\%, 50\,\%) ; il ne dépasse jamais la moitié des besoins, avec un recul marqué en 2022 (1 pt). Leviers de diplomatie économique et humanitaire : organiser une \emph{table ronde des bailleurs} et intensifier le plaidoyer auprès des partenaires traditionnels et émergents ; diversifier les sources de financement (fonds humanitaires régionaux, coopération Sud-Sud, diaspora, secteur privé) (1 pt).
    \item \textbf{(4 pts)} \textbf{Paragraphe argumenté (attendus) :} la gestion de la crise de l'Extrême-Nord illustre la cohérence des principes de la diplomatie camerounaise. D'abord, la \emph{souveraineté} : l'État coordonne l'aide via le MINAT et le MINREX, et aucune assistance ne s'impose sans son consentement. Ensuite, la \emph{solidarité africaine} : le Cameroun accueille des réfugiés nigérians et tchadiens et coopère avec ses voisins dans la Commission Mixte et la Force Multinationale Mixte contre Boko Haram, dans un esprit de solidarité sous-régionale promu par l'UA et la CEMAC. Puis le \emph{partenariat multilatéral} : le Plan de Réponse Humanitaire est mis en œuvre avec le Système des Nations Unies (HCR, OCHA, PAM, UNICEF) et les ONG internationales, traduisant l'attachement du Cameroun au multilatéralisme. Le \emph{règlement pacifique} : le Cameroun privilégie le dialogue et les mécanismes conjoints de sécurité transfrontalière plutôt que l'escalade unilatérale. Enfin, la \emph{protection des droits humains} : fidèle à son préambule constitutionnel, l'État protège les déplacés, les minorités et les populations vulnérables, avec l'appui des clusters protection. Ainsi, la crise de l'Extrême-Nord montre une diplomatie qui concilie souveraineté nationale, ouverture internationale et primauté du droit.
\end{enumerate}
\textbf{Points de vigilance :} citer explicitement le consentement de l'État comme clé d'articulation entre non-ingérence et assistance ; vérifier les calculs ligne par ligne ; en Q4, mobiliser les cinq principes nommés dans la consigne avec un exemple concret pour chacun.
\end{corrige}

\begin{corrige}[Exercice 10 — Sujet Type Bac : Diplomatie économique et médiation régionale (10 pts)]
\textbf{Barème indicatif :} Q1 : 2 pts ; Q2 : 3 pts ; Q3 : 3 pts ; Q4 : 2 pts.

\begin{enumerate}
    \item \textbf{(2 pts)} TCMA des IDE 2019-2023, $n = 4$ :
    \[ TCMA = \left(\frac{390}{320}\right)^{\frac{1}{4}} - 1 = (\{1,21875\})^{0,25} - 1 \approx \{1,0507\} - 1 = \{0,0507\} \]
    Soit environ \textbf{5,1\,\%} par an (1 pt). Interprétation : malgré le recul de 2020 (crise sanitaire), les IDE progressent en moyenne, signe que la diplomatie économique (plaidoyer, forums économiques, accords bilatéraux) porte des fruits, mais la croissance reste modérée et volatile (390 en 2023 contre 410 en 2022) (1 pt).
    \item \textbf{(3 pts)} Proposer une Commission Mixte de Développement Transfrontalier relève de la \emph{diplomatie préventive} : il s'agit d'agir \emph{avant} que le différend maritime ne dégénère en crise ouverte, en créant un cadre permanent de dialogue (1,5 pt). Elle relève aussi du \emph{règlement pacifique des différends} consacré par la Charte de l'ONU (art. 33), l'Acte constitutif de l'UA et les textes de la CEMAC : négociation, médiation et mécanismes conjoints remplacent le recours à la force, à l'image du succès de la CMCN après Bakassi (1,5 pt).
    \item \textbf{(3 pts)} Trois avantages comparatifs : (a) les \emph{infrastructures} : port en eau profonde de Kribi, unique en Afrique centrale, futur axe ferroviaire vers Brazzaville ; (b) les \emph{ressources} : hydrocarbures, minerais (fer, bauxite), potentiel agricole et forestier ; (c) les \emph{marchés} : position de hub de la CEMAC (environ 50 millions de consommateurs), façade atlantique desservant les pays enclavés (Tchad, Centrafrique) (1 pt par avantage justifié).
    \item \textbf{(2 pts)} \textbf{Lettre officielle (extrait type) :}
    \textit{Objet : Proposition de tenue d'un Sommet bilatéral « Paix, Sécurité et Développement partagé ».}
    Monsieur le Ministre, j'ai l'honneur de porter à votre connaissance la volonté du Gouvernement de la République du Cameroun de renforcer les liens d'amitié et de coopération qui unissent nos deux pays frères, dans un contexte où la stabilité du Golfe de Guinée conditionne notre développement commun (exposé des motifs). À cette fin, le Cameroun propose la tenue, à Yaoundé ou à Malabo selon votre convenance, d'un Sommet bilatéral « Paix, Sécurité et Développement partagé », assorti de la création d'une Commission Mixte de Développement Transfrontalier chargée d'examiner les questions de délimitation maritime et d'exploitation conjointe des ressources (proposition concrète). Convaincu que le dialogue est la voie privilégiée du règlement pacifique des différends, conformément à l'Acte constitutif de l'Union Africaine, je vous invite à accepter cette initiative et vous prie d'agréer, Monsieur le Ministre, l'expression de ma haute considération (appel à l'action).
\end{enumerate}
\textbf{Points de vigilance :} dans le TCMA, $n = 4$ (nombre d'années écoulées entre 2019 et 2023) ; citer au moins un texte juridique en Q2 ; la lettre officielle emploie la première personne (« j'ai l'honneur »), contrairement à la note verbale.
\end{corrige}

\begin{corrige}[Exercice 11 — Sujet Type Bac : La place du Cameroun dans la gouvernance mondiale (8 pts)]
\textbf{Barème indicatif :} Q1 : 2 pts ; Q2 : 2 pts ; Q3 : 2 pts ; Q4 : 2 pts.

\begin{enumerate}
    \item \textbf{(2 pts)} Le Cameroun plaide pour une \emph{réforme du Conseil de Sécurité} afin de le rendre plus juste et représentatif : attribution à l'Afrique de deux sièges permanents avec droit de veto, dans le cadre d'un multilatéralisme rénové (1 pt). Cette revendication est portée par le \emph{Groupe africain} (position commune africaine dite « Consensus d'Ezulwini », portée par le Comité des Dix, C10) (1 pt).
    \item \textbf{(2 pts)} En déployant 850 casques bleus et policiers (MINUSCA, MONUSCO), le Cameroun met ses forces au service de la paix d'autres peuples : c'est la \emph{solidarité internationale} en action, et une contribution à la \emph{responsabilité de protéger} les populations civiles contre les violences massives (1 pt). Contrainte majeure : l'insuffisance des \emph{équipements} (matériel, logistique) ou le besoin de \emph{formation} adaptée aux mandats robustes de protection des civils (1 pt).
    \item \textbf{(2 pts)} La lutte contre les discours de haine répond aux fractures internes (crise anglophone, stigmatisations dans l'Extrême-Nord) où la haine attise les violences ; la promouvoir au CDH prolonge l'action intérieure de cohésion sociale (1 pt). Elle rejoint le \emph{Pacte de Marrakech (2018)}, qui engage les États à combattre la xénophobie et les discours de haine envers les migrants, et à promouvoir des migrations sûres, ordonnées et régulières (1 pt).
    \item \textbf{(2 pts)} \textbf{Synthèse (attendus) :} l'action internationale du Cameroun relève d'une diplomatie d'influence : en réclamant la réforme du Conseil de Sécurité, il défend à la fois l'intérêt africain et sa propre stature ; en contribuant aux OMP, il gagne en crédibilité et en poids dans les enceintes onusiennes tout en servant la paix collective ; en candidatant au Conseil des Droits de l'Homme avec des priorités cohérentes avec ses défis internes, il valorise son image et défend ses intérêts nationaux. Influence extérieure et intérêt national se renforcent ainsi mutuellement au service de la paix.
\end{enumerate}
\textbf{Points de vigilance :} nommer précisément le Consensus d'Ezulwini / C10 ; ne pas confondre R2P (protection des populations) et ingérence ; en synthèse, respecter la limite de 10 lignes et articuler les trois documents.
\end{corrige}

\newpage

\coursec{Fiche bilan}

\begin{popfigure}
\begin{tikzpicture}[
  mindmap,
  concept color=popblue,
  grow cyclic,
  level 1/.style={level distance=4.5cm,sibling angle=72},
  level 2/.style={level distance=3cm,sibling angle=45},
  every node/.style={concept, execute at begin node=\small\bfseries},
  root concept/.append style={concept color=popdark, font=\large\bfseries, text width=3.5cm, align=center},
  concept 1/.append style={concept color=popblue, text width=2.5cm},
  concept 2/.append style={concept color=popgreen, text width=2.5cm},
  concept 3/.append style={concept color=poporange, text width=2.5cm},
  concept 4/.append style={concept color=poppurple, text width=2.5cm},
  concept 5/.append style={concept color=poppink, text width=2.5cm},
]
\node[root concept, align=center]{Diplomatie\\ camerounaise\\ \& Relations\\ internationales}
  child[concept 1] { node[align=center]{1. Relations\\ internationales\\ (concept, acteurs,\\ enjeux, souveraineté)} }
  child[concept 2] { node[align=center]{2. Coopération\\ internationale\\ (bilatérale, multi-\\ latérale, décentralisée)} }
  child[concept 3] { node[align=center]{3. Assistance\\ humanitaire\\ (principes, acteurs,\\ droit humanitaire)} }
  child[concept 4] { node[align=center]{4. Principes de la\\ diplomatie\\ camerounaise\\ (non-ingérence, paix,\\ intégration, droit)} }
  child[concept 5] { node[align=center]{5. Mise en situation\\ : rédaction d'une\\ note diplomatique\\ (structure, style)} };
\end{tikzpicture}
\legende{Carte mentale de la leçon : les principes de la diplomatie camerounaise. Les cinq branches correspondent aux cinq parties du plan de la leçon.}
\end{popfigure}

\coursec{Rappel des définitions clés}

\begin{definition}[Les notions fondamentales de la leçon]
\begin{itemize}
  \item \cle{Relations internationales} : ensemble des interactions (politiques, économiques, culturelles, juridiques, militaires) entre les acteurs de la scène mondiale, qu'ils soient étatiques (États) ou non étatiques (organisations internationales, ONG, firmes multinationales).
  \item \cle{Coopération internationale} : action concertée entre États ou organisations internationales pour atteindre des objectifs communs (développement, paix, santé, éducation, sécurité).
  \item \cle{Assistance humanitaire} : aide d'urgence (vivres, soins, abris, eau potable) apportée aux populations en détresse à la suite de conflits ou de catastrophes, selon les principes d'humanité, de neutralité, d'impartialité et d'indépendance.
  \item \cle{Diplomatie} : ensemble des méthodes par lesquelles un État se représente et négocie auprès d'autres États ou d'organisations internationales afin de défendre ses intérêts nationaux par des moyens pacifiques.
  \item \cle{Note diplomatique} : document écrit officiel (formel, impersonnel, courtois) adressé par une mission diplomatique au Ministère des Relations Extérieures du pays hôte (ou inversement) pour exposer une position, protester, demander ou informer.
\end{itemize}
\end{definition}

\coursec{Principes et repères à connaître par cœur}

Le tableau ci-dessous rassemble les repères essentiels de la leçon. Chaque ligne correspond à une partie du plan ; les exemples cités doivent pouvoir être restitués le jour de l'examen.

\begin{center}
\begin{tabularx}{\linewidth}{|l|X|X|}
\hline
\rowcolor{popblueL}
\textbf{Domaine} & \textbf{Principes / points clés} & \textbf{Illustration / base juridique} \\
\hline
Relations internationales & Souveraineté des États, égalité juridique des États, non-ingérence, règlement pacifique des différends, coopération entre nations & Charte des Nations Unies (art. 2), Acte constitutif de l'Union africaine \\
\hline
Coopération internationale & Bilatérale (d'État à État), multilatérale (ONU, UA, CEMAC, CEEAC, OIF, Commonwealth), décentralisée (collectivités territoriales), Sud-Sud et triangulaire & Accords de partenariat UE--ACP (Yaoundé, Lomé, Cotonou) ; coopération Cameroun--Chine, Cameroun--France \\
\hline
Assistance humanitaire & \textbf{1. Humanité} (sauver des vies, soulager la souffrance) ; \textbf{2. Neutralité} (ne pas prendre parti dans le conflit) ; \textbf{3. Impartialité} (aider selon les seuls besoins) ; \textbf{4. Indépendance} (autonomie par rapport aux pouvoirs politiques) & Droit international humanitaire (Conventions de Genève de 1949), Code de conduite du CICR et des ONG \\
\hline
Diplomatie camerounaise & 1. Respect de la souveraineté et non-ingérence ; 2. Règlement pacifique des conflits (négociation, médiation, recours à la justice internationale) ; 3. Intégration régionale et panafricanisme (CEMAC, CEEAC, UA) ; 4. Promotion des droits humains et du droit humanitaire ; 5. Coopération gagnant-gagnant et diversification des partenariats & Préambule de la Constitution du 18 janvier 1996 ; arrêt de la Cour internationale de Justice du 10 octobre 2002 (différend Cameroun--Nigeria sur Bakassi) \\
\hline
Note diplomatique & Structure : en-tête (référence, numéro, lieu et date) $\to$ adresse (Excellence Monsieur le Ministre...) $\to$ objet $\to$ corps (exposé des motifs, argumentation, position) $\to$ formule de politesse rituelle $\to$ signature du chef de mission & Usages : protestation, demande d'appui, notification d'une nomination, rappel d'accords \\
\hline
\end{tabularx}
\end{center}

\begin{attention}[Erreurs fréquentes à éviter]
\begin{itemize}
  \item Ne pas confondre \cle{neutralité} (ne pas prendre parti dans un conflit) et \cle{impartialité} (distribuer l'aide selon les seuls besoins, sans discrimination).
  \item La non-ingérence n'interdit pas la médiation : le Cameroun peut proposer ses bons offices dans une crise voisine (par exemple en Centrafrique) sans s'immiscer dans les affaires intérieures d'un État.
  \item Une note diplomatique n'emploie jamais la première personne du singulier ni un ton familier : elle est rédigée à la troisième personne (\og L'Ambassade du Cameroun présente ses compliments... \fg).
  \item Ne pas attribuer au Cameroun des fonctions qu'il n'occupe pas : vérifier les exemples d'actualité avant de les citer dans une copie.
\end{itemize}
\end{attention}

\coursec{Méthodes express}

\begin{methode}[Analyser un accord de coopération]
\begin{enumerate}
  \item Identifier le type d'accord (bilatéral / multilatéral) et les parties signataires.
  \item Repérer l'objet de l'accord (secteur : éducation, santé, défense, commerce, infrastructures) et sa durée.
  \item Lister les engagements réciproques (obligations, financements, transferts de compétences).
  \item Vérifier la conformité de l'accord aux principes constitutionnels (souveraineté, intérêt national).
  \item Conclure sur les retombées concrètes pour les populations (emplois, infrastructures, formation).
\end{enumerate}
\end{methode}

\begin{methode}[Rédiger une note diplomatique (exigible au Bac)]
\begin{enumerate}
  \item \textbf{En-tête} : référence, numéro d'ordre, lieu et date (ex. : Yaoundé, le 15 octobre 2025).
  \item \textbf{Adresse} : \og Excellence Monsieur le Ministre des Relations Extérieures \fg{} (ou Madame/Monsieur l'Ambassadeur).
  \item \textbf{Objet} : clair, précis, en majuscules (ex. : DEMANDE D'APPUI LOGISTIQUE POUR LE RAPATRIEMENT DE RESSORTISSANTS).
  \item \textbf{Corps} : formule d'entrée (\og L'Ambassade... présente ses compliments... \fg), exposé des faits (contexte, base juridique), argumentation (intérêt commun, réciprocité, urgence humanitaire), demande ou position explicite.
  \item \textbf{Formule de politesse finale} : \og L'Ambassade... renouvelle à... les assurances de sa haute considération. \fg
  \item \textbf{Signature} : nom, qualité du signataire (chef de mission), cachet.
\end{enumerate}
\end{methode}

\begin{aretenir}[Checklist avant le Bac]
$\square$ Je sais définir les relations internationales et citer quatre acteurs majeurs (États, organisations internationales comme l'ONU, ONG, firmes multinationales).\\
$\square$ Je distingue coopération bilatérale, multilatérale et décentralisée, avec un exemple camerounais pour chacune.\\
$\square$ Je récite les quatre principes fondamentaux de l'action humanitaire (humanité, neutralité, impartialité, indépendance) et je sais les distinguer.\\
$\square$ Je cite les cinq piliers de la diplomatie camerounaise (souveraineté et non-ingérence, paix, intégration régionale, droits humains, coopération gagnant-gagnant).\\
$\square$ J'explique le principe de non-ingérence et je donne un cas d'application (ex. : la médiation du Cameroun dans la crise centrafricaine).\\
$\square$ Je situe le Cameroun dans au moins trois organisations d'intégration (UA, CEMAC, CEEAC, OIF, Commonwealth).\\
$\square$ Je sais structurer une note diplomatique complète (en-tête, adresse, objet, corps, formule de politesse, signature).\\
$\square$ Je maîtrise les formules de politesse d'usage (entrée et finale) sans les confondre.\\
$\square$ Je peux analyser un cas d'assistance humanitaire au Cameroun (ex. : personnes déplacées internes de l'Extrême-Nord ou des régions du Nord-Ouest et du Sud-Ouest) : acteurs, principes respectés, défis.\\
$\square$ Je justifie l'attachement du Cameroun au multilatéralisme (membre de l'ONU, participation aux opérations de maintien de la paix, siège dans les organes de l'UA).\\
$\square$ Je relie le Préambule de la Constitution (attachement à la paix, à l'autodétermination des peuples, à l'unité africaine) à la conduite de la politique étrangère camerounaise.
\end{aretenir}

\findecours

\end{document}
